\documentclass[12pt,a4paper]{article}
%
%	STARWARE - Stile base per la scrittura di documenti interni in LaTex
%

%
%	Pacchetti globali
%	Fare qui eventuali aggiunte!
%
\usepackage[utf8]{inputenc}
\usepackage[italian]{babel}
\usepackage[babel]{csquotes}
\usepackage{url}
\usepackage{graphicx}
\usepackage[colorlinks]{hyperref}
\usepackage{lastpage}
\usepackage{fancyhdr}
\usepackage[top=1cm,bottom=4cm,left=80pt,right=80pt]{geometry} %disegna la linea
\usepackage{listings} %per grandi porzioni di codice
\usepackage{color}
\usepackage[table]{xcolor}
\usepackage{booktabs,tabularx}
\usepackage{makeidx}
\usepackage{fixltx2e}
\usepackage{hyperref}
\usepackage{enumitem}
\usepackage{color}
\usepackage[T1]{fontenc}
\usepackage{float}
\usepackage{svg}
\usepackage{amsmath}
\usepackage[toc]{glossaries}
\usepackage{dirtree}
\usepackage{listings}
\usepackage{eurosym}

\makeglossaries

\bibliographystyle{alpha}

%
%	VARIABILI GLOBALI
%
\newcommand{\nomeGruppo}{StarWare}
\newcommand{\mailGruppo}{starware.swe@gmail.com}
\newcommand{\uni}{Universit\`{a} degli Studi di Padova}
\newcommand{\uniAA}{2015/2016}
\newcommand{\Cardin}{Prof. Riccardo Cardin}
\newcommand{\Vardanega}{Prof. Tullio Vardanega}
\newcommand{\Zucchetti}{Zucchetti S.p.a.}
\newcommand{\prj}{Quizzipedia}
\newcommand{\prjL}{Quizzipedia: software per la gestione di questionari}

\newcommand{\AVI}{Alessio Vitella}
\newcommand{\AVE}{Andrea Venier}
\newcommand{\NDC}{Nicola De Cao}
\newcommand{\IB}{Igor Baylyak}
\newcommand{\WS}{Walter Sandon}
\newcommand{\TP}{Thomas Pigarelli}
\newcommand{\AB}{Anna Bonaldo}

\newcommand{\mgls}[1]{\gls{#1}\textsubscript{G}}
\newcommand{\mglspl}[1]{\glspl{#1}\textsubscript{G}}
\newcommand{\mGls}[1]{\Gls{#1}\textsubscript{G}}
\newcommand{\mGlspl}[1]{\Glspl{#1}\textsubscript{G}}

\newcommand{\AM}{\emph{\mGls{amministratore}}}
\newcommand{\AN}{\emph{\mGls{analista}}}
\newcommand{\PG}{\emph{\mGls{progettista}}}
\newcommand{\PR}{\emph{\mGls{programmatore}}}
\newcommand{\VR}{\emph{\mGls{verificatore}}}
\newcommand{\RE}{\emph{\mGls{responsabile}}}
\newcommand{\PM}{\RE}

\newcommand{\AMpl}{\emph{\mGlspl{amministratore}}}
\newcommand{\ANpl}{\emph{\mGlspl{analista}}}
\newcommand{\PGpl}{\emph{\mGlspl{progettista}}}
\newcommand{\PRpl}{\emph{\mGlspl{programmatore}}}
\newcommand{\VRpl}{\emph{\mGlspl{verificatore}}}
\newcommand{\REpl}{\emph{\mGlspl{responsabile}}}
\newcommand{\PMpl}{\REpl}

\newcommand{\RR}{\emph{\mGls{revisione dei requisiti}}}
\newcommand{\RA}{\emph{\mGls{revisione di accettazione}}}
\newcommand{\RP}{\emph{\mGls{revisione di progettazione}}}
\newcommand{\RQ}{\emph{\mGls{revisione di qualifica}}}

\newcommand{\NdP}{\emph{\mGls{norme di progetto}}}
\newcommand{\SdF}{\emph{\mGls{studio di fattibilita}}}
\newcommand{\AdR}{\emph{\mGls{analisi dei requisiti}}}
\newcommand{\PdP}{\emph{\mGls{piano di progetto}}}
\newcommand{\PdQ}{\emph{\mGls{piano di qualifica}}}

\newcommand{\latex}[1]{\texttt{#1}}
\newcommand{\fileName}[1]{\texttt{#1}}
\newcommand{\filePath}[1]{\texttt{#1}}
\newcommand{\TODO}[1]{\texttt{\large \color{red} \underline{TODO: #1}}}

\newcommand{\licenza}{GNU GENERAL PUBLIC LICENSE V2}

%per compilare il template usare questi sotto e commentare glia altri:
%\newcommand{\logoLungo}{../imgs/logoLungo.png}
%\newcommand{\logoGrande}{../imgs/logoGrande.png}
\newcommand{\logoLungo}{../../../template/imgs/logoLungo.png}
\newcommand{\logoGrande}{../../../template/imgs/logoGrande.png}

%
%	Setup stili
%

\newcommand{\HRule}{\rule{\linewidth}{0.5mm}}

\definecolor{dkgreen}{rgb}{0,0.6,0}
\definecolor{gray}{rgb}{0.5,0.5,0.5}
\definecolor{mauve}{rgb}{0.58,0,0.82}
\definecolor{light}{RGB}{255,255,190}

%
%	Setup di pagina
%

%colorazione link
\hypersetup
{
	colorlinks=true,
	linkcolor=black,
	urlcolor=blue,
	citecolor=blue
}

%	Setup Header + Footer

\pagestyle{fancy}
\setlength{\headheight}{2cm} %settato grandezza header

\renewcommand{\footrulewidth}{0.5pt} %ridefinisco il valore della riga di intestazione
\renewcommand{\headrulewidth}{0.5pt} %ridefinisco il valore della riga di pie' di pagina
\addtolength{\headwidth}{\marginparsep}
\addtolength{\headwidth}{\marginparwidth}

\fancyhead{} %annulla head di default
\fancyfoot{} %annulla foot di default

%	Logo intestazione
\lhead{\includegraphics[scale=0.12]{\logoLungo}}


%	footer
\cfoot{
	\uni \ - \uniAA \\
	\href{mailto:\mailGruppo}{\mailGruppo}\\
	{\tiny Questo documento è distribuito sotto licenza {\licenza}}
}
\rfoot{
	\thepage\ di \pageref{LastPage}
}


% test subsubsubsection

\usepackage{titlesec}
\usepackage{hyperref}

\titleclass{\subsubsubsection}{straight}[\subsection]

\newcounter{subsubsubsection}[subsubsection]
\renewcommand\thesubsubsubsection{\thesubsubsection.\arabic{subsubsubsection}}
\renewcommand\theparagraph{\thesubsubsubsection.\arabic{paragraph}} % optional; useful if paragraphs are to be numbered

\titleformat{\subsubsubsection}
  {\normalfont\normalsize\bfseries}{\thesubsubsubsection}{1em}{}
\titlespacing*{\subsubsubsection}
{0pt}{3.25ex plus 1ex minus .2ex}{1.5ex plus .2ex}

\makeatletter
\renewcommand\paragraph{\@startsection{paragraph}{5}{\z@}%
  {3.25ex \@plus1ex \@minus.2ex}%
  {-1em}%
  {\normalfont\normalsize\bfseries}}
\renewcommand\subparagraph{\@startsection{subparagraph}{6}{\parindent}%
  {3.25ex \@plus1ex \@minus .2ex}%
  {-1em}%
  {\normalfont\normalsize\bfseries}}
\def\toclevel@subsubsubsection{4}
\def\toclevel@paragraph{5}
\def\toclevel@paragraph{6}
\def\l@subsubsubsection{\@dottedtocline{4}{7em}{4em}}
\def\l@paragraph{\@dottedtocline{5}{10em}{5em}}
\def\l@subparagraph{\@dottedtocline{6}{14em}{6em}}
\makeatother

\setcounter{secnumdepth}{4}
\setcounter{tocdepth}{4}


%pdflatex -synctex=1 -interaction=nonstopmode %.tex|makeglossaries %|pdflatex -synctex=1 -interaction=nonstopmode %.tex|pdflatex -synctex=1 -interaction=nonstopmode %.tex
\makeglossary

\newglossaryentry{responsabile} {
	name=Responsabile,
	description={è il responsabile della gestione, pianificazione e realizzazione del progetto.},
	plural=Responsabili
}

\newglossaryentry{verificatore} {
	name=Verificatore,
	description={è il responsabile dell'attività di verifica.},
	plural=Verificatori
}

\newglossaryentry{programmatore} {
	name=Programmatore,
	description={è responsabile delle attività di codifica miranti alla realizzazione del prodotto e delle componenti di ausilio necessarie per l'esecuzione delle prove di verifica e validazione.},
	plural=programmatori
}

\newglossaryentry{progettista} {
	name=Progettista,
	description={è responsabile delle attività di progettazione.},
	plural=Progettisti
}

\newglossaryentry{analista} {
	name=Analista,
	description={è responsabile delle attività di analisi. },
	plural=Analisti
}

\newglossaryentry{amministratore} {
	name=Amministratore,
	description={è responsabile dell'efficienza e dell'operatività dell'ambiente di sviluppo; si occupa della redazione e attuazione di piani e procedure di gestione della qualità; inoltre gestisce l'archivio della documentazione del progetto.},
	plural=Amministratori
}

\newglossaryentry{revisione dei requisiti} {
	name=Revisione dei requisiti,
	description={è una revisione formale che determina l'accesso del gruppo al progetto didattico e la concordanza con il cliente di una visione condivisa del prodotto atteso.}
}

\newglossaryentry{revisione di accettazione} {
	name=Revisione di accettazione,
	description={è una revisione formale per l'accertamento del soddisfacimento di tutti i requisiti e il completamento del progetto.}
}

\newglossaryentry{revisione di progettazione} {
	name=Revisione di progettazione,
	description={è una revisione di progresso che accerta la realizzabilità del prodotto e informa il cliente sulle caratteristiche del prodotto.}
}

\newglossaryentry{revisione di qualifica} {
	name=Revisione di qualifica,
	description={è una revisione di progresso che approva l'esito finale delle verifiche e attiva la fase di validazione.}
}

\newglossaryentry{analisi} {
	name=Analisi,
	description={è il periodo di preparazione e produzioni di documenti che precede la Revisione dei requisiti.}
}

\newglossaryentry{progettazione architetturale e di dettaglio} {
	name=Progettazione architetturale e di dettaglio,
	description={è il periodo che intercorre tra la Revisione dei requisiti e la Revisione di progettazione.}
}

\newglossaryentry{codifica} {
	name=Codifica,
	description={è il periodo che intercorre tra la Revisione di progettazione e la Revisione di qualifica.}
}

\newglossaryentry{verifica e validazione} {
	name=Verifica e validazione,
	description={è il periodo che intercorre tra la Revisione di qualifica e la Revisione di accettazione.}
}

\newglossaryentry{repository} {
	name=Repository,
	description={è dove i file sono memorizzati, spesso su un server.}
}

\newglossaryentry{ruolo} {
	name=Ruolo,
	description={una delle  gure professionali che una persona  sica interpreta nel corso del progetto. I ruoli sono: responsabile, amministratore, analista, progettista, programmatore e veri catore.},
    plural=ruoli
}

\newglossaryentry{SVG} {
	name=SVG,
	description={è un formato per la visualizzazione di oggetti in grafica vettoriale. Per maggiori informazioni si veda \href{https://it.wikipedia.org/wiki/Scalable_Vector_Graphics}{qui}.}
}

\newglossaryentry{PNG} {
	name=PNG,
	description={abbreviazione di Portable Network Graphics, è un formato di file per memorizzare immagini. Per ulteriori informazioni si veda http://it.wikipedia.org/wiki/Portable_Network_Graphics.}
}

\newglossaryentry{PDF} {
	name=PDF,
	description={è un formato di file basato su un linguaggio di descrizione di pagina sviluppato da Adobe Systems nel 1993 per rappresentare documenti in modo indipendente dall’hardware e dal software utilizzati per generarli o per visualizzarli. Per ulteriori informazioni si veda http: //it.wikipedia.org/wiki/Portable_Document_Format.}
}

\newglossaryentry{UML} {
	name=UML,
	description={è un linguaggio di modellazione e specifica basato sul paradigma object-oriented. Per ulteriori informazioni si veda http://it.wikipedia.org/wiki/Unified_Modeling_Language.}
}

\newglossaryentry{walkthrough} {
	name=Walkthrough,
    description={consiste nella lettura di un documento o codice cercando errori ed anomalie senza un'idea precisa di quali tipi di errori sarà possibile trovare.}
}

\newglossaryentry{lista di controllo} {
	name=Lista di controllo,
	description={è un elenco di cose da fare per eseguire una determinata attività.}
}

\newglossaryentry{inspection} {
	name=Inspection,
	description={è la lettura mirata di un documento o codice cercando errori specifici.}
}

\newglossaryentry{milestone} {
	name=Milestone,
	description={momento saliente nello sviluppo di un prodotto software per la quale devono essere pronti documenti e/o funzionalità.}
}

\newglossaryentry{ticket} {
	name=Ticket,
	description={rappresenta un compito nell'organizzazione e distribuzione del lavoro all'interno del progetto.},
	plural=tickets
}

\newglossaryentry{commit} {
	name=Commit,
	description={è la copia di modifiche fatte su file locali verso la repository remota. Esso rappresenta anche un particolare stato della repository nel tempo.}
}

\newglossaryentry{versionamento} {
	name=Versionamento,
	description={è la gestione di un versioni multiple di un insieme di informazioni. Per maggiori informazioni si veda \href{http://it. wikipedia.org/wiki/Controllo_versione}{qui}.}
}

\newglossaryentry{task} {
	name=Task,
	description={è un compito secondo la definizione dello standard IEEE 12207.},
	plural=tasks
}

\newglossaryentry{attivita} {
	name=attivita,
	description={è un insieme di task.}
}

\newglossaryentry{redattore} {
	name=Redattore,
	description={colui che redige un documento.},
	plural=redattori
}

\newglossaryentry{proponente} {
	name=Proponente,
	description={colui che ha proposto al committente un capitolato d'appalto.}
}

\newglossaryentry{committente} {
	name=Committente,
	description={colui che assegna un compito. In questo caso è il Professor Tullio Vardanega.}
}

\newglossaryentry{quality assurance} {
	name=Quality assurance,
	description={è l'insieme delle attività volte a garantire il soddisfacimento degli obiettivi della qualità.}
}

\newglossaryentry{telegram} {
	name=Telegram,
	description={è un servizio di messaggistica istantanea utilizzato dal gruppo per comunicazioni interne. Per maggiori informazioni si veda \href{https://it.wikipedia.org/wiki/Telegram_(software)}{qui}.}
}

\newglossaryentry{browser} {
	name=Browser,
	description={è un'applicazione per il recupero, la presentazione e la navigazione di risorse web.}
}

\newglossaryentry{google drive} {
	name=Google Drive,
	description={è un servizio di memorizzazione e sincronizzazione online introdotto da Google il 24 aprile 2012. Per maggiori informazioni si veda \href{https://it.wikipedia.org/wiki/Google_Drive}{qui}.}
}

\newglossaryentry{skype} {
	name=Skype,
	description={è un software proprietario freeware di messaggistica istantanea e VoIP. Per maggiori informazioni si veda \href{https://it.wikipedia.org/wiki/Skype}{qui}.}
}

\newglossaryentry{gantt} {
	name=Gantt,
	description={è un diagramma di supporto alla gestione dei progetti.}
}

\newglossaryentry{projectlibre} {
	name=ProjectLibre,
	description={è un software di gestione progettuale.}
}

\newglossaryentry{PERT} {
	name=PERT,
	description={è uno strumento volto alla programmazione delle attività che compongono il progetto e, più in generale, alla gestione degli aspetti temporali di quest'ultimo.}
}

\newglossaryentry{subtask} {
	name=Subtask,
	description={è un task compreso all'interno di un altro task. La totalità di tutti i subtasks costituisce un intero task.}
	plural=subtasks
}

\newglossaryentry{ticketing} {
	name=Ticketing,
	description={procedura con la quale il Responsabile assegna un task.}
}

\newglossaryentry{git} {
	name=Git,
	description={è un sistema software di controllo di versione distribuito}
}

\newglossaryentry{Quizzpedia} {
	name=Quizzpedia,
	description={è il nome del prodotto software richiesto dal capitolato d'appalto scelto.}
}

\newglossaryentry{schierabile} {
	name=Schierabile,
	description={è la capacità di rilasciare al cliente, con relativa installazione e messa in funzione o esercizio, di una applicazione o di un sistema software tipicamente all'interno di un sistema informatico aziendale.}
}

\newglossaryentry{cross-platform} {
	name=Cross-platform,
	description={può essere riferito ad un linguaggio di programmazione, ad un'applicazione software o ad un dispositivo hardware che funziona su più di un sistema.}
}

\newglossaryentry{QML} {
	name=QML,
	description={è un "Domain Specific Language" richiesto dal capitolato d'appalto per la definizione delle domande all'interno del sistema.}
}

\newglossaryentry{Tomcat} {
	name=Tomcat,
	description={è un application server nella forma di contenitore servlet open source sviluppato dalla Apache Software Foundation. Per maggiori informazioni si veda \href{https://it.wikipedia.org/wiki/Apache_Tomcat}{qui}.}
}

\newglossaryentry{Java} {
	name=Java,
	description={è un linguaggio di programmazione orientato agli oggetti, specificatamente progettato per essere il più possibile indipendente dalla piattaforma di esecuzione. Per maggiori informazioni si veda \href{https://it.wikipedia.org/wiki/Java_(linguaggio_di_programmazione)}{qui}.}
}

\newglossaryentry{Node.js} {
	name=Node.js,
	description={\TODO{}}
}

\newglossaryentry{JavaScript} {
	name=JavaScript,
	description={\TODO{}}
}

\newglossaryentry{PostgreSQL} {
	name=PostgreSQL,
	description={\TODO{}}
}

\newglossaryentry{MongoDB} {
	name=MongoDB,
	description={\TODO{}}
}

\newglossaryentry{HTML5} {
	name=HTML5,
	description={\TODO{}}
}

\newglossaryentry{CSS3} {
	name=CSS3,
	description={\TODO{}}
}

\newglossaryentry{XML} {
	name=XML,
	description={\TODO{}}
}

\newglossaryentry{NoSQL} {
	name=NoSQL,
	description={\TODO{}}
}

\newglossaryentry{Scala} {
	name=Scala,
	description={\TODO{}}
}

\newglossaryentry{Akka} {
	name=Akka,
	description={\TODO{}}
}

\newglossaryentry{BLE} {
	name=BLE,
	description={\TODO{}}
}

\newglossaryentry{MQTT} {
	name=MQTT,
	description={\TODO{}}
}

\newglossaryentry{AWS} {
	name=AWS,
	description={\TODO{}}
}

\newglossaryentry{Heroku} {
	name=Heroku,
	description={\TODO{}}
}

\newglossaryentry{Github} {
	name=Github,
	description={\TODO{}}
}

\newglossaryentry{template} {
	name=template,
	description={\TODO{}}
	plural=templates
}

\newglossaryentry{Dropbox} {
	name=Dropbox,
	description={\TODO{}}
}

\newglossaryentry{open-source} {
	name=open-source,
	description={\TODO{}}
}

\newglossaryentry{project management} {
	name=project management,
	description={\TODO{}}
}

\newglossaryentry{Linux} {
	name=Linux,
	description={\TODO{}}
}

\newglossaryentry{Windows} {
	name=Windows,
	description={\TODO{}}
}

\newglossaryentry{Mac OS} {
	name=Mac OS,
	description={\TODO{}}
}

\newglossaryentry{program evaluation and review techniques} {
	name=program evaluation and review techniques,
	description={\TODO{}}
}

\newglossaryentry{schedule variance} {
	name=schedule variance,
	description={\TODO{}}
}

\newglossaryentry{cost variance} {
	name=cost variance,
	description={\TODO{}}
}


\newglossaryentry{merge} {
	name=merge,
	description={\TODO{}}
}

\newglossaryentry{slack} {
	name=slack,
	description={\TODO{}}
}

\newglossaryentry{baseline} {
	name=baseline,
	description={\TODO{}}
}

\newglossaryentry{Asana} {
	name=Asana,
	description={\TODO{}}
}

\newglossaryentry{deadline} {
	name=deadline,
	description={\TODO{}}
}

\newglossaryentry{revert} {
	name=revert,
	description={\TODO{}}
}

\newglossaryentry{backup} {
	name=backup,
	description={\TODO{}}
}

\newglossaryentry{push} {
	name=push,
	description={\TODO{}}
}

\newglossaryentry{teamwork} {
	name=teamwork,
	description={\TODO{}}
}

\newglossaryentry{evento} {
	name=evento,
	description={\TODO{}}
}

\newglossaryentry{etichetta} {
	name=etichetta,
	description={\TODO{}}
	plural=etichette
}

\newglossaryentry{bug} {
	name=bug,
	description={\TODO{}}
	plural=bugs
}

\newglossaryentry{software} {
	name=software,
	description={\TODO{}}
}

\newglossaryentry{desktop} {
	name=desktop,
	description={\TODO{}}
}

\newglossaryentry{draw.io} {
	name=draw.io,
	description={\TODO{}}
}

\newglossaryentry{definizione di prodotto} {
	name=definizione di prodotto,
	description={\TODO{}}
}

\newglossaryentry{specifica tecnica} {
	name=specifica tecnica,
	description={\TODO{}}
}

\newglossaryentry{manuale utente} {
	name=manuale utente,
	description={\TODO{}}
}

\newglossaryentry{tracy} {
	name=tracy,
	description={\TODO{}}
}

\newglossaryentry{package} {
	name=package,
	description={\TODO{}}
}

\newglossaryentry{stakeholder} {
	name=stakeholder,
	description={\TODO{}}
}

\newglossaryentry{branch} {
	name=branch,
	description={\TODO{}}
	plural=branches
}

\newglossaryentry{pull} {
	name=pull,
	description={\TODO{}}
}

\newglossaryentry{underscore} {
	name=underscore,
	description={\TODO{}}
}

\newglossaryentry{spelling} {
	name=spelling,
	description={\TODO{}}
}

\newglossaryentry{cloud} {
	name=cloud,
	description={\TODO{}}
}

\newglossaryentry{smartphone} {
	name=smartphone,
	description={\TODO{}}
	plural=smartphones
}

\newglossaryentry{checkbox} {
	name=checkbox,
	description={\TODO{}}
}

\newglossaryentry{appointed} {
	name=appointed,
	description={\TODO{}}
}

\newglossaryentry{makefile} {
	name=makefile,
	description={\TODO{}}
}

\newglossaryentry{gulpease} {
	name=Gulpease,
	description={\TODO{}}
}


%Titolo documento
\newcommand{\titoloDocumento}{Analisi dei requisiti}

%Prima data di creazione del documento
\newcommand{\dataCreazione}{30 Novembre 2015}

%Inserite la versione attuale del documento
\newcommand{\versione}{3.0.0}

%Stato in cui si trova il documento: Formale solo all'atto di consegna
\newcommand{\stato}{Formale}

%Uso del documento
\newcommand{\uso}{Esterno}

\rhead{\titoloDocumento}
\lfoot{Versione: \versione}
\title{\titoloDocumento}

\begin{document}
\begin{titlepage}
\begin{center}
\intestazione{\titoloDocumento}

\begin{table}[h]
\begin{center}
\begin{tabular}{r | l}
\multicolumn{2}{c}{\textbf{Informazioni sul documento}}\\
\midrule
\textbf{Nome Documento} & \titoloDocumento \\
\textbf{Versione} & \versione \\
\textbf{Stato} & \emph{\stato} \\
\textbf{Uso} & \emph{\uso} \\
\textbf{Data Creazione} & \dataCreazione \\
\textbf{Data Ultima Modifica} & \today \\
\textbf{Redazione} & \AVE\\
\  & \TP\\
\textbf{Verifica} &  \AVI\\
\textbf{Approvazione} &  \WS\\
\textbf{Lista Distribuzione} & \nomeGruppo \\
\  & \Vardanega \\
\  & \Cardin \\
\  & Il proponente \Zucchetti \\

\end{tabular}
\end{center}
\end{table}

\end{center}
\end{titlepage}
\newpage

\Large{\textbf{Registro delle modifiche}}\\
\normalsize

\begin{center}
	\begin{longtable}[H]{p{0.12\textwidth} p{0.2\textwidth} p{0.18\textwidth} p{0.39\textwidth}}
		\toprule
		\textbf{Versione}	&	\textbf{Autore}	&	\textbf{Data}	&	\textbf{Descrizione}\\*
		\midrule
		\midrule
		3.0.0 & \WS{} & 2016-05-06 & Approvazione \\*
		\midrule
		2.1.0 & \AVI{} & 2016-05-05 & Verifica \\*
		\midrule
		2.0.3 & \AVE{} & 2016-05-04 & Aggiunta riferimenti informativi \\*
		\midrule
		2.0.2 & \AVE{} & 2016-05-03 & Correzioni tracciamento UC-Requisiti \\*
		\midrule
		2.0.1 & \TP{} & 2016-04-28 & Riordino dei casi d'uso per tipologia di attore, aggiunta diagrammi di generalizzazione per UC 18.1 \\*
		\midrule
		2.0.0 & \AB{} & 2016-02-26 & Approvazione \\*
		\midrule
		1.1.0 & \AVE{} \& \AVI{} & 2016-02-26 & Verifica \\*
		\midrule
		1.0.7 & \NDC{} & 2016-02-26 & Modifica priorità requisiti \\*
		\midrule
        1.0.6 & \TP{} & 2016-02-26 & Ridefiniti casi d'uso autenticazione \\*
		\midrule
		1.0.4 & \NDC{} & 2016-02-24 &  Correzioni segnalazioni RR \\*
		\midrule
		1.0.3 & \IB{} & 2016-02-23 &  Modifica glossario \\*
		\midrule
		1.0.2 & \IB{} & 2016-02-22 &  Aggiunta riferimenti normativi e informativi \\*
		\midrule
		1.0.1 & \NDC{} & 2016-02-20 &  Corretti diagrammi use case \\*
		\midrule
		1.0.0 & \NDC{} & 2016-01-21 &  Approvazione \\*
		\midrule
		0.1.0 & \AVI{} & 2016-01-20 &  Prima verifica \\*
		\midrule
		0.0.4 & \TP{} & 2016-01-19 & Aggiunta diagrammi casi d'uso\\*
		\midrule
		0.0.3 & \AVE{} & 2016-01-19 &  Introduzione documento e delle sezioni  \\*
		\midrule
		0.0.2 & \AB{} & 2016-01-19 &  Aggiunta tabella dei requisiti \\*
		\midrule
		0.0.1 & \WS{} & 2016-01-19 &  Aggiunta tabella dei casi d'uso \\*
		\midrule
		0.0.0 & \NDC{} & 2016-01-19 &  Creazione documento \\*
		\bottomrule
		\caption{Versionamento del documento}
		\label{tabVers1}
	\end{longtable}
\end{center}

\newpage
\tableofcontents
\newpage
\listoftables
\listoffigures
\newpage

\section{Introduzione}

\subsection{Scopo}
L’analisi del capitolato d’appalto e il successivo incontro con il proponente ha evidenziato un insieme di requisiti che il presente documento ha lo scopo di elencare e descrivere in modo dettagliato. Lo scopo di tale documento è quindi presentare le funzionalità che offrirà il prodotto.

\subsection{Descrizione}
Il documento presenterà una breve descrizione degli attori coinvolti nell'iterazione con il sistema proposto. Successivamente verranno elencati i casi d'uso individuati come essenziali per la creazione di un sistema che corrisponda alle esigenze del proponente.

\subsection{Descrizione Prodotto}
\descrizioneProdotto

\subsection{Glossario}\label{glossario}
\glossarioPrint

\subsection{Riferimenti}\label{riferimenti}
\subsubsection{Normativi}
\begin{itemize}
	\item Capitolato d’Appalto \prjL: \url{http://www.math.unipd.it/~tullio/IS-1/2015/Progetto/C5.pdf};
	\item Vincoli di organigramma e dettagli tecnico-economici: \url{http://www.math.unipd.it/~tullio/IS-1/2015/Progetto/PD01b.html};
	\item \NdP: \NdPv{} . 
\end{itemize}
\newpage

\subsubsection{Informativi}
	\begin{itemize}
		\item Slide d’insegnamento di Ingegneria del Software modulo A, Ingegneria dei requisiti: \url{http://www.math.unipd.it/~tullio/IS-1/2015/};
		\item Software Engineering - Ian Sommerville - 9th Edition (2010)
		\begin{itemize}
			\item Chapter 4: Requirements engineering;
		\end{itemize}
		\item \textit{IEEE 830-1998}: \url{http://www.math.uaa.alaska.edu/~afkjm/cs401/IEEE830.pdf};
		\item Standard \mgls{uml} 2.0 :
		\url{http://www.omg.org/spec/UML/2.0/}.
	\end{itemize}


\section{Attori}

\begin{figure}[H]
    \centering
    \begin{tikzpicture}
        \umlactor[x=0, y=0]{Studente}
        \umlactor[x=4, y=0]{Docente}
        \umlactor[x=4, y=4]{Amministratore}
        \umlactor[x=0, y=4]{Proprietario}
        \umlactor[x=0, y=-4]{Ospite}
        \umlinherit{Docente}{Studente}
        \umlinherit{Amministratore}{Docente}
        \umlinherit{Proprietario}{Amministratore}
    \end{tikzpicture}
    \caption{Attori}
    \label{fig:actors} 
\end{figure}

Di seguito viene fornita una breve descrizione degli attori utilizzati nei casi d'uso:

\begin{itemize}
    \item \textbf{Ospite:} utente che non ha ancora effettuato l'accesso ed è quindi limitato
        ad accedere o registrarsi
    \item \textbf{Studente:} utente che è registrato al servizio, ha effettuato l'accesso al
        sistema e esegue questionari
    \item \textbf{Docente:} utente in grado di gestire studenti, classi, argomenti, questionari
        e domande
    \item \textbf{Amministratore:} utente che gestisce i docenti
    \item \textbf{Proprietario:} amministratore iniziale che è in grado di creare nuovi amministratori
        e cancellarne di esistenti
\end{itemize}

Nei diagrammi seguenti, le generalizzazioni tra attori non vengono specificate perché
appesantiscono troppo l'immagine. Come riferimento è possibile utilizzare il diagramma mostrato in questa
sezione.

\newpage
\section{Casi d'uso}

\subsection{Struttura di un caso d'uso}
I casi d'uso sono creati secondo lo standard UML 2.4 e sono identificati dalla seguente notazione:
\begin{center}
	UC[codice]: [Titolo]
\end{center}
dove il \textbf{codice} è un numero progressivo identificativo di ogni requisito, gerarchico nel  caso di sotto-casi d'uso tramite la notazione \textit{CodiceUCPadre.CodiceSottoUC}. Per ogni caso d'uso deve inoltre essere indicato:
\begin{itemize}
	\item \textbf{Titolo:} breve ma non ambiguo
	\item \textbf{Attori:} principali e secondari coinvolti
	\item \textbf{Precondizione}
	\item \textbf{Flusso principale degli eventi:} descrizione, eventualmente per punti, del main scenario del Caso d'uso
	\item \textbf{Postcondizione}
\end{itemize}

% Output di tracy sotto a:
% tracy-tpiga.rhcloud.com/tracy/latex/UseCases?type=latex&raw=1
\hypertarget{UC1}{}
\subsection{Caso d'uso UC1: Autenticazione}
        \begin{figure}[H]
            \centering
            \begin{resizedtikzpicture}{\textwidth}
		\umlactor[x=0, y=-1]{Ospite}
		\umlactor[x=0, y=1]{Google}
		\begin{umlsystem}[x=0, fill=lightgray!20]{Quizzpedia}
			\umlusecase[x=4, y=0, fill=white, width=3cm, name=2]{\textbf{UC1:} Autenticazione}
			\umlassoc{Ospite}{2}
			\umlassoc{Google}{2}
			\umlusecase[x=11, y=2.3333333333333, fill=white, width=3cm, name=170]{\textbf{UC2:} Autenticazione con Google}
			\umlinherit{170}{2}
			\umlusecase[x=11, y=0, fill=white, width=3cm, name=173]{\textbf{UC3:} Autenticazione con Facebook}
			\umlinherit{173}{2}
			\umlusecase[x=11, y=-2.3333333333333, fill=white, width=3cm, name=174]{\textbf{UC4:} Autenticazione con Windows}
			\umlinherit{174}{2}
		\end{umlsystem}
            \end{resizedtikzpicture}
            \caption{Caso d'uso UC1: Autenticazione}
            \label{fig:UC1} 
        \end{figure}
    \begin{description}
\item[Attori:] Ospite, Google;
\item[Scopo e descrizione:] L'ospite effettua l'autenticazione presso il sistema
      \item[Precondizione:] L'ospite non è autenticato presso il sistema;

        \item[Flusso principale degli eventi:] \begin{enumerate}
          \item L'ospite si autentica selezionando uno dei provider forniti;

      \end{enumerate}
    \item[Postcondizione:] L'ospite è autenticato presso il sistema.
  \end{description}
\hypertarget{UC1.1}{}
\subsection{Caso d'uso UC1.1: Registrazione}\begin{description}
\item[Attori:] Ospite;
\item[Scopo e descrizione:] Il sistema crea un'account automaticamente alla prima autenticazione, deducendo le informazioni personali con quelle fornite dal provider 
      \item[Precondizione:] L'ospite non è autenticato nel sistema e non ha un account ;

        \item[Flusso principale degli eventi:] \begin{enumerate}
          \item Viene creato automaticamente un account per l'ospite nel sistema;

      \end{enumerate}
    \item[Postcondizione:] L’ospite possiede un’account di tipo studente presso il sistema.
  \end{description}
\hypertarget{UC1.2}{}
\subsection{Caso d'uso UC1.2: Errore autenticazione annullata}\begin{description}
\item[Attori:] Amministratore;
\item[Scopo e descrizione:] Viene visualizzato un messaggio d'errore nel caso in cui l'autenticazione venga annullata dall'ospite
      \item[Precondizione:] L'ospite ha annullato l'autenticazione;

        \item[Flusso principale degli eventi:] \begin{enumerate}
          \item Viene visualizzato un messaggio d'errore;

      \end{enumerate}
    \item[Postcondizione:] Viene mostrato un messaggio d'errore.
  \end{description}
\hypertarget{UC2}{}
\subsection{Caso d'uso UC2: Autenticazione con Google}
        \begin{figure}[H]
            \centering
            \begin{resizedtikzpicture}{\textwidth}
		\umlactor[x=0, y=0]{Ospite}
		\begin{umlsystem}[x=0, fill=lightgray!20]{Quizzpedia}
			\umlusecase[x=4, y=0, fill=white, width=3cm, name=170]{\textbf{UC2:} Autenticazione con Google}
			\umlassoc{Ospite}{170}
			\umlusecase[x=11, y=1.2916666666667, fill=white, width=3cm, name=1]{\textbf{UC1.1:} Registrazione}
			\umlusecase[x=11, y=-1.2916666666667, fill=white, width=3cm, name=172]{\textbf{UC1.2:} Errore autenticazione annullata}
			\umlextend{1}{170}
			\umlextend{172}{170}
		\end{umlsystem}
            \end{resizedtikzpicture}
            \caption{Caso d'uso UC2: Autenticazione con Google}
            \label{fig:UC2} 
        \end{figure}
    \begin{description}
\item[Attori:] Ospite;
\item[Scopo e descrizione:] L'ospite seleziona la funzionalità di login tramite credenziali Google
      \item[Precondizione:] L'ospite non è autenticato presso il sistema;
    \item[Estensioni:]
      \begin{enumerate}
          \item Se l'ospite ha annullato l'autenticazione viene visualizzato un messaggio di errore (\hyperlink{UC1.2}{UC1.2});
          \item Viene creato un'account alla prima autenticazione (\hyperlink{UC1.1}{UC1.1});

      \end{enumerate}
    \item[Postcondizione:] L'ospite è autenticato con le credenziali di Google.
  \end{description}
\hypertarget{UC3}{}
\subsection{Caso d'uso UC3: Autenticazione con Facebook}
        \begin{figure}[H]
            \centering
            \begin{resizedtikzpicture}{\textwidth}
		\umlactor[x=0, y=-1]{Ospite}
		\umlactor[x=0, y=1]{Facebook}
		\begin{umlsystem}[x=0, fill=lightgray!20]{Quizzpedia}
			\umlusecase[x=4, y=0, fill=white, width=3cm, name=173]{\textbf{UC3:} Autenticazione con Facebook}
			\umlassoc{Ospite}{173}
			\umlassoc{Facebook}{173}
			\umlusecase[x=11, y=1.2916666666667, fill=white, width=3cm, name=1]{\textbf{UC1.1:} Registrazione}
			\umlusecase[x=11, y=-1.2916666666667, fill=white, width=3cm, name=172]{\textbf{UC1.2:} Errore autenticazione annullata}
			\umlextend{1}{173}
			\umlextend{172}{173}
		\end{umlsystem}
            \end{resizedtikzpicture}
            \caption{Caso d'uso UC3: Autenticazione con Facebook}
            \label{fig:UC3} 
        \end{figure}
    \begin{description}
\item[Attori:] Ospite, Facebook;
\item[Scopo e descrizione:] L'ospite seleziona la funzionalità di login tramite credenziali Facebook
      \item[Precondizione:] L'ospite non è autenticato presso il sistema;
    \item[Estensioni:]
      \begin{enumerate}
          \item Se l'ospite ha annullato l'autenticazione viene visualizzato un messaggio di errore (\hyperlink{UC1.2}{UC1.2});
          \item Viene creato un'account alla prima autenticazione	 (\hyperlink{UC1.1}{UC1.1});

      \end{enumerate}
    \item[Postcondizione:] L'ospite è autenticato con le credenziali di Facebook.
  \end{description}
\hypertarget{UC4}{}
\subsection{Caso d'uso UC4: Autenticazione con Windows}
        \begin{figure}[H]
            \centering
            \begin{resizedtikzpicture}{\textwidth}
		\umlactor[x=0, y=-1]{Ospite}
		\umlactor[x=0, y=1]{Windows}
		\begin{umlsystem}[x=0, fill=lightgray!20]{Quizzpedia}
			\umlusecase[x=4, y=0, fill=white, width=3cm, name=174]{\textbf{UC4:} Autenticazione con Windows}
			\umlassoc{Ospite}{174}
			\umlassoc{Windows}{174}
			\umlusecase[x=11, y=1.2916666666667, fill=white, width=3cm, name=1]{\textbf{UC1.1:} Registrazione}
			\umlusecase[x=11, y=-1.2916666666667, fill=white, width=3cm, name=172]{\textbf{UC1.2:} Errore autenticazione annullata}
			\umlextend{1}{174}
			\umlextend{172}{174}
		\end{umlsystem}
            \end{resizedtikzpicture}
            \caption{Caso d'uso UC4: Autenticazione con Windows}
            \label{fig:UC4} 
        \end{figure}
    \begin{description}
\item[Attori:] Ospite, Windows;
\item[Scopo e descrizione:] L'ospite seleziona la funzionalità di login tramite credenziali Windows
      \item[Precondizione:] L'ospite non è autenticato presso il sistema
;
    \item[Estensioni:]
      \begin{enumerate}
          \item Se l'ospite ha annullato l'autenticazione viene visualizzato un messaggio di errore (\hyperlink{UC1.2}{UC1.2});
          \item Viene creato un'account alla prima autenticazione	 (\hyperlink{UC1.1}{UC1.1});

      \end{enumerate}
    \item[Postcondizione:] L'ospite è autenticato con le credenziali di Windows.
  \end{description}
\hypertarget{UC5}{}
\subsection{Caso d'uso UC5: Gestione account}
        \begin{figure}[H]
            \centering
            \begin{resizedtikzpicture}{\textwidth}
		\umlactor[x=0, y=0]{Utente}
		\begin{umlsystem}[x=0, fill=lightgray!20]{Quizzpedia}
			\umlusecase[x=6, y=0, fill=white, width=3cm, name=41]{\textbf{UC5.1:} Modifica informazioni personali}
			\umlassoc{Utente}{41}
		\end{umlsystem}
            \end{resizedtikzpicture}
            \caption{Caso d'uso UC5: Gestione account}
            \label{fig:UC5} 
        \end{figure}
    \begin{description}
\item[Attori:] Utente;
\item[Scopo e descrizione:] L'utente gestisce il proprio account
      \item[Precondizione:] L'utente è autenticato nel sistema;

        \item[Flusso principale degli eventi:] \begin{enumerate}
          \item L’utente visualizza le proprie informazioni personali;
          \item L'utente modifica le sue informazioni personali (\hyperlink{UC5.1}{UC5.1});
          \item L'utente conferma le modifiche effettuate;

      \end{enumerate}
    \item[Postcondizione:] Il sistema ha apportato le modifiche all'account dell'utente.
  \end{description}
\hypertarget{UC5.1}{}
\subsection{Caso d'uso UC5.1: Modifica informazioni personali}
        \begin{figure}[H]
            \centering
            \begin{resizedtikzpicture}{\textwidth}
		\umlactor[x=0, y=0]{Utente}
		\begin{umlsystem}[x=0, fill=lightgray!20]{Quizzpedia}
			\umlusecase[x=6, y=0, fill=white, width=3cm, name=129]{\textbf{UC5.1.1:} Modifica nome}
			\umlassoc{Utente}{129}
		\end{umlsystem}
            \end{resizedtikzpicture}
            \caption{Caso d'uso UC5.1: Modifica informazioni personali}
            \label{fig:UC5.1} 
        \end{figure}
    \begin{description}
\item[Attori:] Utente;
\item[Scopo e descrizione:] L'utente modifica le proprie informazione personali
      \item[Precondizione:] L'utente è autenticato nel sistema;

        \item[Flusso principale degli eventi:] \begin{enumerate}
          \item L’utente modifica il proprio nome completo (\hyperlink{UC5.1.1}{UC5.1.1});

      \end{enumerate}
    \item[Estensioni:]
      \begin{enumerate}
          \item Se l'utente ha cancellato uno dei campi delle informazioni personali viene visualizzato un messaggio di errore (\hyperlink{UC5.2}{UC5.2});

      \end{enumerate}
    \item[Postcondizione:] L'utente ha modificato le proprie informazioni personali.
  \end{description}
\hypertarget{UC5.1.1}{}
\subsection{Caso d'uso UC5.1.1: Modifica nome}\begin{description}
\item[Attori:] Utente;
\item[Scopo e descrizione:] L'utente modifica il proprio nome nelle impostazioni del profilo
      \item[Precondizione:] L'utente è autenticato nel sistema;

        \item[Flusso principale degli eventi:] \begin{enumerate}
          \item L'utente modifica il nome;

      \end{enumerate}
    \item[Postcondizione:] L'utente ha modificato il nome.
  \end{description}
\hypertarget{UC5.2}{}
\subsection{Caso d'uso UC5.2: Errore cancellazione campi obbligatori}\begin{description}
\item[Attori:] Utente;
\item[Scopo e descrizione:] Il sistema avvisa l'utente che non è possibile cancellare le informazioni personali obbligatorie
      \item[Precondizione:] L'utente non ha inserito tutti i campi;

        \item[Flusso principale degli eventi:] \begin{enumerate}
          \item Viene visualizzato un messaggio d'errore;

      \end{enumerate}
    \item[Postcondizione:] Non viene effettuata la modifica.
  \end{description}
\hypertarget{UC6}{}
\subsection{Caso d'uso UC6: Logout}\begin{description}
\item[Attori:] Utente;
\item[Scopo e descrizione:] L'utente effettua il logout
      \item[Precondizione:] L'utente è autenticato nel sistema;

        \item[Flusso principale degli eventi:] \begin{enumerate}
          \item L'utente seleziona la funzionalità di Logout;

      \end{enumerate}
    \item[Postcondizione:] L'utente non è più autenticato nel sistema.
  \end{description}
\hypertarget{UC7}{}
\subsection{Caso d'uso UC7: Gestione domande}
        \begin{figure}[H]
            \centering
            \begin{resizedtikzpicture}{\textwidth}
		\umlactor[x=0, y=0]{Docente}
		\begin{umlsystem}[x=0, fill=lightgray!20]{Quizzpedia}
			\umlusecase[x=6, y=2.0208333333333, fill=white, width=3cm, name=43]{\textbf{UC7.1:} Inserimento domanda}
			\umlassoc{Docente}{43}
			\umlusecase[x=6, y=-0.0625, fill=white, width=3cm, name=46]{\textbf{UC7.2:} Modifica domanda}
			\umlassoc{Docente}{46}
			\umlusecase[x=6, y=-2.0208333333333, fill=white, width=3cm, name=50]{\textbf{UC7.3:} Elimina domanda}
			\umlassoc{Docente}{50}
		\end{umlsystem}
            \end{resizedtikzpicture}
            \caption{Caso d'uso UC7: Gestione domande}
            \label{fig:UC7} 
        \end{figure}
    \begin{description}
\item[Attori:] Docente;
\item[Scopo e descrizione:] Il docente gestisce le proprie domande
      \item[Precondizione:] Il docente è autenticato nel sistema;

        \item[Flusso principale degli eventi:] \begin{enumerate}
          \item Il docente può creare una nuova domanda (\hyperlink{UC7.1}{UC7.1});
          \item Il docente può modificare una domanda (\hyperlink{UC7.2}{UC7.2});
          \item Il docente può eliminare una domanda (\hyperlink{UC7.3}{UC7.3});

      \end{enumerate}
    \item[Postcondizione:] Il sistema ha ottenuto le informazioni sulle operazioni che il docente desidera eseguire sulle domande.
  \end{description}
\hypertarget{UC7.1}{}
\subsection{Caso d'uso UC7.1: Inserimento domanda}
        \begin{figure}[H]
            \centering
            \begin{resizedtikzpicture}{\textwidth}
		\umlactor[x=0, y=0]{Docente}
		\begin{umlsystem}[x=0, fill=lightgray!20]{Quizzpedia}
			\umlusecase[x=4, y=0, fill=white, width=3cm, name=43]{\textbf{UC7.1:} Inserimento domanda}
			\umlassoc{Docente}{43}
			\umlusecase[x=11, y=7.9791666666667, fill=white, width=3cm, name=148]{\textbf{UC7.6:} Inserimento domanda di tipo vero/falso}
			\umlinherit{148}{43}
			\umlusecase[x=11, y=5.1041666666667, fill=white, width=3cm, name=149]{\textbf{UC7.7:} Inserimento domanda a scelta multipla}
			\umlinherit{149}{43}
			\umlusecase[x=11, y=2.1875, fill=white, width=3cm, name=150]{\textbf{UC7.8:} Inserimento domanda a risposta multipla}
			\umlinherit{150}{43}
			\umlusecase[x=11, y=-1.2291666666667, fill=white, width=3cm, name=151]{\textbf{UC7.9:} Inserimento domanda di tipo testo con parole omesse}
			\umlinherit{151}{43}
			\umlusecase[x=11, y=-4.6458333333333, fill=white, width=3cm, name=152]{\textbf{UC7.10:} Inserimento domanda con l'associazione di parole}
			\umlinherit{152}{43}
			\umlusecase[x=11, y=-7.9791666666667, fill=white, width=3cm, name=153]{\textbf{UC7.11:} Inserimento domanda a risposta aperta}
			\umlinherit{153}{43}
		\end{umlsystem}
            \end{resizedtikzpicture}
            \caption{Caso d'uso UC7.1: Inserimento domanda}
            \label{fig:UC7.1} 
        \end{figure}
    \begin{description}
\item[Attori:] Docente;
\item[Scopo e descrizione:] Il docente compone una domanda in linguaggio QML, che verrà salvata nel sistema e potrà essere utilizzata nei questionari
      \item[Precondizione:] Il docente è autenticato nel sistema;

        \item[Flusso principale degli eventi:] \begin{enumerate}
          \item Il docente seleziona gli argomenti della domanda (\hyperlink{UC7.1.1}{UC7.1.1});
          \item Il docente compone la domanda in QML  (\hyperlink{UC7.1.2}{UC7.1.2});
          \item Il docente compone la domanda attraverso l'interfaccia grafica (\hyperlink{UC7.1.3}{UC7.1.3});
          \item Il docente conferma la creazione della domanda;

      \end{enumerate}
    \item[Estensioni:]
      \begin{enumerate}
          \item Se il codice QML inserito non è valido viene visualizzato un messaggio di errore (\hyperlink{UC7.4}{UC7.4});
          \item Se non è stato selezionato nessun argomento viene visualizzato un messaggio di errore (\hyperlink{UC7.5}{UC7.5});

      \end{enumerate}
    \item[Postcondizione:] È stata creata una nuova domanda.
  \end{description}
\hypertarget{UC7.1.1}{}
\subsection{Caso d'uso UC7.1.1: Selezione argomenti nuova domanda}\begin{description}
\item[Attori:] Docente;
\item[Scopo e descrizione:] Il docente seleziona gli argomenti corrispondenti alla domanda selezionata
      \item[Precondizione:] Il docente sta creando una nuova domanda;

        \item[Flusso principale degli eventi:] \begin{enumerate}
          \item Il docente seleziona gli argomenti della domanda;

      \end{enumerate}
    \item[Postcondizione:] Il docente ha definito gli argomenti per classificare la domanda
.
  \end{description}
\hypertarget{UC7.1.2}{}
\subsection{Caso d'uso UC7.1.2: Scrittura domanda in QML della nuova domanda}\begin{description}
\item[Attori:] Docente;
\item[Scopo e descrizione:] Il docente compone una domanda in linguaggio QML
      \item[Precondizione:] Il docente sta creando una nuova domanda;

        \item[Flusso principale degli eventi:] \begin{enumerate}
          \item Il docente compone la domanda in QML;

      \end{enumerate}
    \item[Postcondizione:] La domanda è stata scritta in linguaggio QML.
  \end{description}
\hypertarget{UC7.1.3}{}
\subsection{Caso d'uso UC7.1.3: Scrittura nuova domanda da interfaccia grafica}\begin{description}
\item[Attori:] Docente;
\item[Scopo e descrizione:] Il docente compone la domanda attraverso un interfaccia grafica
      \item[Precondizione:] Il docente sta creando una nuova domanda;

        \item[Flusso principale degli eventi:] \begin{enumerate}
          \item Il docente compone la domanda attraverso l'interfaccia grafica;

      \end{enumerate}
    \item[Postcondizione:] La domanda è stata composta attraverso l'interfaccia grafica.
  \end{description}
\hypertarget{UC7.2}{}
\subsection{Caso d'uso UC7.2: Modifica domanda}
        \begin{figure}[H]
            \centering
            \begin{resizedtikzpicture}{\textwidth}
		\umlactor[x=0, y=0]{Docente}
		\begin{umlsystem}[x=0, fill=lightgray!20]{Quizzpedia}
			\umlusecase[x=6, y=3.5416666666667, fill=white, width=3cm, name=163]{\textbf{UC7.2.1:} Selezione argomenti modifica domanda}
			\umlassoc{Docente}{163}
			\umlusecase[x=6, y=0, fill=white, width=3cm, name=164]{\textbf{UC7.2.2:} Scrittura domanda in QML della domanda da modificare}
			\umlassoc{Docente}{164}
			\umlusecase[x=6, y=-3.5416666666667, fill=white, width=3cm, name=165]{\textbf{UC7.2.3:} Modifica domanda da interfaccia grafica}
			\umlassoc{Docente}{165}
		\end{umlsystem}
            \end{resizedtikzpicture}
            \caption{Caso d'uso UC7.2: Modifica domanda}
            \label{fig:UC7.2} 
        \end{figure}
    \begin{description}
\item[Attori:] Docente;
\item[Scopo e descrizione:] Il docente effettua delle modifiche ad una domanda 
      \item[Precondizione:] Il docente è autenticato nel sistema;

        \item[Flusso principale degli eventi:] \begin{enumerate}
          \item Ricerca della domanda da modificare (\hyperlink{UC24}{UC24});
          \item Selezione della domanda da modificare;
          \item Il docente seleziona gli argomenti della domanda	 (\hyperlink{UC7.2.1}{UC7.2.1});
          \item Il docente modifica la domanda in QML	 (\hyperlink{UC7.2.2}{UC7.2.2});
          \item Il docente compone la domanda attraverso l'interfaccia grafica (\hyperlink{UC7.2.3}{UC7.2.3});
          \item Conferma delle modifiche;

      \end{enumerate}
    \item[Estensioni:]
      \begin{enumerate}
          \item Se il codice QML inserito non è valido viene visualizzato un messaggio di errore (\hyperlink{UC7.4}{UC7.4});
          \item Se non è stato selezionato nessun argomento viene visualizzato un messaggio di errore (\hyperlink{UC7.5}{UC7.5});

      \end{enumerate}
    \item[Postcondizione:] La domanda è stata modificata.
  \end{description}
\hypertarget{UC7.2.1}{}
\subsection{Caso d'uso UC7.2.1: Selezione argomenti modifica domanda}\begin{description}
\item[Attori:] Docente;
\item[Scopo e descrizione:] Il docente seleziona gli argomenti corrispondenti alla domanda selezionata
      \item[Precondizione:] Il docente sta modificando una domanda
;

        \item[Flusso principale degli eventi:] \begin{enumerate}
          \item Il docente seleziona gli argomenti della domanda;

      \end{enumerate}
    \item[Postcondizione:] Il docente ha definito gli argomenti per classificare la domanda.
  \end{description}
\hypertarget{UC7.2.2}{}
\subsection{Caso d'uso UC7.2.2: Scrittura domanda in QML della domanda da modificare}\begin{description}
\item[Attori:] Docente;
\item[Scopo e descrizione:] Il docente compone una domanda in linguaggio QML
      \item[Precondizione:] Il docente sta modificando una domanda;

        \item[Flusso principale degli eventi:] \begin{enumerate}
          \item Il docente compone la domanda in QML;

      \end{enumerate}
    \item[Postcondizione:] La domanda è stata scritta in linguaggio QML.
  \end{description}
\hypertarget{UC7.2.3}{}
\subsection{Caso d'uso UC7.2.3: Modifica domanda da interfaccia grafica}\begin{description}
\item[Attori:] Docente;
\item[Scopo e descrizione:] Il docente modifica la domanda attraverso un interfaccia grafica

      \item[Precondizione:] Il docente sta modificando una domanda
;

        \item[Flusso principale degli eventi:] \begin{enumerate}
          \item Il docente modifica la domanda attraverso l'interfaccia grafica;

      \end{enumerate}
    \item[Postcondizione:] La domanda è stata modificata attraverso l'interfaccia grafica.
  \end{description}
\hypertarget{UC7.3}{}
\subsection{Caso d'uso UC7.3: Elimina domanda}\begin{description}
\item[Attori:] Docente;
\item[Scopo e descrizione:] Il docente rimuove dal sistema una domanda da lui creata
      \item[Precondizione:] Il docente è autenticato nel sistema;

        \item[Flusso principale degli eventi:] \begin{enumerate}
          \item Ricerca della domanda da eliminare (\hyperlink{UC24}{UC24});
          \item Selezione della domanda da eliminare;
          \item Conferma per l'eliminazione della domanda	;

      \end{enumerate}
    \item[Postcondizione:] La domanda è stata eliminata.
  \end{description}
\hypertarget{UC7.4}{}
\subsection{Caso d'uso UC7.4: Errore QML non valido}\begin{description}
\item[Attori:] Docente;
\item[Scopo e descrizione:] Il sistema avvisa il docente di uno o più errori nel codice QML inserito
      \item[Precondizione:] Il QML non è valido;

        \item[Flusso principale degli eventi:] \begin{enumerate}
          \item Viene visualizzato un messaggio d'errore;

      \end{enumerate}
    \item[Postcondizione:] Non viene inserita la domanda.
  \end{description}
\hypertarget{UC7.5}{}
\subsection{Caso d'uso UC7.5: Errore argomento mancante}\begin{description}
\item[Attori:] Docente;
\item[Scopo e descrizione:] Il sistema avvisa il docente che deve essere inserito almeno un argomento
      \item[Precondizione:] Non è stato selezionato almeno un argomento;

        \item[Flusso principale degli eventi:] \begin{enumerate}
          \item Viene visualizzato un messaggio d'errore;

      \end{enumerate}
    \item[Postcondizione:] Non viene inserita la domanda.
  \end{description}
\hypertarget{UC7.6}{}
\subsection{Caso d'uso UC7.6: Inserimento domanda di tipo vero/falso}\begin{description}
\item[Attori:] Docente;
\item[Scopo e descrizione:] Il docente inserisce una domanda di tipo vero/falso
      \item[Precondizione:] Il docente è autenticato presso il sistema;

        \item[Flusso principale degli eventi:] \begin{enumerate}
          \item Il docente seleziona gli argomenti della domanda (\hyperlink{UC7.1.1}{UC7.1.1});
          \item Il docente compone la domanda vero/falso in QML specificando la risposta esatta (\hyperlink{UC7.1.2}{UC7.1.2});
          \item Il docente conferma la creazione della domanda;

      \end{enumerate}
    \item[Postcondizione:] È stata creata una nuova domanda di tipo vero/falso.
  \end{description}
\hypertarget{UC7.7}{}
\subsection{Caso d'uso UC7.7: Inserimento domanda a scelta multipla}\begin{description}
\item[Attori:] Docente;
\item[Scopo e descrizione:] Il docente inserisce una domanda a scelta multipla
      \item[Precondizione:] Il docente è autenticato presso il sistema;

        \item[Flusso principale degli eventi:] \begin{enumerate}
          \item Il docente seleziona gli argomenti della domanda (\hyperlink{UC7.1.1}{UC7.1.1});
          \item Il docente compone la domanda a scelta multipla in QML specificando testo della domanda, risposte possibili e risposta corretta (\hyperlink{UC7.1.2}{UC7.1.2});
          \item Il docente conferma la creazione della domanda;

      \end{enumerate}
    \item[Postcondizione:] È stata creata una nuova domanda a scelta multipla.
  \end{description}
\hypertarget{UC7.8}{}
\subsection{Caso d'uso UC7.8: Inserimento domanda a risposta multipla}\begin{description}
\item[Attori:] Docente;
\item[Scopo e descrizione:] Il docente inserisce una domanda a risposta multipla
      \item[Precondizione:] Il docente è autenticato presso il sistema;

        \item[Flusso principale degli eventi:] \begin{enumerate}
          \item Il docente seleziona gli argomenti della domanda (\hyperlink{UC7.1.1}{UC7.1.1});
          \item Il docente compone la domanda a risposta multipla in QML specificando il testo della domanda, risposte possibili e risposta/e esatte (\hyperlink{UC7.1.2}{UC7.1.2});
          \item Il docente conferma la creazione della domanda;

      \end{enumerate}
    \item[Postcondizione:] È stata creata una nuova domanda a risposta multipla.
  \end{description}
\hypertarget{UC7.9}{}
\subsection{Caso d'uso UC7.9: Inserimento domanda di tipo testo con parole omesse}\begin{description}
\item[Attori:] Docente;
\item[Scopo e descrizione:] Il docente inserisce una domanda di tipo testo con parole omesse
      \item[Precondizione:] Il docente è autenticato presso il sistema;

        \item[Flusso principale degli eventi:] \begin{enumerate}
          \item Il docente seleziona gli argomenti della domanda (\hyperlink{UC7.1.1}{UC7.1.1});
          \item Il docente compone il testo della domanda in QML specificando quali parole possano essere scelte e dove vadano messe  (\hyperlink{UC7.1.2}{UC7.1.2});
          \item Il docente conferma la creazione della domanda;

      \end{enumerate}
    \item[Postcondizione:] È stata creata una domanda di tipo testo con parole omesse.
  \end{description}
\hypertarget{UC7.10}{}
\subsection{Caso d'uso UC7.10: Inserimento domanda con l'associazione di parole}\begin{description}
\item[Attori:] Docente;
\item[Scopo e descrizione:] Il docente inserisce una domanda con associazione di parole
      \item[Precondizione:] Il docente è autenticato presso il sistema;

        \item[Flusso principale degli eventi:] \begin{enumerate}
          \item Il docente seleziona gli argomenti della domanda (\hyperlink{UC7.1.1}{UC7.1.1});
          \item Il docente compone la domanda in QML specificando le parole che possono essere combinate e le giuste combinazioni (\hyperlink{UC7.1.2}{UC7.1.2});
          \item Il docente conferma la creazione della domanda;

      \end{enumerate}
    \item[Postcondizione:] Stata inserita una domanda con associazione di parole.
  \end{description}
\hypertarget{UC7.11}{}
\subsection{Caso d'uso UC7.11: Inserimento domanda a risposta aperta}\begin{description}
\item[Attori:] Docente;
\item[Scopo e descrizione:] Il docente inserisce una domanda a risposta aperta
      \item[Precondizione:] Il docente è autenticato presso il sistema;

        \item[Flusso principale degli eventi:] \begin{enumerate}
          \item Il docente seleziona gli argomenti della domanda (\hyperlink{UC7.1.1}{UC7.1.1});
          \item Il docente compone la domanda in QML specificando il testo della domanda (ma non la risposta) (\hyperlink{UC7.1.2}{UC7.1.2});
          \item Il docente conferma la creazione della domanda;

      \end{enumerate}
    \item[Postcondizione:] È stata inserita una domanda a risposta aperta.
  \end{description}
\hypertarget{UC8}{}
\subsection{Caso d'uso UC8: Gestione questionari}
        \begin{figure}[H]
            \centering
            \begin{resizedtikzpicture}{\textwidth}
		\umlactor[x=0, y=0]{Docente}
		\begin{umlsystem}[x=0, fill=lightgray!20]{Quizzpedia}
			\umlusecase[x=6, y=2.1875, fill=white, width=3cm, name=42]{\textbf{UC8.1:} Inserisci questionario}
			\umlassoc{Docente}{42}
			\umlusecase[x=6, y=-0.020833333333333, fill=white, width=3cm, name=47]{\textbf{UC8.2:} Modifica questionario}
			\umlassoc{Docente}{47}
			\umlusecase[x=6, y=-2.1875, fill=white, width=3cm, name=48]{\textbf{UC8.3:} Elimina questionario}
			\umlassoc{Docente}{48}
		\end{umlsystem}
            \end{resizedtikzpicture}
            \caption{Caso d'uso UC8: Gestione questionari}
            \label{fig:UC8} 
        \end{figure}
    \begin{description}
\item[Attori:] Docente;
\item[Scopo e descrizione:] Il docente gestisce i propri questionari
      \item[Precondizione:] Il docente è autenticato nel sistema;

        \item[Flusso principale degli eventi:] \begin{enumerate}
          \item Il docente può creare un nuovo questionario (\hyperlink{UC8.1}{UC8.1});
          \item Il docente può modificare un questionario (\hyperlink{UC8.2}{UC8.2});
          \item Il docente può eliminare un questionario (\hyperlink{UC8.3}{UC8.3});

      \end{enumerate}
    \item[Postcondizione:] Il sistema ha ottenuto le informazioni sulle operazioni che il docente desidera eseguire su un questionario.
  \end{description}
\hypertarget{UC8.1}{}
\subsection{Caso d'uso UC8.1: Inserisci questionario}
        \begin{figure}[H]
            \centering
            \begin{resizedtikzpicture}{\textwidth}
		\umlactor[x=0, y=0]{Docente}
		\begin{umlsystem}[x=0, fill=lightgray!20]{Quizzpedia}
			\umlusecase[x=6, y=3.125, fill=white, width=3cm, name=91]{\textbf{UC8.1.1:} Aggiungi domanda in un nuovo questionario }
			\umlassoc{Docente}{91}
			\umlusecase[x=6, y=0, fill=white, width=3cm, name=93]{\textbf{UC8.1.2:} Elimina domanda da un nuovo questionario}
			\umlassoc{Docente}{93}
			\umlusecase[x=6, y=-3.125, fill=white, width=3cm, name=118]{\textbf{UC8.1.3:} Seleziona argomenti del nuovo questionario}
			\umlassoc{Docente}{118}
		\end{umlsystem}
            \end{resizedtikzpicture}
            \caption{Caso d'uso UC8.1: Inserisci questionario}
            \label{fig:UC8.1} 
        \end{figure}
    \begin{description}
\item[Attori:] Docente;
\item[Scopo e descrizione:] Il docente crea un questionario
      \item[Precondizione:] Il docente è autenticato nel sistema;

        \item[Flusso principale degli eventi:] \begin{enumerate}
          \item Il docente può aggiungere domande al questionario (\hyperlink{UC8.1.1}{UC8.1.1});
          \item Il docente può togliere domande precedentemente aggiunte al questionario (\hyperlink{UC8.1.2}{UC8.1.2});
          \item Il docente seleziona gli argomenti del questionario (\hyperlink{UC8.1.3}{UC8.1.3});
          \item Il docente aggiunge il titolo del questionario;
          \item Il docente conferma la creazione del questionario;

      \end{enumerate}
    \item[Estensioni:]
      \begin{enumerate}
          \item Se il questionario non ha domande viene visualizzato un messaggio d'errore (\hyperlink{UC8.4}{UC8.4});

      \end{enumerate}
    \item[Postcondizione:] È stato creato un nuovo questionario.
  \end{description}
\hypertarget{UC8.1.1}{}
\subsection{Caso d'uso UC8.1.1: Aggiungi domanda in un nuovo questionario }\begin{description}
\item[Attori:] Docente;
\item[Scopo e descrizione:] Il docente ricerca e seleziona una domanda da inserire in un nuovo questionario
      \item[Precondizione:] Il docente sta creando un nuovo questionario;

        \item[Flusso principale degli eventi:] \begin{enumerate}
          \item Viene ricercata un domanda (\hyperlink{UC24}{UC24});
          \item Selezione della domanda;
          \item Conferma inserimento domanda;

      \end{enumerate}
    \item[Postcondizione:] È stata aggiunta una domanda al nuovo questionario.
  \end{description}
\hypertarget{UC8.1.2}{}
\subsection{Caso d'uso UC8.1.2: Elimina domanda da un nuovo questionario}\begin{description}
\item[Attori:] Docente;
\item[Scopo e descrizione:] Il docente elimina una domanda da un nuovo questionario
      \item[Precondizione:] Il docente sta creando un nuovo questionario e ha selezionato una domanda da eliminare;

        \item[Flusso principale degli eventi:] \begin{enumerate}
          \item Viene eliminata la domanda;

      \end{enumerate}
    \item[Postcondizione:] È stato eliminata la domanda dal nuovo questionario.
  \end{description}
\hypertarget{UC8.1.3}{}
\subsection{Caso d'uso UC8.1.3: Seleziona argomenti del nuovo questionario}\begin{description}
\item[Attori:] Docente;
\item[Scopo e descrizione:] Il docente seleziona gli argomenti relativi al nuovo questionario 
      \item[Precondizione:] Il docente sta creando un nuovo questionario;

        \item[Flusso principale degli eventi:] \begin{enumerate}
          \item Il docente aggiunge un argomento al questionario;
          \item Il docente toglie un argomento al questionario;

      \end{enumerate}
    \item[Postcondizione:] Il docente ha selezionato gli argomenti del nuovo questionario.
  \end{description}
\hypertarget{UC8.2}{}
\subsection{Caso d'uso UC8.2: Modifica questionario}
        \begin{figure}[H]
            \centering
            \begin{resizedtikzpicture}{\textwidth}
		\umlactor[x=0, y=0]{Docente}
		\begin{umlsystem}[x=0, fill=lightgray!20]{Quizzpedia}
			\umlusecase[x=6, y=3.4166666666667, fill=white, width=3cm, name=166]{\textbf{UC8.2.1:} Aggiungi domanda in un questionario}
			\umlassoc{Docente}{166}
			\umlusecase[x=6, y=0, fill=white, width=3cm, name=167]{\textbf{UC8.2.2:} Elimina domanda da un questionario da  modificare}
			\umlassoc{Docente}{167}
			\umlusecase[x=6, y=-3.4166666666667, fill=white, width=3cm, name=168]{\textbf{UC8.2.3:} Selezione argomenti modifica questionario}
			\umlassoc{Docente}{168}
		\end{umlsystem}
            \end{resizedtikzpicture}
            \caption{Caso d'uso UC8.2: Modifica questionario}
            \label{fig:UC8.2} 
        \end{figure}
    \begin{description}
\item[Attori:] Docente;
\item[Scopo e descrizione:] Il docente modifica il questionario che ha selezionato potendo aggiungere e rimuovere domande e modificare gli argomenti del questionario
      \item[Precondizione:] Il docente è autenticato nel sistema;

        \item[Flusso principale degli eventi:] \begin{enumerate}
          \item Ricerca del questionario da modificare (\hyperlink{UC18}{UC18});
          \item Selezione del questionario da modificare;
          \item Il docente può aggiungere domande al questionario (\hyperlink{UC8.2.1}{UC8.2.1});
          \item Il docente può togliere domande precedentemente aggiunte al questionario (\hyperlink{UC8.2.2}{UC8.2.2});
          \item Il docente può modificare gli argomenti del questionario (\hyperlink{UC8.2.3}{UC8.2.3});
          \item Il docente conferma la modifica del questionario;

      \end{enumerate}
    \item[Estensioni:]
      \begin{enumerate}
          \item Se il questionario non ha domande viene visualizzato un messaggio d'errore (\hyperlink{UC8.4}{UC8.4});

      \end{enumerate}
    \item[Postcondizione:] È stato modificato il questionario.
  \end{description}
\hypertarget{UC8.2.1}{}
\subsection{Caso d'uso UC8.2.1: Aggiungi domanda in un questionario}\begin{description}
\item[Attori:] Docente;
\item[Scopo e descrizione:] Il docente inserisce una domanda da un questionario da modificare
      \item[Precondizione:] Il docente ha selezionato un questionario da modificare e una domanda da inserire;

        \item[Flusso principale degli eventi:] \begin{enumerate}
          \item Viene ricercata una domanda (\hyperlink{UC24}{UC24});
          \item Selezione della domanda	;
          \item Conferma inserimento domanda;

      \end{enumerate}
    \item[Postcondizione:] È stata aggiunta una domanda al questionario.
  \end{description}
\hypertarget{UC8.2.2}{}
\subsection{Caso d'uso UC8.2.2: Elimina domanda da un questionario da  modificare}\begin{description}
\item[Attori:] Docente;
\item[Scopo e descrizione:] Il docente elimina una domanda da un questionario da modificare
      \item[Precondizione:] Il docente ha selezionato un questionario da modificare e una domanda da eliminare;

        \item[Flusso principale degli eventi:] \begin{enumerate}
          \item Viene eliminata la domanda;

      \end{enumerate}
    \item[Postcondizione:] È stato eliminata la domanda questionario da modificare.
  \end{description}
\hypertarget{UC8.2.3}{}
\subsection{Caso d'uso UC8.2.3: Selezione argomenti modifica questionario}\begin{description}
\item[Attori:] Docente;
\item[Scopo e descrizione:] Il docente seleziona gli argomenti corrispondenti al questionario selezionato
      \item[Precondizione:] Il docente sta modificando un questionario;

        \item[Flusso principale degli eventi:] \begin{enumerate}
          \item Il docente aggiunge un argomento al questionario;
          \item Il docente toglie un argomento al questionario;

      \end{enumerate}
    \item[Postcondizione:] Il docente ha definito gli argomenti per classificare il questionario.
  \end{description}
\hypertarget{UC8.3}{}
\subsection{Caso d'uso UC8.3: Elimina questionario}\begin{description}
\item[Attori:] Docente;
\item[Scopo e descrizione:] Il docente rimuove un questionario dal sistema
      \item[Precondizione:] Il docente è autenticato nel sistema;

        \item[Flusso principale degli eventi:] \begin{enumerate}
          \item Ricerca il questionario da eliminare (\hyperlink{UC18}{UC18});
          \item Selezione del questionario da eliminare;
          \item Il docente conferma l'eliminazione;

      \end{enumerate}
    \item[Postcondizione:] È stato eliminato il questionario.
  \end{description}
\hypertarget{UC8.4}{}
\subsection{Caso d'uso UC8.4: Errore questionario vuoto}\begin{description}
\item[Attori:] Docente;
\item[Scopo e descrizione:] Il sistema avvisa il docente che nel questionario deve essere presente almeno una domanda
      \item[Precondizione:] Il questionario selezionato non ha domande;

        \item[Flusso principale degli eventi:] \begin{enumerate}
          \item Viene visualizzato un messaggio di errore;

      \end{enumerate}
    \item[Postcondizione:] Il questionario non viene inserito nel sistema.
  \end{description}
\hypertarget{UC9}{}
\subsection{Caso d'uso UC9: Gestione classi}
        \begin{figure}[H]
            \centering
            \begin{resizedtikzpicture}{\textwidth}
		\umlactor[x=0, y=0]{Docente}
		\begin{umlsystem}[x=0, fill=lightgray!20]{Quizzpedia}
			\umlusecase[x=6, y=1.9375, fill=white, width=3cm, name=44]{\textbf{UC9.1:} Inserisci classe}
			\umlassoc{Docente}{44}
			\umlusecase[x=6, y=-0.020833333333333, fill=white, width=3cm, name=45]{\textbf{UC9.2:} Modifica classe}
			\umlassoc{Docente}{45}
			\umlusecase[x=6, y=-1.9375, fill=white, width=3cm, name=49]{\textbf{UC9.3:} Elimina classe}
			\umlassoc{Docente}{49}
		\end{umlsystem}
            \end{resizedtikzpicture}
            \caption{Caso d'uso UC9: Gestione classi}
            \label{fig:UC9} 
        \end{figure}
    \begin{description}
\item[Attori:] Docente;
\item[Scopo e descrizione:] Il docente gestisce le proprie classi
      \item[Precondizione:] Il docente è autenticato nel sistema;

        \item[Flusso principale degli eventi:] \begin{enumerate}
          \item Il docente può creare una nuova classe (\hyperlink{UC9.1}{UC9.1});
          \item Il docente può modificare una classe (\hyperlink{UC9.2}{UC9.2});
          \item Il docente può eliminare una classe (\hyperlink{UC9.3}{UC9.3});

      \end{enumerate}
    \item[Postcondizione:] Il sistema ha ottenuto le informazioni sulle operazioni che il docente desidera eseguire sulla classe.
  \end{description}
\hypertarget{UC9.1}{}
\subsection{Caso d'uso UC9.1: Inserisci classe}
        \begin{figure}[H]
            \centering
            \begin{resizedtikzpicture}{\textwidth}
		\umlactor[x=0, y=0]{Docente}
		\begin{umlsystem}[x=0, fill=lightgray!20]{Quizzpedia}
			\umlusecase[x=6, y=2.4583333333333, fill=white, width=3cm, name=90]{\textbf{UC9.1.1:} Inserisci nome classe}
			\umlassoc{Docente}{90}
			\umlusecase[x=6, y=0, fill=white, width=3cm, name=94]{\textbf{UC9.1.2:} Inserisci argomenti classe}
			\umlassoc{Docente}{94}
			\umlusecase[x=6, y=-2.4583333333333, fill=white, width=3cm, name=95]{\textbf{UC9.1.3:} Inserisci password classe}
			\umlassoc{Docente}{95}
		\end{umlsystem}
            \end{resizedtikzpicture}
            \caption{Caso d'uso UC9.1: Inserisci classe}
            \label{fig:UC9.1} 
        \end{figure}
    \begin{description}
\item[Attori:] Docente;
\item[Scopo e descrizione:] Il docente crea una nuova classe alla quale gli studenti potranno iscriversi 
      \item[Precondizione:] Il docente è autenticato nel sistema;

        \item[Flusso principale degli eventi:] \begin{enumerate}
          \item Viene inserito il nome della classe (\hyperlink{UC9.1.1}{UC9.1.1});
          \item Vengono inseriti gli argomenti della classe (\hyperlink{UC9.1.2}{UC9.1.2});
          \item Viene inserita la password della classe (\hyperlink{UC9.1.3}{UC9.1.3});
          \item Viene confermato l'inserimento della classe;

      \end{enumerate}
    \item[Postcondizione:] È stato creata una nuova classe.
  \end{description}
\hypertarget{UC9.1.1}{}
\subsection{Caso d'uso UC9.1.1: Inserisci nome classe}
        \begin{figure}[H]
            \centering
            \begin{resizedtikzpicture}{\textwidth}
		\umlactor[x=0, y=0]{Docente}
		\begin{umlsystem}[x=0, fill=lightgray!20]{Quizzpedia}
			\umlusecase[x=4, y=0, fill=white, width=3cm, name=90]{\textbf{UC9.1.1:} Inserisci nome classe}
			\umlassoc{Docente}{90}
			\umlusecase[x=11, y=0, fill=white, width=3cm, name=92]{\textbf{UC9.4:} Errore nome classe già presente}
			\umlextend{92}{90}
		\end{umlsystem}
            \end{resizedtikzpicture}
            \caption{Caso d'uso UC9.1.1: Inserisci nome classe}
            \label{fig:UC9.1.1} 
        \end{figure}
    \begin{description}
\item[Attori:] Docente;
\item[Scopo e descrizione:] Il docente inserisce il nome scelto per la classe
      \item[Precondizione:] Il docente ha selezionato la classe;

        \item[Flusso principale degli eventi:] \begin{enumerate}
          \item Il docente inserisce un nome per la classe non ancora presente nel sistema;

      \end{enumerate}
    \item[Estensioni:]
      \begin{enumerate}
          \item Se il nome della classe è già presente nel sistema viene visualizzato un messaggio di errore (\hyperlink{UC9.4}{UC9.4});

      \end{enumerate}
    \item[Postcondizione:] Il docente ha inserito il nome della classe.
  \end{description}
\hypertarget{UC9.1.2}{}
\subsection{Caso d'uso UC9.1.2: Inserisci argomenti classe}\begin{description}
\item[Attori:] Docente;
\item[Scopo e descrizione:] Il docente specifica di quali argomenti tratta la classe selezionata
      \item[Precondizione:] Il docente ha selezionato una classe;

        \item[Flusso principale degli eventi:] \begin{enumerate}
          \item Il docente specifica di quali argomenti tratta la classe selezionata;

      \end{enumerate}
    \item[Postcondizione:] Il docente ha inserito gli argomenti della classe.
  \end{description}
\hypertarget{UC9.1.3}{}
\subsection{Caso d'uso UC9.1.3: Inserisci password classe}\begin{description}
\item[Attori:] Docente;
\item[Scopo e descrizione:] Il docente specifica la password che verrà fornita poi agli studenti in classe o per altre vie per poter accedere alla classe
      \item[Precondizione:] Il docente ha selezionato una classe;

        \item[Flusso principale degli eventi:] \begin{enumerate}
          \item Il docente inserisce una password che servirà per accedere alla classe che verrà creata;

      \end{enumerate}
    \item[Postcondizione:] Il docente ha inserito la password per accedere alla classe.
  \end{description}
\hypertarget{UC9.2}{}
\subsection{Caso d'uso UC9.2: Modifica classe}\begin{description}
\item[Attori:] Docente;
\item[Scopo e descrizione:] Il docente modifica degli attributi o dei dati legati alla classe scelta
      \item[Precondizione:] Il docente è autenticato nel sistema, è presente nel sistema la classe da modificare;

        \item[Flusso principale degli eventi:] \begin{enumerate}
          \item Viene selezionata una delle proprie classi da modificare;
          \item Può venire modificato il nome (\hyperlink{UC9.1.1}{UC9.1.1});
          \item Possono venire modificati gli argomenti (\hyperlink{UC9.1.2}{UC9.1.2});
          \item Può venire modificata la password (\hyperlink{UC9.1.3}{UC9.1.3});
          \item Viene confermata la modifica;

      \end{enumerate}
    \item[Postcondizione:] È stata modificata la classe.
  \end{description}
\hypertarget{UC9.3}{}
\subsection{Caso d'uso UC9.3: Elimina classe}\begin{description}
\item[Attori:] Docente;
\item[Scopo e descrizione:] Il docente rimuove dal sistema una classe da lui creata 
      \item[Precondizione:] Il docente è autenticato nel sistema, è presente nel sistema la classe da eliminare;

        \item[Flusso principale degli eventi:] \begin{enumerate}
          \item Viene sezionata una delle proprie classi;
          \item Viene sezionata la funzionalità elimina;
          \item Viene confermata l'eliminazione;

      \end{enumerate}
    \item[Postcondizione:] È stata eliminata la classe.
  \end{description}
\hypertarget{UC9.4}{}
\subsection{Caso d'uso UC9.4: Errore nome classe già presente}\begin{description}
\item[Attori:] Docente;
\item[Scopo e descrizione:] Viene visualizzato un errore nel caso il nome della classe che si sta tentando di inserire esiste già
      \item[Precondizione:] Il nome della classe è già presente nel sistema;

        \item[Flusso principale degli eventi:] \begin{enumerate}
          \item Viene visualizzato un messaggio di errore;

      \end{enumerate}
    \item[Postcondizione:] La classe non viene inserita nel sistema.
  \end{description}
\hypertarget{UC10}{}
\subsection{Caso d'uso UC10: Gestione argomenti}
        \begin{figure}[H]
            \centering
            \begin{resizedtikzpicture}{\textwidth}
		\umlactor[x=0, y=0]{Docente}
		\begin{umlsystem}[x=0, fill=lightgray!20]{Quizzpedia}
			\umlusecase[x=6, y=3.2916666666667, fill=white, width=3cm, name=135]{\textbf{UC10.1:} Esplorazione argomenti}
			\umlassoc{Docente}{135}
			\umlusecase[x=6, y=1.0416666666667, fill=white, width=3cm, name=137]{\textbf{UC10.2:} Crea argomento}
			\umlassoc{Docente}{137}
			\umlusecase[x=6, y=-1.0416666666667, fill=white, width=3cm, name=141]{\textbf{UC10.4:} Modifica argomento}
			\umlassoc{Docente}{141}
			\umlusecase[x=6, y=-3.2916666666667, fill=white, width=3cm, name=146]{\textbf{UC10.5:} Eliminazione argomento}
			\umlassoc{Docente}{146}
		\end{umlsystem}
            \end{resizedtikzpicture}
            \caption{Caso d'uso UC10: Gestione argomenti}
            \label{fig:UC10} 
        \end{figure}
    \begin{description}
\item[Attori:] Docente;
\item[Scopo e descrizione:] Il docente gestisce gli argomenti
      \item[Precondizione:] Il docente è autenticato nel sistema;

        \item[Flusso principale degli eventi:] \begin{enumerate}
          \item Il docente può esplorare gli argomenti (\hyperlink{UC10.1}{UC10.1});
          \item Il docente può creare un argomento (\hyperlink{UC10.2}{UC10.2});
          \item Il docente può modificare un argomento (\hyperlink{UC10.4}{UC10.4});
          \item Il docente può eliminare un argomento (\hyperlink{UC10.5}{UC10.5});

      \end{enumerate}
    \item[Postcondizione:] Il sistema ha ottenuto le informazioni sulle operazioni che il docente desidera eseguire su un argomento.
  \end{description}
\hypertarget{UC10.1}{}
\subsection{Caso d'uso UC10.1: Esplorazione argomenti}\begin{description}
\item[Attori:] Docente;
\item[Scopo e descrizione:] Il docente visualizza gli argomenti presenti nel sistema
      \item[Precondizione:] Il docente è autenticato nel sistema;

        \item[Flusso principale degli eventi:] \begin{enumerate}
          \item Il docente visualizza l'albero degli argomenti;

      \end{enumerate}
    \item[Postcondizione:] Il docente ha visualizzato l'albero degli argomenti.
  \end{description}
\hypertarget{UC10.2}{}
\subsection{Caso d'uso UC10.2: Crea argomento}
        \begin{figure}[H]
            \centering
            \begin{resizedtikzpicture}{\textwidth}
		\umlactor[x=0, y=0]{Docente}
		\begin{umlsystem}[x=0, fill=lightgray!20]{Quizzpedia}
			\umlusecase[x=4, y=0, fill=white, width=3cm, name=137]{\textbf{UC10.2:} Crea argomento}
			\umlassoc{Docente}{137}
			\umlusecase[x=11, y=0, fill=white, width=3cm, name=140]{\textbf{UC10.3:} Argomento già presente nel sistema}
			\umlextend{140}{137}
		\end{umlsystem}
            \end{resizedtikzpicture}
            \caption{Caso d'uso UC10.2: Crea argomento}
            \label{fig:UC10.2} 
        \end{figure}
    \begin{description}
\item[Attori:] Docente;
\item[Scopo e descrizione:] Il docente aggiunge un nuovo argomento nel sistema
      \item[Precondizione:] Il docente è autenticato nel sistema;

        \item[Flusso principale degli eventi:] \begin{enumerate}
          \item Il docente inserisce il nome dell'argomento;
          \item Il docente seleziona l'argomento padre (\hyperlink{UC10.1}{UC10.1});
          \item Il docente conferma la creazione dell'argomento;

      \end{enumerate}
    \item[Estensioni:]
      \begin{enumerate}
          \item Se l'argomento già presente nel sistema viene visualizzato un messaggio di errore (\hyperlink{UC10.3}{UC10.3});

      \end{enumerate}
    \item[Postcondizione:] L'argomento inserito è ora presente tra gli argomenti del sistema.
  \end{description}
\hypertarget{UC10.3}{}
\subsection{Caso d'uso UC10.3: Argomento già presente nel sistema}\begin{description}
\item[Attori:] Docente;
\item[Scopo e descrizione:] Il sistema avvisa il docente che non è possibile inserire un argomento già presente
      \item[Precondizione:] L'argomento è già presente nel sistema;

        \item[Flusso principale degli eventi:] \begin{enumerate}
          \item Viene visualizzato un messaggio d'errore;

      \end{enumerate}
    \item[Postcondizione:] L'argomento non viene inserito nel sistema.
  \end{description}
\hypertarget{UC10.4}{}
\subsection{Caso d'uso UC10.4: Modifica argomento}
        \begin{figure}[H]
            \centering
            \begin{resizedtikzpicture}{\textwidth}
		\umlactor[x=0, y=0]{Docente}
		\begin{umlsystem}[x=0, fill=lightgray!20]{Quizzpedia}
			\umlusecase[x=4, y=0, fill=white, width=3cm, name=141]{\textbf{UC10.4:} Modifica argomento}
			\umlassoc{Docente}{141}
			\umlusecase[x=11, y=0, fill=white, width=3cm, name=140]{\textbf{UC10.3:} Argomento già presente nel sistema}
			\umlextend{140}{141}
		\end{umlsystem}
            \end{resizedtikzpicture}
            \caption{Caso d'uso UC10.4: Modifica argomento}
            \label{fig:UC10.4} 
        \end{figure}
    \begin{description}
\item[Attori:] Docente;
\item[Scopo e descrizione:] Il docente modifica il nome di un argomento
      \item[Precondizione:] Il docente è autenticato nel sistema;

        \item[Flusso principale degli eventi:] \begin{enumerate}
          \item Il docente seleziona un argomento da modificare (\hyperlink{UC10.1}{UC10.1});
          \item Il docente può modificare il nome dell'argomento;
          \item Il docente può cambiare l'argomento padre (\hyperlink{UC10.1}{UC10.1});
          \item Il docente conferma la modifica dell'argomento;

      \end{enumerate}
    \item[Estensioni:]
      \begin{enumerate}
          \item Se l'argomento già presente nel sistema viene visualizzato un messaggio di errore (\hyperlink{UC10.3}{UC10.3});

      \end{enumerate}
    \item[Postcondizione:] È stato modificato un argomento.
  \end{description}
\hypertarget{UC10.5}{}
\subsection{Caso d'uso UC10.5: Eliminazione argomento}
        \begin{figure}[H]
            \centering
            \begin{resizedtikzpicture}{\textwidth}
		\umlactor[x=0, y=0]{Docente}
		\begin{umlsystem}[x=0, fill=lightgray!20]{Quizzpedia}
			\umlusecase[x=4, y=0, fill=white, width=3cm, name=146]{\textbf{UC10.5:} Eliminazione argomento}
			\umlassoc{Docente}{146}
			\umlusecase[x=11, y=0, fill=white, width=3cm, name=147]{\textbf{UC10.6:} Errore l'argomento ha domande o questionari}
			\umlextend{147}{146}
		\end{umlsystem}
            \end{resizedtikzpicture}
            \caption{Caso d'uso UC10.5: Eliminazione argomento}
            \label{fig:UC10.5} 
        \end{figure}
    \begin{description}
\item[Attori:] Docente;
\item[Scopo e descrizione:] Il docente elimina un argomento già presente nel sistema
      \item[Precondizione:] Il docente è autenticato nel sistema;

        \item[Flusso principale degli eventi:] \begin{enumerate}
          \item Il docente seleziona l'argomento da eliminare (\hyperlink{UC10.1}{UC10.1});
          \item Il docente conferma l'eliminazione;

      \end{enumerate}
    \item[Estensioni:]
      \begin{enumerate}
          \item Se l'argomento ha domande o questionari viene visualizzato un messaggio d'errore (\hyperlink{UC10.6}{UC10.6});

      \end{enumerate}
    \item[Postcondizione:] È stato eliminato un argomento.
  \end{description}
\hypertarget{UC10.6}{}
\subsection{Caso d'uso UC10.6: Errore l'argomento ha domande o questionari}\begin{description}
\item[Attori:] Docente;
\item[Scopo e descrizione:] Il sistema avvisa il docente che non può essere eliminato un argomento di cui esistano ancora domande o questionari
      \item[Precondizione:] L'argomento ha domande o questionari al suo interno;

        \item[Flusso principale degli eventi:] \begin{enumerate}
          \item Viene visualizzato un messaggio d'errore;

      \end{enumerate}
    \item[Postcondizione:] L'argomento non viene eliminato.
  \end{description}
\hypertarget{UC11}{}
\subsection{Caso d'uso UC11: Visualizza statistiche}
        \begin{figure}[H]
            \centering
            \begin{resizedtikzpicture}{\textwidth}
		\umlactor[x=0, y=0]{Docente}
		\begin{umlsystem}[x=0, fill=lightgray!20]{Quizzpedia}
			\umlusecase[x=4, y=0, fill=white, width=3cm, name=10]{\textbf{UC11:} Visualizza statistiche}
			\umlassoc{Docente}{10}
			\umlusecase[x=11, y=2.7083333333333, fill=white, width=3cm, name=11]{\textbf{UC12:} Visualizza statistiche domanda}
			\umlinherit{11}{10}
			\umlusecase[x=11, y=0, fill=white, width=3cm, name=12]{\textbf{UC13:} Visualizza statistiche questionario}
			\umlinherit{12}{10}
			\umlusecase[x=11, y=-2.7083333333333, fill=white, width=3cm, name=13]{\textbf{UC14:} Visualizza statistiche classe}
			\umlinherit{13}{10}
		\end{umlsystem}
            \end{resizedtikzpicture}
            \caption{Caso d'uso UC11: Visualizza statistiche}
            \label{fig:UC11} 
        \end{figure}
    \begin{description}
\item[Attori:] Docente;
\item[Scopo e descrizione:] Il docente visualizza le statistiche
      \item[Precondizione:] Il docente è autenticato presso il sistema;

        \item[Flusso principale degli eventi:] \begin{enumerate}
          \item Il docente può visualizzare le statistiche;

      \end{enumerate}
    \item[Postcondizione:] Il docente ha visualizzato le statistiche a cui era interessato.
  \end{description}
\hypertarget{UC12}{}
\subsection{Caso d'uso UC12: Visualizza statistiche domanda}\begin{description}
\item[Attori:] Docente;
\item[Scopo e descrizione:] Il docente visualizza il numero di risposte totali, risposte corrette, risposte errate e la percentuale di risposte corrette sul totale
      \item[Precondizione:] Il docente è autenticato presso il sistema;

        \item[Flusso principale degli eventi:] \begin{enumerate}
          \item Il docente ricerca una domanda (\hyperlink{UC24}{UC24});
          \item Il docente seleziona la domanda interessata;

      \end{enumerate}
    \item[Postcondizione:] Il docente ha visualizzato le statistiche relative alla domanda a cui era interessato.
  \end{description}
\hypertarget{UC13}{}
\subsection{Caso d'uso UC13: Visualizza statistiche questionario}\begin{description}
\item[Attori:] Docente;
\item[Scopo e descrizione:] Il docente visualizza per ogni punteggio la percentuale degli studenti che ottenuto tale punteggio e la media del punteggio
      \item[Precondizione:] Il docente è autenticato presso il sistema;

        \item[Flusso principale degli eventi:] \begin{enumerate}
          \item Il docente ricerca un questionario (\hyperlink{UC18}{UC18});
          \item Il docente seleziona il questionario interessato;

      \end{enumerate}
    \item[Postcondizione:] Il docente ha visualizzato le statistiche relative al questionario a cui era interessato.
  \end{description}
\hypertarget{UC14}{}
\subsection{Caso d'uso UC14: Visualizza statistiche classe}
        \begin{figure}[H]
            \centering
            \begin{resizedtikzpicture}{\textwidth}
		\umlactor[x=0, y=0]{Docente}
		\begin{umlsystem}[x=0, fill=lightgray!20]{Quizzpedia}
			\umlusecase[x=6, y=4.7916666666667, fill=white, width=3cm, name=107]{\textbf{UC14.1:} Visualizza risultati domande della classe}
			\umlassoc{Docente}{107}
			\umlusecase[x=6, y=1.5833333333333, fill=white, width=3cm, name=108]{\textbf{UC14.2:} Visualizza risultati questionari della classe}
			\umlassoc{Docente}{108}
			\umlusecase[x=6, y=-1.625, fill=white, width=3cm, name=109]{\textbf{UC14.3:} Visualizza sommario statistiche classe}
			\umlassoc{Docente}{109}
			\umlusecase[x=6, y=-4.7916666666667, fill=white, width=3cm, name=110]{\textbf{UC14.4:} Visualizza statistiche studente della classe}
			\umlassoc{Docente}{110}
		\end{umlsystem}
            \end{resizedtikzpicture}
            \caption{Caso d'uso UC14: Visualizza statistiche classe}
            \label{fig:UC14} 
        \end{figure}
    \begin{description}
\item[Attori:] Docente;
\item[Scopo e descrizione:] Il docente visualizza le statistiche di una sua classe
      \item[Precondizione:] Il docente è autenticato presso il sistema;

        \item[Flusso principale degli eventi:] \begin{enumerate}
          \item Il docente seleziona la classe di cui vuole visualizzare le statistiche;
          \item Il docente può visualizzare i risultati delle sue domande della classe (\hyperlink{UC14.1}{UC14.1});
          \item Il docente può visualizzare i risultati dei suoi questionari della classe (\hyperlink{UC14.2}{UC14.2});
          \item Il docente può visualizzare un sommario delle statistiche della classe (\hyperlink{UC14.3}{UC14.3});
          \item Il docente può visualizzare i risultati dei test di uno studente della classe (\hyperlink{UC14.4}{UC14.4});

      \end{enumerate}
    \item[Postcondizione:] Il docente ha visualizzato le statistiche relative ad una delle sue classi.
  \end{description}
\hypertarget{UC14.1}{}
\subsection{Caso d'uso UC14.1: Visualizza risultati domande della classe}\begin{description}
\item[Attori:] Docente;
\item[Scopo e descrizione:] Il docente visualizza i risultati e le statistiche relative alle domande della classe selezionata
      \item[Precondizione:] Il docente ha selezionato una classe;

        \item[Flusso principale degli eventi:] \begin{enumerate}
          \item Il docente cerca una domanda tra quelle della classe (\hyperlink{UC24}{UC24});
          \item Il docente seleziona la domanda per visualizzarne i risultati e le statistiche;

      \end{enumerate}
    \item[Postcondizione:] Il docente ha visualizzato i risultati e le statistiche relative alle domande desiderate.
  \end{description}
\hypertarget{UC14.2}{}
\subsection{Caso d'uso UC14.2: Visualizza risultati questionari della classe}\begin{description}
\item[Attori:] Docente;
\item[Scopo e descrizione:] Il docente visualizza i risultati e le statistiche relative ai questionari della classe
      \item[Precondizione:] Il docente ha selezionato una classe;

        \item[Flusso principale degli eventi:] \begin{enumerate}
          \item Il docente cerca un questionario tra quelli della classe (\hyperlink{UC18}{UC18});
          \item Il docente seleziona il questionario per visualizzarne i risultati e le statistiche;

      \end{enumerate}
    \item[Postcondizione:] Il docente ha visualizzato i risultati e le statistiche relative ai questionari desiderate.
  \end{description}
\hypertarget{UC14.3}{}
\subsection{Caso d'uso UC14.3: Visualizza sommario statistiche classe}\begin{description}
\item[Attori:] Docente;
\item[Scopo e descrizione:] Il docente visualizza le statistiche generali relative alla classe selezionata
      \item[Precondizione:] Il docente ha selezionato una classe;

        \item[Flusso principale degli eventi:] \begin{enumerate}
          \item Il docente seleziona la funzionalità per visualizzare il sommario delle statistiche della classe;

      \end{enumerate}
    \item[Postcondizione:] Il docente ha visualizzato le statistiche generali relative alla classe selezionata.
  \end{description}
\hypertarget{UC14.4}{}
\subsection{Caso d'uso UC14.4: Visualizza statistiche studente della classe}
        \begin{figure}[H]
            \centering
            \begin{resizedtikzpicture}{\textwidth}
		\umlactor[x=0, y=0]{Docente}
		\begin{umlsystem}[x=0, fill=lightgray!20]{Quizzpedia}
			\umlusecase[x=6, y=0, fill=white, width=3cm, name=111]{\textbf{UC14.4.1:} Visualizza risultati questionari dello studente}
			\umlassoc{Docente}{111}
		\end{umlsystem}
            \end{resizedtikzpicture}
            \caption{Caso d'uso UC14.4: Visualizza statistiche studente della classe}
            \label{fig:UC14.4} 
        \end{figure}
    \begin{description}
\item[Attori:] Docente;
\item[Scopo e descrizione:] Il docente visualizza i risultati e le statistiche relative alla classe selezionata
      \item[Precondizione:] Il docente ha selezionato una classe;

        \item[Flusso principale degli eventi:] \begin{enumerate}
          \item Il docente seleziona lo studente di cui vuole vedere i risultati e le statistiche;
          \item Il docente visualizza i risultati e le statistiche di un questionario della classe eseguito dallo studente selezionato (\hyperlink{UC14.4.1}{UC14.4.1});

      \end{enumerate}
    \item[Postcondizione:] Il docente ha visualizzato i risultati e le statistiche relative allo studente desiderate.
  \end{description}
\hypertarget{UC14.4.1}{}
\subsection{Caso d'uso UC14.4.1: Visualizza risultati questionari dello studente}\begin{description}
\item[Attori:] Docente;
\item[Scopo e descrizione:] Il docente visualizza i risultati e le statistiche relative ai questionari dello studente selezionato
      \item[Precondizione:] Il docente ha selezionato uno studente della classe;

        \item[Flusso principale degli eventi:] \begin{enumerate}
          \item Il docente cerca un questionario tra quelli dello studente nella classe (\hyperlink{UC18}{UC18});
          \item Il docente seleziona il questionario per visualizzarne i risultati e le statistiche;

      \end{enumerate}
    \item[Postcondizione:] Il docente ha visualizzato i risultati e le statistiche relative ai questionari dello studente selezionato.
  \end{description}
\hypertarget{UC15}{}
\subsection{Caso d'uso UC15: Esegui questionario}
        \begin{figure}[H]
            \centering
            \begin{resizedtikzpicture}{\textwidth}
		\umlactor[x=0, y=0]{Studente}
		\begin{umlsystem}[x=0, fill=lightgray!20]{Quizzpedia}
			\umlusecase[x=6, y=2.25, fill=white, width=3cm, name=124]{\textbf{UC15.1:} Rispondi domanda}
			\umlassoc{Studente}{124}
			\umlusecase[x=6, y=0, fill=white, width=3cm, name=160]{\textbf{UC15.10:} Feedback questionario}
			\umlassoc{Studente}{160}
			\umlusecase[x=6, y=-2.25, fill=white, width=3cm, name=132]{\textbf{UC15.3:} Conferma questionario}
			\umlassoc{Studente}{132}
		\end{umlsystem}
            \end{resizedtikzpicture}
            \caption{Caso d'uso UC15: Esegui questionario}
            \label{fig:UC15} 
        \end{figure}
    \begin{description}
\item[Attori:] Studente;
\item[Scopo e descrizione:] Lo studente effettua un questionario
      \item[Precondizione:] Lo studente è autenticato nel sistema;

        \item[Flusso principale degli eventi:] \begin{enumerate}
          \item Lo studente risponde ad una domanda (\hyperlink{UC15.1}{UC15.1});
          \item Lo studente può andare alla domanda successiva;
          \item Lo studente può andare alla domanda precedente;
          \item Lo studente conferma il questionario (\hyperlink{UC15.3}{UC15.3});
          \item Lo studente può lasciare un feedback (\hyperlink{UC15.10}{UC15.10});

      \end{enumerate}
    \item[Postcondizione:] Il questionario è stato compilato e vengono visualizzati i risultati.
  \end{description}
\hypertarget{UC15.1}{}
\subsection{Caso d'uso UC15.1: Rispondi domanda}
        \begin{figure}[H]
            \centering
            \begin{resizedtikzpicture}{\textwidth}
		\umlactor[x=0, y=0]{Studente}
		\begin{umlsystem}[x=0, fill=lightgray!20]{Quizzpedia}
			\umlusecase[x=4, y=0, fill=white, width=3cm, name=124]{\textbf{UC15.1:} Rispondi domanda}
			\umlassoc{Studente}{124}
			\umlusecase[x=11, y=7.6875, fill=white, width=3cm, name=154]{\textbf{UC15.4:} Rispondi domanda vero/falso}
			\umlinherit{154}{124}
			\umlusecase[x=11, y=4.9375, fill=white, width=3cm, name=155]{\textbf{UC15.5:} Rispondi domanda a scelta multipla}
			\umlinherit{155}{124}
			\umlusecase[x=11, y=2.1041666666667, fill=white, width=3cm, name=156]{\textbf{UC15.6:} Rispondi domanda a risposta multipla}
			\umlinherit{156}{124}
			\umlusecase[x=11, y=-1.2291666666667, fill=white, width=3cm, name=157]{\textbf{UC15.7:} Rispondi domanda di tipo testo con parole omesse}
			\umlinherit{157}{124}
			\umlusecase[x=11, y=-4.5625, fill=white, width=3cm, name=158]{\textbf{UC15.8:} Rispondi domanda con associazione di parole}
			\umlinherit{158}{124}
			\umlusecase[x=11, y=-7.6875, fill=white, width=3cm, name=159]{\textbf{UC15.9:} Rispondi domanda a risposta aperta}
			\umlinherit{159}{124}
		\end{umlsystem}
            \end{resizedtikzpicture}
            \caption{Caso d'uso UC15.1: Rispondi domanda}
            \label{fig:UC15.1} 
        \end{figure}
    \begin{description}
\item[Attori:] Studente;
\item[Scopo e descrizione:] Lo studente riporta la risposta alla domanda corrente
      \item[Precondizione:] Lo studente ha iniziato l'esecuzione di un questionario;

        \item[Flusso principale degli eventi:] \begin{enumerate}
          \item Lo studente seleziona la risposta che ritiene esatta con le modalità previste dal tipo di domanda;
          \item Lo studente può lasciare un feedback per dire che la domanda gli piace oppure per segnalare un errore (\hyperlink{UC15.10}{UC15.10});

      \end{enumerate}
    \item[Postcondizione:] Lo studente ha risposto alla domanda.
  \end{description}
\hypertarget{UC15.2}{}
\subsection{Caso d'uso UC15.2: Errore domanda non risposta}\begin{description}
\item[Attori:] Studente;
\item[Scopo e descrizione:] Il sistema non permette allo studente di terminare un questionario in cui non tutte le domande hanno una risposta.
      \item[Precondizione:] Lo studente non ha riposto ad alcune domande nel momento in cui conferma il questionario;

        \item[Flusso principale degli eventi:] \begin{enumerate}
          \item Viene visualizzato un messaggio di errore;

      \end{enumerate}
    \item[Postcondizione:] Il questionario non viene terminato.
  \end{description}
\hypertarget{UC15.3}{}
\subsection{Caso d'uso UC15.3: Conferma questionario}
        \begin{figure}[H]
            \centering
            \begin{resizedtikzpicture}{\textwidth}
		\umlactor[x=0, y=0]{Studente}
		\begin{umlsystem}[x=0, fill=lightgray!20]{Quizzpedia}
			\umlusecase[x=4, y=0, fill=white, width=3cm, name=132]{\textbf{UC15.3:} Conferma questionario}
			\umlassoc{Studente}{132}
			\umlusecase[x=11, y=0, fill=white, width=3cm, name=125]{\textbf{UC15.2:} Errore domanda non risposta}
			\umlextend{125}{132}
		\end{umlsystem}
            \end{resizedtikzpicture}
            \caption{Caso d'uso UC15.3: Conferma questionario}
            \label{fig:UC15.3} 
        \end{figure}
    \begin{description}
\item[Attori:] Studente;
\item[Scopo e descrizione:] Lo studente conferma il questionario eseguito
      \item[Precondizione:] Lo studente è autenticato nel sistema;

        \item[Flusso principale degli eventi:] \begin{enumerate}
          \item Lo studente conferma il questionario;

      \end{enumerate}
    \item[Estensioni:]
      \begin{enumerate}
          \item Se quando lo studente conferma il questionario ci sono domande non risposte viene visualizzato un errore (\hyperlink{UC15.2}{UC15.2});

      \end{enumerate}
    \item[Postcondizione:] Il questionario viene confermato.
  \end{description}
\hypertarget{UC15.4}{}
\subsection{Caso d'uso UC15.4: Rispondi domanda vero/falso}\begin{description}
\item[Attori:] Studente;
\item[Scopo e descrizione:] Lo studente risponde ad una domanda di tipo vero/falso
      \item[Precondizione:] È stata iniziata l'esecuzione di un questionario;

        \item[Flusso principale degli eventi:] \begin{enumerate}
          \item Lo studente seleziona vero o falso;

      \end{enumerate}
    \item[Postcondizione:] Lo studente ha risposto alla domanda di tipo vero/falso.
  \end{description}
\hypertarget{UC15.5}{}
\subsection{Caso d'uso UC15.5: Rispondi domanda a scelta multipla}\begin{description}
\item[Attori:] Studente;
\item[Scopo e descrizione:] Lo studente risponde ad una domanda a scelta multipla
      \item[Precondizione:] È stata iniziata l'esecuzione di un questionario;

        \item[Flusso principale degli eventi:] \begin{enumerate}
          \item Lo studente seleziona la risposta che ritiene corretta;

      \end{enumerate}
    \item[Postcondizione:] Lo studente ha risposto alla domanda a scelta multipla.
  \end{description}
\hypertarget{UC15.6}{}
\subsection{Caso d'uso UC15.6: Rispondi domanda a risposta multipla}\begin{description}
\item[Attori:] Studente;
\item[Scopo e descrizione:] Lo studente risponde ad una domanda a risposta multipla
      \item[Precondizione:] È stata iniziata l'esecuzione di un questionario;

        \item[Flusso principale degli eventi:] \begin{enumerate}
          \item Lo studente e seleziona la o le risposte che ritiene corrette;

      \end{enumerate}
    \item[Postcondizione:] Lo studente ha risposto alla domanda a risposta multipla.
  \end{description}
\hypertarget{UC15.7}{}
\subsection{Caso d'uso UC15.7: Rispondi domanda di tipo testo con parole omesse}\begin{description}
\item[Attori:] Studente;
\item[Scopo e descrizione:] Lo studente risponde ad una domanda di tipo testo con parole omesse
      \item[Precondizione:] È stata iniziata l'esecuzione di un questionario;

        \item[Flusso principale degli eventi:] \begin{enumerate}
          \item Lo studente riempie il testo con parole omesse scegliendole da una lista;

      \end{enumerate}
    \item[Postcondizione:] Lo studente ha risposto alla domanda di tipo testo con parole omesse.
  \end{description}
\hypertarget{UC15.8}{}
\subsection{Caso d'uso UC15.8: Rispondi domanda con associazione di parole}\begin{description}
\item[Attori:] Studente;
\item[Scopo e descrizione:] Lo studente risponde ad una domanda con associazione di parole
      \item[Precondizione:] È stata iniziata l'esecuzione di un questionario;

        \item[Flusso principale degli eventi:] \begin{enumerate}
          \item Lo studente e seleziona le associazioni di parole che ritiene corrette;

      \end{enumerate}
    \item[Postcondizione:] Lo studente ha risposto alla domanda con associazione di parole.
  \end{description}
\hypertarget{UC15.9}{}
\subsection{Caso d'uso UC15.9: Rispondi domanda a risposta aperta}\begin{description}
\item[Attori:] Studente;
\item[Scopo e descrizione:] Lo studente risponde ad una domanda a risposta aperta
      \item[Precondizione:] È stata iniziata l'esecuzione di un questionario;

        \item[Flusso principale degli eventi:] \begin{enumerate}
          \item Lo studente inserisce la risposta come testo;

      \end{enumerate}
    \item[Postcondizione:] Lo studente ha risposto alla domanda a risposta aperta.
  \end{description}
\hypertarget{UC15.10}{}
\subsection{Caso d'uso UC15.10: Feedback questionario}\begin{description}
\item[Attori:] Studente;
\item[Scopo e descrizione:] Lo studente lascia un feedback positivo al questionario
      \item[Precondizione:] Lo studente ha completato il questionario e vede i suoi risultati;

        \item[Flusso principale degli eventi:] \begin{enumerate}
          \item Lo studente può lasciare un feedback positivo al questionario;

      \end{enumerate}
    \item[Postcondizione:] Il feedback dello studente viene memorizzato per il questionario svolto .
  \end{description}
\hypertarget{UC16}{}
\subsection{Caso d'uso UC16: Iscrizione ad una classe}
        \begin{figure}[H]
            \centering
            \begin{resizedtikzpicture}{\textwidth}
		\umlactor[x=0, y=0]{Studente}
		\begin{umlsystem}[x=0, fill=lightgray!20]{Quizzpedia}
			\umlusecase[x=6, y=0, fill=white, width=3cm, name=97]{\textbf{UC16.1:} Inserisci password classe}
			\umlassoc{Studente}{97}
		\end{umlsystem}
            \end{resizedtikzpicture}
            \caption{Caso d'uso UC16: Iscrizione ad una classe}
            \label{fig:UC16} 
        \end{figure}
    \begin{description}
\item[Attori:] Studente;
\item[Scopo e descrizione:] Lo studente si iscrive ad una classe
      \item[Precondizione:] Lo studente è autenticato presso il sistema;

        \item[Flusso principale degli eventi:] \begin{enumerate}
          \item Lo studente cerca una classe (\hyperlink{UC29}{UC29});
          \item Lo studente inserisce la password (che gli deve essere già stata consegnata personalmente dal docente della classe) per iscriversi alla classe (\hyperlink{UC16.1}{UC16.1});
          \item Lo studente conferma l'iscrizione alla classe;

      \end{enumerate}
    \item[Postcondizione:] Lo studente si è iscritto alla classe desiderata.
  \end{description}
\hypertarget{UC16.1}{}
\subsection{Caso d'uso UC16.1: Inserisci password classe}
        \begin{figure}[H]
            \centering
            \begin{resizedtikzpicture}{\textwidth}
		\umlactor[x=0, y=0]{Studente}
		\begin{umlsystem}[x=0, fill=lightgray!20]{Quizzpedia}
			\umlusecase[x=4, y=0, fill=white, width=3cm, name=97]{\textbf{UC16.1:} Inserisci password classe}
			\umlassoc{Studente}{97}
			\umlusecase[x=11, y=0, fill=white, width=3cm, name=98]{\textbf{UC16.2:} Errore password classe}
			\umlextend{98}{97}
		\end{umlsystem}
            \end{resizedtikzpicture}
            \caption{Caso d'uso UC16.1: Inserisci password classe}
            \label{fig:UC16.1} 
        \end{figure}
    \begin{description}
\item[Attori:] Studente;
\item[Scopo e descrizione:] Lo studente inserisce la password per l'iscrizione alla classe
      \item[Precondizione:] Lo studente è autenticato presso il sistema;

        \item[Flusso principale degli eventi:] \begin{enumerate}
          \item Lo studente inserisce la password corretta (che gli deve essere già stata consegnata personalmente dal docente della classe) per iscriversi alla classe;

      \end{enumerate}
    \item[Estensioni:]
      \begin{enumerate}
          \item Se la password inserita non è corretta viene visualizzato un messaggio di errore (\hyperlink{UC16.2}{UC16.2});

      \end{enumerate}
    \item[Postcondizione:] Lo studente ha inserito la password per iscriversi alla classe.
  \end{description}
\hypertarget{UC16.2}{}
\subsection{Caso d'uso UC16.2: Errore password classe}\begin{description}
\item[Attori:] Studente;
\item[Scopo e descrizione:] Il sistema avverte lo studente che ha inserito una password errata per l'iscrizione alla classe
      \item[Precondizione:] Lo studente ha inserito una password errata per iscriversi alla classe;

        \item[Flusso principale degli eventi:] \begin{enumerate}
          \item Viene visualizzato un messaggio di errore;

      \end{enumerate}
    \item[Postcondizione:] Lo studente non viene iscritto alla classe.
  \end{description}
\hypertarget{UC17}{}
\subsection{Caso d'uso UC17: Visualizza storico studente}
        \begin{figure}[H]
            \centering
            \begin{resizedtikzpicture}{\textwidth}
		\umlactor[x=0, y=0]{Studente}
		\begin{umlsystem}[x=0, fill=lightgray!20]{Quizzpedia}
			\umlusecase[x=6, y=3.125, fill=white, width=3cm, name=136]{\textbf{UC17.1:} Visualizza statistiche domande studente}
			\umlassoc{Studente}{136}
			\umlusecase[x=6, y=0, fill=white, width=3cm, name=138]{\textbf{UC17.2:} Visualizza statistiche questionari studente}
			\umlassoc{Studente}{138}
			\umlusecase[x=6, y=-3.125, fill=white, width=3cm, name=139]{\textbf{UC17.3:} Visualizza sommario statistiche studente}
			\umlassoc{Studente}{139}
		\end{umlsystem}
            \end{resizedtikzpicture}
            \caption{Caso d'uso UC17: Visualizza storico studente}
            \label{fig:UC17} 
        \end{figure}
    \begin{description}
\item[Attori:] Studente;
\item[Scopo e descrizione:] Lo studente visualizza il proprio storico
      \item[Precondizione:] Lo studente è autenticato presso il sistema;

        \item[Flusso principale degli eventi:] \begin{enumerate}
          \item Lo studente può visualizzare le sue risposte e le statistiche generali relative alle domande eseguite (\hyperlink{UC17.1}{UC17.1});
          \item Lo studente può visualizzare le sue risposte e le statistiche generali relative ai questionari eseguiti (\hyperlink{UC17.2}{UC17.2});
          \item Lo studente può visualizzare un sommario delle proprie statistiche (\hyperlink{UC17.3}{UC17.3});

      \end{enumerate}
    \item[Postcondizione:] Lo studente ha visualizzato le sue risposte e le statistiche generali relative ai test eseguiti.
  \end{description}
\hypertarget{UC17.1}{}
\subsection{Caso d'uso UC17.1: Visualizza statistiche domande studente}\begin{description}
\item[Attori:] Studente;
\item[Scopo e descrizione:] Lo studente visualizza la percentuale delle proprie risposte corrette e le statistiche generali della domanda
      \item[Precondizione:] Lo studente è autenticato presso il sistema;

        \item[Flusso principale degli eventi:] \begin{enumerate}
          \item Lo studente cerca una domanda tra quelle che ha eseguito in passato (\hyperlink{UC24}{UC24});
          \item Lo studente seleziona la domanda per visualizzarne le risposte date e le statistiche generali;
          \item Lo studente visualizza le risposte già date in passato alla domanda selezionata;
          \item Lo studente visualizza la difficoltà della domanda calcolata come funzione delle risposte date da tutti gli utenti;

      \end{enumerate}
    \item[Postcondizione:] Lo studente ha visualizzato i risultati e le statistiche relative alle domande desiderate.
  \end{description}
\hypertarget{UC17.2}{}
\subsection{Caso d'uso UC17.2: Visualizza statistiche questionari studente}\begin{description}
\item[Attori:] Studente;
\item[Scopo e descrizione:] Lo studente visualizza i punteggi che ha ottenuto nel questionario e le statistiche generali dello stesso
      \item[Precondizione:] Lo studente è autenticato presso il sistema;

        \item[Flusso principale degli eventi:] \begin{enumerate}
          \item Lo studente cerca un questionario tra quelli che ha eseguito in passato (\hyperlink{UC18}{UC18});
          \item Lo studente seleziona il questionario per visualizzarne le risposte date e le statistiche generali;
          \item Lo studente può vedere le statistiche di ogni singola domanda del questionario selezionato (\hyperlink{UC17.1}{UC17.1});
          \item Lo studente vede le proprie risposte al questionario e il punteggio totale;
          \item Lo studente vede la difficoltà del questionario calcolata come funzione delle risposte date da tutti gli utenti;

      \end{enumerate}
    \item[Postcondizione:] Lo studente ha visualizzato i risultati e le statistiche relative ai questionari desiderati.
  \end{description}
\hypertarget{UC17.3}{}
\subsection{Caso d'uso UC17.3: Visualizza sommario statistiche studente}\begin{description}
\item[Attori:] Studente;
\item[Scopo e descrizione:] Lo studente visualizza il totale delle risposte corrette sul totale delle risposte date e la media dei punteggi su tutti i questionari
      \item[Precondizione:] Lo studente è autenticato presso il sistema;

        \item[Flusso principale degli eventi:] \begin{enumerate}
          \item Lo studente visualizza il totale delle domande eseguite e il numero di risposte corrette;
          \item Lo studente visualizza il totale dei questionari eseguiti e la media dei risultati ottenuti;
          \item Lo studente visualizza la media delle difficoltà delle domande a cui ha risposto correttamente;
          \item Lo studente visualizza la media delle difficoltà delle domande a cui non ha risposto correttamente;
          \item Per ogni argomento lo studente visualizza il totale delle domande eseguite, il numero di risposte corrette, la media delle difficoltà delle risposte corrette e la media delle difficoltà delle risposte non corrette;

      \end{enumerate}
    \item[Postcondizione:] Lo studente ha visualizzato i risultati e le statistiche relative alle domande desiderate.
  \end{description}
\hypertarget{UC18}{}
\subsection{Caso d'uso UC18: Ricerca questionario}
        \begin{figure}[H]
            \centering
            \begin{resizedtikzpicture}{\textwidth}
		\umlactor[x=0, y=0]{Utente}
		\begin{umlsystem}[x=0, fill=lightgray!20]{Quizzpedia}
			\umlusecase[x=4, y=0, fill=white, width=3cm, name=17]{\textbf{UC18:} Ricerca questionario}
			\umlassoc{Utente}{17}
			\umlusecase[x=11, y=5.3125, fill=white, width=3cm, name=18]{\textbf{UC19:} Ricerca questionario per titolo}
			\umlinherit{18}{17}
			\umlusecase[x=11, y=2.7708333333333, fill=white, width=3cm, name=19]{\textbf{UC20:} Ricerca questionario per classe}
			\umlinherit{19}{17}
			\umlusecase[x=11, y=0.10416666666667, fill=white, width=3cm, name=20]{\textbf{UC21:} Ricerca questionario per argomento}
			\umlinherit{20}{17}
			\umlusecase[x=11, y=-2.5625, fill=white, width=3cm, name=21]{\textbf{UC22:} Ricerca questionario per docente}
			\umlinherit{21}{17}
			\umlusecase[x=11, y=-5.3125, fill=white, width=3cm, name=22]{\textbf{UC23:} Ricerca questionario per difficoltà}
			\umlinherit{22}{17}
		\end{umlsystem}
            \end{resizedtikzpicture}
            \caption{Caso d'uso UC18: Ricerca questionario}
            \label{fig:UC18} 
        \end{figure}
    \begin{description}
\item[Attori:] Utente;
\item[Scopo e descrizione:] L'utente ricerca un questionario
      \item[Precondizione:] L'utente è autenticato presso il sistema;

        \item[Flusso principale degli eventi:] \begin{enumerate}
          \item L'utente inserisce i dati per la ricerca;

      \end{enumerate}
    \item[Postcondizione:] Il sistema mostra la lista dei questionari che soddisfano la ricerca.
  \end{description}
\hypertarget{UC19}{}
\subsection{Caso d'uso UC19: Ricerca questionario per titolo}\begin{description}
\item[Attori:] Utente;
\item[Scopo e descrizione:] L'utente ricerca un questionario per titolo
      \item[Precondizione:] L'utente è autenticato presso il sistema;

        \item[Flusso principale degli eventi:] \begin{enumerate}
          \item L'utente inserisce il titolo del questionario che vuole cercare;

      \end{enumerate}
    \item[Postcondizione:] Il sistema mostra la lista dei questionari il cui titolo corrisponde al titolo ricercato.
  \end{description}
\hypertarget{UC20}{}
\subsection{Caso d'uso UC20: Ricerca questionario per classe}\begin{description}
\item[Attori:] Utente;
\item[Scopo e descrizione:] L'utente ricerca il questionario per classe
      \item[Precondizione:] L'utente è autenticato presso il sistema;

        \item[Flusso principale degli eventi:] \begin{enumerate}
          \item L'utente seleziona la classe di cui vuole visualizzare i questionari;

      \end{enumerate}
    \item[Postcondizione:] Il sistema mostra la lista dei questionari la cui classe corrisponde alla classe ricercata.
  \end{description}
\hypertarget{UC21}{}
\subsection{Caso d'uso UC21: Ricerca questionario per argomento}\begin{description}
\item[Attori:] Utente;
\item[Scopo e descrizione:] L'utente ricerca questionario per argomento
      \item[Precondizione:] L'utente è autenticato presso il sistema;

        \item[Flusso principale degli eventi:] \begin{enumerate}
          \item L'utente seleziona gli argomento di cui vuole visualizzare i questionari;

      \end{enumerate}
    \item[Postcondizione:] Il sistema mostra la lista dei questionari i cui argomenti corrispondono agli argomenti ricercati.
  \end{description}
\hypertarget{UC22}{}
\subsection{Caso d'uso UC22: Ricerca questionario per docente}\begin{description}
\item[Attori:] Utente;
\item[Scopo e descrizione:] L'utente ricerca questionario per docente
      \item[Precondizione:] L'utente è autenticato presso il sistema;

        \item[Flusso principale degli eventi:] \begin{enumerate}
          \item L'utente inserisce il nome del docente di cui vuole cercare i questionari;

      \end{enumerate}
    \item[Postcondizione:] Il sistema mostra la lista dei questionari il cui docente corrisponde al docente ricercato.
  \end{description}
\hypertarget{UC23}{}
\subsection{Caso d'uso UC23: Ricerca questionario per difficoltà}\begin{description}
\item[Attori:] Utente;
\item[Scopo e descrizione:] L'utente ricerca un questionario per difficoltà
      \item[Precondizione:] L'utente è autenticato presso il sistema;

        \item[Flusso principale degli eventi:] \begin{enumerate}
          \item L'utente inserisce un limite inferiore e superiore di difficoltà di cui vuole cercare questionari;

      \end{enumerate}
    \item[Postcondizione:] Il sistema mostra la lista dei questionari la cui difficoltà corrisponde alla difficoltà ricercata.
  \end{description}
\hypertarget{UC24}{}
\subsection{Caso d'uso UC24: Ricerca domanda}
        \begin{figure}[H]
            \centering
            \begin{resizedtikzpicture}{\textwidth}
		\umlactor[x=0, y=0]{Utente}
		\begin{umlsystem}[x=0, fill=lightgray!20]{Quizzpedia}
			\umlusecase[x=4, y=0, fill=white, width=3cm, name=23]{\textbf{UC24:} Ricerca domanda}
			\umlassoc{Utente}{23}
			\umlusecase[x=11, y=3.7708333333333, fill=white, width=3cm, name=24]{\textbf{UC25:} Ricerca domanda per keywords}
			\umlinherit{24}{23}
			\umlusecase[x=11, y=1.3125, fill=white, width=3cm, name=25]{\textbf{UC26:} Ricerca domanda per argomento}
			\umlinherit{25}{23}
			\umlusecase[x=11, y=-1.2291666666667, fill=white, width=3cm, name=26]{\textbf{UC27:} Ricerca domanda per difficoltà}
			\umlinherit{26}{23}
			\umlusecase[x=11, y=-3.7708333333333, fill=white, width=3cm, name=27]{\textbf{UC28:} Ricerca domanda per docente}
			\umlinherit{27}{23}
		\end{umlsystem}
            \end{resizedtikzpicture}
            \caption{Caso d'uso UC24: Ricerca domanda}
            \label{fig:UC24} 
        \end{figure}
    \begin{description}
\item[Attori:] Utente;
\item[Scopo e descrizione:] L'utente ricerca una domanda
      \item[Precondizione:] L'utente è autenticato presso il sistema;

        \item[Flusso principale degli eventi:] \begin{enumerate}
          \item L'utente inserisce i dati per la ricerca	;

      \end{enumerate}
    \item[Postcondizione:] Il sistema mostra la lista delle domande che soddisfano la ricerca.
  \end{description}
\hypertarget{UC25}{}
\subsection{Caso d'uso UC25: Ricerca domanda per keywords}\begin{description}
\item[Attori:] Utente;
\item[Scopo e descrizione:] L'utente ricerca una domanda per keywords

      \item[Precondizione:] L'utente è autenticato presso il sistema
;

        \item[Flusso principale degli eventi:] \begin{enumerate}
          \item L'utente inserisce le keywords presenti nel testo delle domande che vuole cercare;

      \end{enumerate}
    \item[Postcondizione:] Il sistema mostra la lista delle domande che contengono nel titolo o nel corpo le parole chiave specificate.
  \end{description}
\hypertarget{UC26}{}
\subsection{Caso d'uso UC26: Ricerca domanda per argomento}\begin{description}
\item[Attori:] Utente;
\item[Scopo e descrizione:] L'utente ricerca una domanda per argomento
      \item[Precondizione:] L'utente è autenticato presso il sistema
;

        \item[Flusso principale degli eventi:] \begin{enumerate}
          \item L'utente seleziona gli argomenti di cui vuole visualizzare le domande;

      \end{enumerate}
    \item[Postcondizione:] Il sistema mostra la lista delle domande che contengono gli argomenti selezionati.
  \end{description}
\hypertarget{UC27}{}
\subsection{Caso d'uso UC27: Ricerca domanda per difficoltà}\begin{description}
\item[Attori:] Utente;
\item[Scopo e descrizione:] L'utente ricerca una domanda per difficoltà

      \item[Precondizione:] L'utente è autenticato presso il sistema
;

        \item[Flusso principale degli eventi:] \begin{enumerate}
          \item L'utente inserisce il limite minimo e massimo di difficoltà di cui vuole visualizzare le domande;

      \end{enumerate}
    \item[Postcondizione:] Il sistema mostra la lista delle domande con la difficoltà selezionata.
  \end{description}
\hypertarget{UC28}{}
\subsection{Caso d'uso UC28: Ricerca domanda per docente}\begin{description}
\item[Attori:] Utente;
\item[Scopo e descrizione:] L'utente ricerca una domanda per docente
      \item[Precondizione:] L'utente è autenticato presso il sistema
;

        \item[Flusso principale degli eventi:] \begin{enumerate}
          \item L'utente inserisce il nome del docente di cui vuole visualizzare le domande;

      \end{enumerate}
    \item[Postcondizione:] Il sistema mostra la lista delle domande create dal docente selezionato.
  \end{description}
\hypertarget{UC29}{}
\subsection{Caso d'uso UC29: Ricerca classe}
        \begin{figure}[H]
            \centering
            \begin{resizedtikzpicture}{\textwidth}
		\umlactor[x=0, y=0]{Studente}
		\begin{umlsystem}[x=0, fill=lightgray!20]{Quizzpedia}
			\umlusecase[x=4, y=0, fill=white, width=3cm, name=28]{\textbf{UC29:} Ricerca classe}
			\umlassoc{Studente}{28}
			\umlusecase[x=11, y=1.2083333333333, fill=white, width=3cm, name=29]{\textbf{UC30:} Ricerca classe per docente}
			\umlinherit{29}{28}
			\umlusecase[x=11, y=-1.2083333333333, fill=white, width=3cm, name=30]{\textbf{UC31:} Ricerca classe per argomento}
			\umlinherit{30}{28}
		\end{umlsystem}
            \end{resizedtikzpicture}
            \caption{Caso d'uso UC29: Ricerca classe}
            \label{fig:UC29} 
        \end{figure}
    \begin{description}
\item[Attori:] Studente;
\item[Scopo e descrizione:] Lo studente ricerca una classe
      \item[Precondizione:] Lo studente è autenticato presso il sistema;

        \item[Flusso principale degli eventi:] \begin{enumerate}
          \item Lo studente inserisce i dati per la ricerca;

      \end{enumerate}
    \item[Postcondizione:] Il sistema mostra la lista delle classi che soddisfano la ricerca.
  \end{description}
\hypertarget{UC30}{}
\subsection{Caso d'uso UC30: Ricerca classe per docente}\begin{description}
\item[Attori:] Studente;
\item[Scopo e descrizione:] Lo studente ricerca una classe per docente
      \item[Precondizione:] Lo studente è autenticato presso il sistema;

        \item[Flusso principale degli eventi:] \begin{enumerate}
          \item Lo studente inserisce il nome del docente di cui le classi;

      \end{enumerate}
    \item[Postcondizione:] Il sistema mostra la lista delle classi che soddisfano la ricerca in base al docente selezionato.
  \end{description}
\hypertarget{UC31}{}
\subsection{Caso d'uso UC31: Ricerca classe per argomento}\begin{description}
\item[Attori:] Studente;
\item[Scopo e descrizione:] Lo studente ricerca una classe per argomento
      \item[Precondizione:] Lo studente è autenticato presso il sistema;

        \item[Flusso principale degli eventi:] \begin{enumerate}
          \item Lo studente seleziona gli argomenti di cui vuole visualizzare le classi;

      \end{enumerate}
    \item[Postcondizione:] Il sistema mostra la lista delle classi che soddisfano la ricerca in base agli argomenti selezionati.
  \end{description}
\hypertarget{UC32}{}
\subsection{Caso d'uso UC32: Azioni Amministratore}
        \begin{figure}[H]
            \centering
            \begin{resizedtikzpicture}{\textwidth}
		\umlactor[x=0, y=0]{Amministratore}
		\begin{umlsystem}[x=0, fill=lightgray!20]{Quizzpedia}
			\umlusecase[x=6, y=1, fill=white, width=3cm, name=34]{\textbf{UC32.1:} Cambia ruolo}
			\umlassoc{Amministratore}{34}
			\umlusecase[x=6, y=-1, fill=white, width=3cm, name=35]{\textbf{UC32.2:} Rimuovi  docente}
			\umlassoc{Amministratore}{35}
		\end{umlsystem}
            \end{resizedtikzpicture}
            \caption{Caso d'uso UC32: Azioni Amministratore}
            \label{fig:UC32} 
        \end{figure}
    \begin{description}
\item[Attori:] Amministratore;
\item[Scopo e descrizione:] Descrive le azione che può compiere un aministratore
      \item[Precondizione:] L'amministratore è autenticato nel sistema;

        \item[Flusso principale degli eventi:] \begin{enumerate}
          \item L'amministratore può impostare il ruolo di uno studente a docente o viceversa (\hyperlink{UC32.1}{UC32.1});
          \item L'amministratore può rimuovere un docente o studente (\hyperlink{UC32.2}{UC32.2});

      \end{enumerate}
    \item[Postcondizione:] Il sistema ha ottenuto le informazioni sulle operazioni che l’amministratore desidera eseguire.
  \end{description}
\hypertarget{UC32.1}{}
\subsection{Caso d'uso UC32.1: Cambia ruolo}\begin{description}
\item[Attori:] Amministratore;
\item[Scopo e descrizione:] L'amministratore imposta il ruolo dello studente a docente o viceversa
      \item[Precondizione:] L'utente selezionato è uno studente o docente;

        \item[Flusso principale degli eventi:] \begin{enumerate}
          \item L'amministratore seleziona uno studente;
          \item L'amministratore imposta il ruolo dello studente a docente;
          \item L'amministratore conferma l'operazione da eseguire;

      \end{enumerate}
    \item[Postcondizione:] Il ruolo dell'utente è diventato quello selezionato.
  \end{description}
\hypertarget{UC32.2}{}
\subsection{Caso d'uso UC32.2: Rimuovi  docente}\begin{description}
\item[Attori:] Amministratore;
\item[Scopo e descrizione:] L'amministratore rimuove un docente dal sistema
      \item[Precondizione:] Il docente da rimuovere deve essere presente nel sistema;

        \item[Flusso principale degli eventi:] \begin{enumerate}
          \item Amministratore ricerca il docente da rimuovere;
          \item Amministratore seleziona il docente da rimuovere;
          \item Amministratore conferma la rimozione del docente selezionato;

      \end{enumerate}
    \item[Postcondizione:] Il medesimo docente è stato rimosso dal sistema.
  \end{description}
\hypertarget{UC33}{}
\subsection{Caso d'uso UC33: Azioni proprietario}
        \begin{figure}[H]
            \centering
            \begin{resizedtikzpicture}{\textwidth}
		\umlactor[x=0, y=0]{Proprietario}
		\begin{umlsystem}[x=0, fill=lightgray!20]{Quizzpedia}
			\umlusecase[x=6, y=1.1458333333333, fill=white, width=3cm, name=61]{\textbf{UC33.1:} Aggiungi amministratore}
			\umlassoc{Proprietario}{61}
			\umlusecase[x=6, y=-1.1458333333333, fill=white, width=3cm, name=102]{\textbf{UC33.2:} Elimina amministratore}
			\umlassoc{Proprietario}{102}
		\end{umlsystem}
            \end{resizedtikzpicture}
            \caption{Caso d'uso UC33: Azioni proprietario}
            \label{fig:UC33} 
        \end{figure}
    \begin{description}
\item[Attori:] Proprietario;
\item[Scopo e descrizione:] Descrive le azioni che può compiere un proprietario
      \item[Precondizione:] Il proprietario è autenticato presso il sistema;

        \item[Flusso principale degli eventi:] \begin{enumerate}
          \item Il proprietario può impostare il ruolo di un docente o studente ad amministratore (\hyperlink{UC33.1}{UC33.1});
          \item Il proprietario può rimuovere un amministratore (\hyperlink{UC33.2}{UC33.2});

      \end{enumerate}
    \item[Postcondizione:] Il sistema ha ottenuto le informazioni sulle operazioni che il proprietario desidera eseguire.
  \end{description}
\hypertarget{UC33.1}{}
\subsection{Caso d'uso UC33.1: Aggiungi amministratore}\begin{description}
\item[Attori:] Proprietario;
\item[Scopo e descrizione:] Il proprietario invita un nuovo amministratore nel sistema inserendone la mail
      \item[Precondizione:] Il proprietario è autenticato nel sistema;

        \item[Flusso principale degli eventi:] \begin{enumerate}
          \item Il proprietario inserisce mail del nuovo amministratore;
          \item Il proprietario seleziona la funzionalità di aggiunta amministratore;

      \end{enumerate}
    \item[Postcondizione:] È stato inserito un nuovo account amministratore.
  \end{description}
\hypertarget{UC33.2}{}
\subsection{Caso d'uso UC33.2: Elimina amministratore}\begin{description}
\item[Attori:] Proprietario;
\item[Scopo e descrizione:] Il proprietario sceglie l'amministratore da eliminare e conferma l'eliminazione
      \item[Precondizione:] Il proprietario è autenticato nel sistema;

        \item[Flusso principale degli eventi:] \begin{enumerate}
          \item Il proprietario seleziona l'account amministratore da eliminare;
          \item Il proprietario seleziona la funzionalità di eliminazione amministratore;
          \item Il proprietario conferma l'eliminazione dell'amministratore;

      \end{enumerate}
    \item[Postcondizione:] È stato eliminato l'amministratore selezionato.
  \end{description}


\newpage
\section{Requisiti}
\subsection{Struttura di un requisito}
È compito degli \ANpl\ stilare una lista dei requisiti emersi dal capitolato e da eventuali riunioni con il proponente. Questi requisiti dovranno essere classificati per tipo e importanza, utilizzando la seguente struttura:
\begin{center}
	R[importanza][tipo][codice]
\end{center}
\begin{itemize}
	\item \textbf{Importanza} può assumere i seguenti valori:
	\begin{description}
		\item[1:] requisito desiderabile
		\item[2:] requisito opzionale
		\item[3:] requisito obbligatorio
	\end{description}
	\item \textbf{Tipo} può assumere i seguenti valori:
	\begin{description}
		\item[F:] Funzionale, descrive i servizi o le funzioni offerte dal sistema
		\item[Q:] Di Qualità, descrive i requisiti sulla qualità offerte dal sistema
		\item[P:] Prestazionale, descrive i requisiti sulle prestazioni offerte dal sistema
		\item[V:] Vincolo, descrive i vincoli sui servizi offerti dal sistema
	\end{description}
	\item \textbf{Codice} è un numero progressivo univoco per ogni requisito, indipendente da importanza e tipo. Nel caso si abbia un sotto-requisito codice può anche essere espresso in modo gerarchico tramite la notazione:
	\begin{center}
		\textit{CodiceRequistoPadre.CodiceSottorequisito}
	\end{center}
	\item \textbf{Validazione:} verrà inserito un link al metodo deciso per validare il requisito
\end{itemize}
Ogni requisito deve essere correlato da una sintetica ma precisa descrizione. Per ogni requisito bisogna indicarne le fonti, che posso essere il capitolato o uno o più casi d'uso.
\newpage
\subsection{Tabella dei requisiti/fonti}

% Output di tracy sotto a:
% tracy-tpiga.rhcloud.com/tracy/latex/RequirementSource?type=latex&raw=1
\begin{longtable}{|r l|p{2cm}|p{6cm}|p{2cm}|}
\hline
\multicolumn{2}{|c|}{Requisito} & Tipologia & Descrizione & Fonti\tabularnewline
\hline
 & \hypertarget{R-3V1}{R-3V1} & Vincolo

Obbligatorio & La parte destinata ai creatori di domande e di questionari dovrà essere fruibile attraverso un PC & Capitolato\tabularnewline
\hline
 & \hypertarget{R-3V2}{R-3V2} & Vincolo

Obbligatorio & La parte server deve essere realizzata in Javascript con server node.js
 & Capitolato\tabularnewline
\hline
 & \hypertarget{R-3V3}{R-3V3} & Vincolo

Obbligatorio & La parte client dovrà essere eseguibile in un browser HTML5 & Capitolato\tabularnewline
\hline
\begin{tikzpicture}
\draw [->, thick] (0.2,0.2) -- (0.2,0.1) -- (1,0.1);
\end{tikzpicture} & \hypertarget{R-3V3.1}{R-3V3.1} & Vincolo

Obbligatorio & L'applicazione funziona correttamente su Chrome versione 47.0.2526 e successive & Capitolato\tabularnewline
\hline
\begin{tikzpicture}
\draw [->, thick] (0.2,0.2) -- (0.2,0.1) -- (1,0.1);
\end{tikzpicture} & \hypertarget{R-3V3.2}{R-3V3.2} & Vincolo

Obbligatorio & L'applicazione funziona correttamente su Firefox versione 38.5.2esr e successive & Capitolato\tabularnewline
\hline
\begin{tikzpicture}
\draw [->, thick] (0.2,0.2) -- (0.2,0.1) -- (1,0.1);
\end{tikzpicture} & \hypertarget{R-3V3.3}{R-3V3.3} & Vincolo

Obbligatorio & L'applicazione funziona correttamente su Microsoft Edge versione 25.10586 e successive & Capitolato\tabularnewline
\hline
\begin{tikzpicture}
\draw [->, thick] (0.2,0.2) -- (0.2,0.1) -- (1,0.1);
\end{tikzpicture} & \hypertarget{R-3V3.4}{R-3V3.4} & Vincolo

Obbligatorio & L'applicazione funziona correttamente su Safari versione 9.0.2 e successive & Capitolato\tabularnewline
\hline
 & \hypertarget{R-3V4}{R-3V4} & Vincolo

Obbligatorio & L'applicazione deve utilizzare fogli stile CSS3 per l’ aspetto estetico & Capitolato\tabularnewline
\hline
 & \hypertarget{R-3V5}{R-3V5} & Vincolo

Obbligatorio & L'applicazione deve utilizzare Javascript per la parte attiva del client & Capitolato\tabularnewline
\hline
 & \hypertarget{R-3V6}{R-3V6} & Vincolo

Obbligatorio & La parte destinata agli esaminandi dovrà funzionare con qualunque dispositivo: dal Personal Computer, ai Tablet ed agli SmartPhone attraverso browser & Capitolato\tabularnewline
\hline
\begin{tikzpicture}
\draw [->, thick] (0.2,0.2) -- (0.2,0.1) -- (1,0.1);
\end{tikzpicture} & \hypertarget{R-3V6.1}{R-3V6.1} & Vincolo

Obbligatorio & La parte destinata agli esaminandi dovrà funzionare si smartphone e tablet iOS dalla versione iOS6 & Capitolato\tabularnewline
\hline
\begin{tikzpicture}
\draw [->, thick] (0.2,0.2) -- (0.2,0.1) -- (1,0.1);
\end{tikzpicture} & \hypertarget{R-3V6.2}{R-3V6.2} & Vincolo

Obbligatorio & La parte destinata agli esaminandi dovrà funzionare su Android dalla versione 4.0 & Capitolato\tabularnewline
\hline
 & \hypertarget{R-3F7}{R-3F7} & Funzionale

Obbligatorio & Il sistema deve gestire questionari & Capitolato

\hyperlink{UC6}{UC6}\tabularnewline
\hline
\begin{tikzpicture}
\draw [->, thick] (0.2,0.2) -- (0.2,0.1) -- (1,0.1);
\end{tikzpicture} & \hypertarget{R-3F7.1}{R-3F7.1} & Funzionale

Obbligatorio & L'applicazione deve archiviare questionari in un server, suddivisi per argomento
 & Capitolato\tabularnewline
\hline
\begin{tikzpicture}
\draw [->, thick] (0.2,0.2) -- (0.2,0.1) -- (1,0.1);
\end{tikzpicture} & \hypertarget{R-3F7.2}{R-3F7.2} & Funzionale

Obbligatorio & L'applicazione deve tradurre da QML a HTML le domande archiviate & Capitolato

\hyperlink{UC13.1}{UC13.1}

\hyperlink{UC5.12}{UC5.12}

\hyperlink{UC5}{UC5}

\hyperlink{UC6}{UC6}\tabularnewline
\hline
\begin{tikzpicture}
\draw [->, thick] (0.2,0.2) -- (0.2,0.1) -- (1,0.1);
\end{tikzpicture} & \hypertarget{R-3F7.3}{R-3F7.3} & Funzionale

Obbligatorio & L'applicazione deve proporre questionari preconfezionati & Capitolato

\hyperlink{UC6}{UC6}\tabularnewline
\hline
\begin{tikzpicture}
\draw [->, thick] (0.2,0.2) -- (0.2,0.1) -- (1,0.1);
\end{tikzpicture} & \hypertarget{R-3F7.4}{R-3F7.4} & Funzionale

Obbligatorio & L'applicazione deve valutare le risposte fornite dall'utente & Capitolato

\hyperlink{UC13.11}{UC13.11}

\hyperlink{UC13}{UC13}

\hyperlink{UC15}{UC15}\tabularnewline
\hline
\begin{tikzpicture}
\draw [->, thick] (0.2,0.2) -- (0.2,0.1) -- (1,0.1);
\end{tikzpicture} & \hypertarget{R-3F7.5}{R-3F7.5} & Funzionale

Obbligatorio & Le domande dovranno essere raccolte attraverso uno specifico linguaggio, che chiameremo Quiz Markup Language (QML), studiato per descrivere quiz
 & Capitolato

\hyperlink{UC5.1.2}{UC5.1.2}\tabularnewline
\hline
\begin{tikzpicture}
\draw [->, thick] (0.4,0.2) -- (0.4,0.1) -- (1,0.1);
\end{tikzpicture} & \hypertarget{R-3F7.5.1}{R-3F7.5.1} & Funzionale

Obbligatorio & Il QML deve essere un linguaggio semplice con un parsing ad hoc definito dal gruppo. La sintassi per scrivere il testo della domanda deve essere unica, la sintassi per le risposte è dipendente dal tipo di domanda & Capitolato

\hyperlink{UC5.1.2}{UC5.1.2}\tabularnewline
\hline
\begin{tikzpicture}
\draw [->, thick] (0.6,0.2) -- (0.6,0.1) -- (1,0.1);
\end{tikzpicture} & \hypertarget{R-2F7.5.1.1}{R-2F7.5.1.1} & Funzionale

Opzionale & Il QML deve permettere l'inserimento di testo in grassetto & \hyperlink{UC5.1.2}{UC5.1.2}\tabularnewline
\hline
\begin{tikzpicture}
\draw [->, thick] (0.6,0.2) -- (0.6,0.1) -- (1,0.1);
\end{tikzpicture} & \hypertarget{R-3F7.5.1.2}{R-3F7.5.1.2} & Funzionale

Obbligatorio & Il QML deve permettere l'inserimento di immagini & Capitolato

\hyperlink{UC5.1.2}{UC5.1.2}\tabularnewline
\hline
\begin{tikzpicture}
\draw [->, thick] (0.6,0.2) -- (0.6,0.1) -- (1,0.1);
\end{tikzpicture} & \hypertarget{R-2F7.5.1.3}{R-2F7.5.1.3} & Funzionale

Opzionale & Il QML deve permettere l'inserimento di testo in corsivo & \hyperlink{UC5.1.2}{UC5.1.2}\tabularnewline
\hline
\begin{tikzpicture}
\draw [->, thick] (0.6,0.2) -- (0.6,0.1) -- (1,0.1);
\end{tikzpicture} & \hypertarget{R-2F7.5.1.4}{R-2F7.5.1.4} & Funzionale

Opzionale & Il QML deve permettere l'inserimento di tabelle & \hyperlink{UC5.1.2}{UC5.1.2}\tabularnewline
\hline
\begin{tikzpicture}
\draw [->, thick] (0.6,0.2) -- (0.6,0.1) -- (1,0.1);
\end{tikzpicture} & \hypertarget{R-2F7.5.1.5}{R-2F7.5.1.5} & Funzionale

Opzionale & Il QML deve permettere di definire un punteggio per ogni risposta & \hyperlink{UC5.1.2}{UC5.1.2}\tabularnewline
\hline
\begin{tikzpicture}
\draw [->, thick] (0.6,0.2) -- (0.6,0.1) -- (1,0.1);
\end{tikzpicture} & \hypertarget{R-3F7.5.1.6}{R-3F7.5.1.6} & Funzionale

Obbligatorio & Il QML deve permettere di definire il tipo di domanda & \hyperlink{UC5.1.2}{UC5.1.2}\tabularnewline
\hline
\begin{tikzpicture}
\draw [->, thick] (0.6,0.2) -- (0.6,0.1) -- (1,0.1);
\end{tikzpicture} & \hypertarget{R-2F7.5.1.7}{R-2F7.5.1.7} & Funzionale

Opzionale & Il QML deve permettere l'inserimento di video & \hyperlink{UC5.1.2}{UC5.1.2}

\hyperlink{UC5.1.3}{UC5.1.3}\tabularnewline
\hline
\begin{tikzpicture}
\draw [->, thick] (0.6,0.2) -- (0.6,0.1) -- (1,0.1);
\end{tikzpicture} & \hypertarget{R-2F7.5.1.8}{R-2F7.5.1.8} & Funzionale

Opzionale & Il QML deve permettere l'inserimento di audio & \hyperlink{UC5.1.2}{UC5.1.2}

\hyperlink{UC5.1.3}{UC5.1.3}\tabularnewline
\hline
\begin{tikzpicture}
\draw [->, thick] (0.4,0.2) -- (0.4,0.1) -- (1,0.1);
\end{tikzpicture} & \hypertarget{R-2F7.5.2}{R-2F7.5.2} & Funzionale

Opzionale & Il QML può gestire risposte a scelta multipla, a riempimento di parole omesse, ad associazione, a risposta aperta, con testi e immagini & Capitolato

\hyperlink{UC5.10}{UC5.10}

\hyperlink{UC5.11}{UC5.11}

\hyperlink{UC5.1}{UC5.1}

\hyperlink{UC5.7}{UC5.7}

\hyperlink{UC5.9}{UC5.9}\tabularnewline
\hline
\begin{tikzpicture}
\draw [->, thick] (0.4,0.2) -- (0.4,0.1) -- (1,0.1);
\end{tikzpicture} & \hypertarget{R-3F7.5.3}{R-3F7.5.3} & Funzionale

Obbligatorio & Le domande scritte in QML devono essere editabili dall'utente & Proponente

\hyperlink{UC5.1}{UC5.1}\tabularnewline
\hline
\begin{tikzpicture}
\draw [->, thick] (0.4,0.2) -- (0.4,0.1) -- (1,0.1);
\end{tikzpicture} & \hypertarget{R-2F7.5.4}{R-2F7.5.4} & Funzionale

Opzionale & Le domande possono essere inserite e editabili attraverso interfaccia grafica & Proponente

\hyperlink{UC5.1.3}{UC5.1.3}\tabularnewline
\hline
\begin{tikzpicture}
\draw [->, thick] (0.4,0.2) -- (0.4,0.1) -- (1,0.1);
\end{tikzpicture} & \hypertarget{R-3F7.5.5}{R-3F7.5.5} & Funzionale

Obbligatorio & Il QML deve gestire risposte vero/falso e risposte a scelta multipla con testi e immagini & Capitolato

\hyperlink{UC5.1}{UC5.1}

\hyperlink{UC5.6}{UC5.6}

\hyperlink{UC5.8}{UC5.8}\tabularnewline
\hline
\begin{tikzpicture}
\draw [->, thick] (0.2,0.2) -- (0.2,0.1) -- (1,0.1);
\end{tikzpicture} & \hypertarget{R-2F7.6}{R-2F7.6} & Funzionale

Opzionale & Il sistema deve gestire sia questionari pubblici che privati attraverso la creazione di classi la cui iscrizione è protetta da password & Proponente

\hyperlink{UC14}{UC14}

\hyperlink{UC7.1}{UC7.1}\tabularnewline
\hline
\begin{tikzpicture}
\draw [->, thick] (0.2,0.2) -- (0.2,0.1) -- (1,0.1);
\end{tikzpicture} & \hypertarget{R-3F7.7}{R-3F7.7} & Funzionale

Obbligatorio & I docenti costruiranno questionari scegliendo tra le domande archiviate nel server & Capitolato

\hyperlink{UC6.1}{UC6.1}

\hyperlink{UC6}{UC6}\tabularnewline
\hline
\begin{tikzpicture}
\draw [->, thick] (0.4,0.2) -- (0.4,0.1) -- (1,0.1);
\end{tikzpicture} & \hypertarget{R-3F7.7.1}{R-3F7.7.1} & Funzionale

Obbligatorio & Il sistema deve segnalare un errore nel caso un questionario creato non contenga domande & \hyperlink{UC6.1}{UC6.1}

\hyperlink{UC6.4}{UC6.4}

\hyperlink{UC6}{UC6}\tabularnewline
\hline
\begin{tikzpicture}
\draw [->, thick] (0.2,0.2) -- (0.2,0.1) -- (1,0.1);
\end{tikzpicture} & \hypertarget{R-2F7.8}{R-2F7.8} & Funzionale

Opzionale & Il sistema deve offrire la possibilità di creare dinamicamente i questionari & Capitolato\tabularnewline
\hline
\begin{tikzpicture}
\draw [->, thick] (0.2,0.2) -- (0.2,0.1) -- (1,0.1);
\end{tikzpicture} & \hypertarget{R-2F7.9}{R-2F7.9} & Funzionale

Opzionale & L'applicazione deve essere predisposta ad un sistema di feedback per i questionari & Capitolato

\hyperlink{UC13.10}{UC13.10}\tabularnewline
\hline
\begin{tikzpicture}
\draw [->, thick] (0.2,0.2) -- (0.2,0.1) -- (1,0.1);
\end{tikzpicture} & \hypertarget{R-2F7.10}{R-2F7.10} & Funzionale

Opzionale & Il sistema deve permettere la raccolta dei dati relativi a questionari svolti & Capitolato

\hyperlink{UC15}{UC15}\tabularnewline
\hline
\begin{tikzpicture}
\draw [->, thick] (0.4,0.2) -- (0.4,0.1) -- (1,0.1);
\end{tikzpicture} & \hypertarget{R-2F7.10.1}{R-2F7.10.1} & Funzionale

Opzionale & Il sistema deve permettere la rivisitazione del test svolto e dei risultati ottenuti & Capitolato

\hyperlink{UC15}{UC15}\tabularnewline
\hline
\begin{tikzpicture}
\draw [->, thick] (0.4,0.2) -- (0.4,0.1) -- (1,0.1);
\end{tikzpicture} & \hypertarget{R-2F7.10.2}{R-2F7.10.2} & Funzionale

Opzionale & Il sistema deve offrire la possibilità di analizzare i dati raccolti dai test effettuati & Capitolato

\hyperlink{UC15.2}{UC15.2}

\hyperlink{UC15}{UC15}\tabularnewline
\hline
\begin{tikzpicture}
\draw [->, thick] (0.4,0.2) -- (0.4,0.1) -- (1,0.1);
\end{tikzpicture} & \hypertarget{R-2F7.10.3}{R-2F7.10.3} & Funzionale

Opzionale & Il sistema deve offrire la possibilità di visualizzazione sotto forma di grafico delle statistiche dei dati raccolti & Capitolato

\hyperlink{UC15.3}{UC15.3}

\hyperlink{UC15}{UC15}\tabularnewline
\hline
\begin{tikzpicture}
\draw [->, thick] (0.2,0.2) -- (0.2,0.1) -- (1,0.1);
\end{tikzpicture} & \hypertarget{R-3F7.11}{R-3F7.11} & Funzionale

Obbligatorio & Il sistema deve prevedere la gestione di domande & \hyperlink{UC5}{UC5}

\hyperlink{UC6.1.1}{UC6.1.1}

\hyperlink{UC6.1.2}{UC6.1.2}

\hyperlink{UC6}{UC6}\tabularnewline
\hline
\begin{tikzpicture}
\draw [->, thick] (0.4,0.2) -- (0.4,0.1) -- (1,0.1);
\end{tikzpicture} & \hypertarget{R-3F7.11.1}{R-3F7.11.1} & Funzionale

Obbligatorio & Il sistema deve consentire ad un docente di inserire una nuova domanda & \hyperlink{UC5.1}{UC5.1}

\hyperlink{UC5}{UC5}\tabularnewline
\hline
\begin{tikzpicture}
\draw [->, thick] (0.6,0.2) -- (0.6,0.1) -- (1,0.1);
\end{tikzpicture} & \hypertarget{R-3F7.11.1.1}{R-3F7.11.1.1} & Funzionale

Obbligatorio & Un docente deve poter selezionare gli argomenti di una domanda & \hyperlink{UC5.1.1}{UC5.1.1}

\hyperlink{UC5.1}{UC5.1}

\hyperlink{UC5}{UC5}\tabularnewline
\hline
\begin{tikzpicture}
\draw [->, thick] (0.8,0.2) -- (0.8,0.1) -- (1,0.1);
\end{tikzpicture} & \hypertarget{R-3F7.11.1.1.1}{R-3F7.11.1.1.1} & Funzionale

Obbligatorio & Il sistema deve segnalare un errore nel caso non venga selezionato alcun argomento per una nuova domanda & \hyperlink{UC5.1.1}{UC5.1.1}

\hyperlink{UC5.5}{UC5.5}

\hyperlink{UC5}{UC5}

\hyperlink{UC6.1.1}{UC6.1.1}\tabularnewline
\hline
\begin{tikzpicture}
\draw [->, thick] (0.6,0.2) -- (0.6,0.1) -- (1,0.1);
\end{tikzpicture} & \hypertarget{R-3F7.11.1.2}{R-3F7.11.1.2} & Funzionale

Obbligatorio & Un docente deve comporre la domanda da inserire in QML & \hyperlink{UC5.1.2}{UC5.1.2}

\hyperlink{UC5}{UC5}

\hyperlink{UC6.1.1}{UC6.1.1}\tabularnewline
\hline
\begin{tikzpicture}
\draw [->, thick] (0.8,0.2) -- (0.8,0.1) -- (1,0.1);
\end{tikzpicture} & \hypertarget{R-3F7.11.1.2.1}{R-3F7.11.1.2.1} & Funzionale

Obbligatorio & Il sistema deve segnalare un errore nel caso di inserimento QML non valido & \hyperlink{UC5.1.2}{UC5.1.2}

\hyperlink{UC5.4}{UC5.4}

\hyperlink{UC5}{UC5}

\hyperlink{UC6.1.1}{UC6.1.1}\tabularnewline
\hline
\begin{tikzpicture}
\draw [->, thick] (0.8,0.2) -- (0.8,0.1) -- (1,0.1);
\end{tikzpicture} & \hypertarget{R-3F7.11.1.2.2}{R-3F7.11.1.2.2} & Funzionale

Obbligatorio & Il docente deve poter inserire domande di tipo vero/falso & \hyperlink{UC5.1}{UC5.1}

\hyperlink{UC5.6}{UC5.6}\tabularnewline
\hline
\begin{tikzpicture}
\draw [->, thick] (0.8,0.2) -- (0.8,0.1) -- (1,0.1);
\end{tikzpicture} & \hypertarget{R-3F7.11.1.2.3}{R-3F7.11.1.2.3} & Funzionale

Obbligatorio & Il docente deve poter inserire domande a scelta multipla & \hyperlink{UC5.1}{UC5.1}

\hyperlink{UC5.7}{UC5.7}\tabularnewline
\hline
\begin{tikzpicture}
\draw [->, thick] (0.8,0.2) -- (0.8,0.1) -- (1,0.1);
\end{tikzpicture} & \hypertarget{R-1F7.11.1.2.4}{R-1F7.11.1.2.4} & Funzionale

Desiderabile & Il docente deve poter inserire domande a risposta multipla & \hyperlink{UC5.1}{UC5.1}

\hyperlink{UC5.8}{UC5.8}\tabularnewline
\hline
\begin{tikzpicture}
\draw [->, thick] (0.8,0.2) -- (0.8,0.1) -- (1,0.1);
\end{tikzpicture} & \hypertarget{R-1F7.11.1.2.5}{R-1F7.11.1.2.5} & Funzionale

Desiderabile & Il docente deve poter inserire domande di tipo testo con parole omesse & \hyperlink{UC5.1}{UC5.1}

\hyperlink{UC5.9}{UC5.9}\tabularnewline
\hline
\begin{tikzpicture}
\draw [->, thick] (0.8,0.2) -- (0.8,0.1) -- (1,0.1);
\end{tikzpicture} & \hypertarget{R-1F7.11.1.2.6}{R-1F7.11.1.2.6} & Funzionale

Desiderabile & Il docente deve poter inserire domande con associazione di parole & \hyperlink{UC5.10}{UC5.10}

\hyperlink{UC5.1}{UC5.1}\tabularnewline
\hline
\begin{tikzpicture}
\draw [->, thick] (0.8,0.2) -- (0.8,0.1) -- (1,0.1);
\end{tikzpicture} & \hypertarget{R-2F7.11.1.2.7}{R-2F7.11.1.2.7} & Funzionale

Opzionale & Il docente deve poter inserire domande a risposta aperta & \hyperlink{UC5.11}{UC5.11}

\hyperlink{UC5.1}{UC5.1}\tabularnewline
\hline
\begin{tikzpicture}
\draw [->, thick] (0.4,0.2) -- (0.4,0.1) -- (1,0.1);
\end{tikzpicture} & \hypertarget{R-3F7.11.2}{R-3F7.11.2} & Funzionale

Obbligatorio & Il sistema deve consentire ad un docente di modificare una domanda & \hyperlink{UC5.2}{UC5.2}

\hyperlink{UC5}{UC5}\tabularnewline
\hline
\begin{tikzpicture}
\draw [->, thick] (0.4,0.2) -- (0.4,0.1) -- (1,0.1);
\end{tikzpicture} & \hypertarget{R-3F7.11.3}{R-3F7.11.3} & Funzionale

Obbligatorio & Il sistema deve consentire ad un docente di eliminare una propria domanda presente nel sistema & \hyperlink{UC5.3}{UC5.3}

\hyperlink{UC5}{UC5}\tabularnewline
\hline
\begin{tikzpicture}
\draw [->, thick] (0.2,0.2) -- (0.2,0.1) -- (1,0.1);
\end{tikzpicture} & \hypertarget{R-3F7.12}{R-3F7.12} & Funzionale

Obbligatorio & il sistema deve dare la possibilità al docente di modificare il proprio questionario & \hyperlink{UC6.2}{UC6.2}

\hyperlink{UC6}{UC6}\tabularnewline
\hline
\begin{tikzpicture}
\draw [->, thick] (0.4,0.2) -- (0.4,0.1) -- (1,0.1);
\end{tikzpicture} & \hypertarget{R-3F7.12.1}{R-3F7.12.1} & Funzionale

Obbligatorio & Il sistema deve dare la possibilità di ricercare un questionario & \hyperlink{UC16}{UC16}\tabularnewline
\hline
\begin{tikzpicture}
\draw [->, thick] (0.4,0.2) -- (0.4,0.1) -- (1,0.1);
\end{tikzpicture} & \hypertarget{R-3F7.12.2}{R-3F7.12.2} & Funzionale

Obbligatorio & Il sistema deve dare la possibile al docente di aggiungere al questionario già fatto una domanda. & \hyperlink{UC6.2.1}{UC6.2.1}

\hyperlink{UC6}{UC6}\tabularnewline
\hline
\begin{tikzpicture}
\draw [->, thick] (0.4,0.2) -- (0.4,0.1) -- (1,0.1);
\end{tikzpicture} & \hypertarget{R-3F7.12.3}{R-3F7.12.3} & Funzionale

Obbligatorio & Il sistema deve dare la possibilità al docente di togliere domande precedentemente aggiunte al questionario & \hyperlink{UC6.2.2}{UC6.2.2}

\hyperlink{UC6}{UC6}\tabularnewline
\hline
\begin{tikzpicture}
\draw [->, thick] (0.4,0.2) -- (0.4,0.1) -- (1,0.1);
\end{tikzpicture} & \hypertarget{R-3F7.12.4}{R-3F7.12.4} & Funzionale

Obbligatorio & Il sistema deve consentire ad un docente di modificare gli argomenti di un proprio questionario & \hyperlink{UC6.2.3}{UC6.2.3}

\hyperlink{UC6.2}{UC6.2}

\hyperlink{UC6}{UC6}\tabularnewline
\hline
\begin{tikzpicture}
\draw [->, thick] (0.2,0.2) -- (0.2,0.1) -- (1,0.1);
\end{tikzpicture} & \hypertarget{R-3F7.13}{R-3F7.13} & Funzionale

Obbligatorio & Il sistema deve dare la possibilità al docente di eliminare il proprio questionario & \hyperlink{UC6.3}{UC6.3}

\hyperlink{UC6}{UC6}\tabularnewline
\hline
 & \hypertarget{R-3F8}{R-3F8} & Funzionale

Obbligatorio & L'utente attraverso una specifica interfaccia potrà rispondere alle domande & Capitolato

\hyperlink{UC13.1}{UC13.1}\tabularnewline
\hline
 & \hypertarget{R-3F9}{R-3F9} & Funzionale

Obbligatorio & Il sistema deve avere almeno 3 tipi di utenti: amministratore, docente e studente & Capitolato\tabularnewline
\hline
 & \hypertarget{R-3F10}{R-3F10} & Funzionale

Obbligatorio & Il sistema deve gestire la registrazione di un utente & \hyperlink{UC2}{UC2}\tabularnewline
\hline
\begin{tikzpicture}
\draw [->, thick] (0.2,0.2) -- (0.2,0.1) -- (1,0.1);
\end{tikzpicture} & \hypertarget{R-3F10.1}{R-3F10.1} & Funzionale

Obbligatorio & Il sistema deve già possedere un proprietario all'installazione del sistema & Interno\tabularnewline
\hline
\begin{tikzpicture}
\draw [->, thick] (0.2,0.2) -- (0.2,0.1) -- (1,0.1);
\end{tikzpicture} & \hypertarget{R-3F10.2}{R-3F10.2} & Funzionale

Obbligatorio & La registrazione di un utente richiede l’inserimento di uno username, password e nome completo  & \hyperlink{UC2.1}{UC2.1}

\hyperlink{UC2.2}{UC2.2}

\hyperlink{UC2.3}{UC2.3}\tabularnewline
\hline
\begin{tikzpicture}
\draw [->, thick] (0.4,0.2) -- (0.4,0.1) -- (1,0.1);
\end{tikzpicture} & \hypertarget{R-3F10.2.1}{R-3F10.2.1} & Funzionale

Obbligatorio & Il sistema deve segnalare un errore nel caso non vengano inseriti tutti i campi necessari per la registrazione & \hyperlink{UC2.4}{UC2.4}

\hyperlink{UC2.5}{UC2.5}

\hyperlink{UC2.6}{UC2.6}

\hyperlink{UC2.7}{UC2.7}\tabularnewline
\hline
\begin{tikzpicture}
\draw [->, thick] (0.4,0.2) -- (0.4,0.1) -- (1,0.1);
\end{tikzpicture} & \hypertarget{R-3F10.2.2}{R-3F10.2.2} & Funzionale

Obbligatorio & La registrazione di uno utente deve richiedere l'inserimento delle informazioni personali & \hyperlink{UC2.1}{UC2.1}\tabularnewline
\hline
\begin{tikzpicture}
\draw [->, thick] (0.2,0.2) -- (0.2,0.1) -- (1,0.1);
\end{tikzpicture} & \hypertarget{R-3F10.3}{R-3F10.3} & Funzionale

Obbligatorio & Il sistema deve permettere la visualizzazione di un messaggio di errore se l'username dell' utente è già presente nel sistema & \hyperlink{UC2.7}{UC2.7}\tabularnewline
\hline
\begin{tikzpicture}
\draw [->, thick] (0.2,0.2) -- (0.2,0.1) -- (1,0.1);
\end{tikzpicture} & \hypertarget{R-3F10.4}{R-3F10.4} & Funzionale

Obbligatorio & Il sistema deve permettere la visualizzazione di un messaggio di errore se la password dell'utente è troppo corta & \hyperlink{UC2.6}{UC2.6}\tabularnewline
\hline
\begin{tikzpicture}
\draw [->, thick] (0.2,0.2) -- (0.2,0.1) -- (1,0.1);
\end{tikzpicture} & \hypertarget{R-3F10.5}{R-3F10.5} & Funzionale

Obbligatorio & Il sistema deve segnalare un errore se viene inserito un username già utilizzato & \hyperlink{UC2.7}{UC2.7}\tabularnewline
\hline
 & \hypertarget{R-3V11}{R-3V11} & Vincolo

Obbligatorio & Il database server side deve essere MongoDB & Capitolato

Interno\tabularnewline
\hline
 & \hypertarget{R-3F12}{R-3F12} & Funzionale

Obbligatorio & Il sistema deve gestire le azioni di amministratore & \hyperlink{UC30}{UC30}\tabularnewline
\hline
\begin{tikzpicture}
\draw [->, thick] (0.2,0.2) -- (0.2,0.1) -- (1,0.1);
\end{tikzpicture} & \hypertarget{R-3F12.1}{R-3F12.1} & Funzionale

Obbligatorio & L'amministratore deve poter cambiare il ruolo di un utente con grado inferiore al proprio in un ruolo con grado inferiore al proprio & \hyperlink{UC30.1}{UC30.1}

\hyperlink{UC30}{UC30}\tabularnewline
\hline
\begin{tikzpicture}
\draw [->, thick] (0.2,0.2) -- (0.2,0.1) -- (1,0.1);
\end{tikzpicture} & \hypertarget{R-3F12.2}{R-3F12.2} & Funzionale

Obbligatorio & Il sistema deve dare la possibilità ad un amministratore o proprietario di eliminare un utente di ruolo inferiore al proprio & \hyperlink{UC30.2}{UC30.2}

\hyperlink{UC30}{UC30}\tabularnewline
\hline
 & \hypertarget{R-2F13}{R-2F13} & Funzionale

Opzionale & Il sistema deve prevedere la gestione di classi da parte di un docente & \hyperlink{UC7}{UC7}\tabularnewline
\hline
\begin{tikzpicture}
\draw [->, thick] (0.2,0.2) -- (0.2,0.1) -- (1,0.1);
\end{tikzpicture} & \hypertarget{R-2F13.1}{R-2F13.1} & Funzionale

Opzionale & Un docente deve avere la possibilità di creare una nuova classe & \hyperlink{UC7.1}{UC7.1}

\hyperlink{UC7}{UC7}\tabularnewline
\hline
\begin{tikzpicture}
\draw [->, thick] (0.4,0.2) -- (0.4,0.1) -- (1,0.1);
\end{tikzpicture} & \hypertarget{R-2F13.1.1}{R-2F13.1.1} & Funzionale

Opzionale & Un docente deve poter specificare il nome della nuova classe & \hyperlink{UC7.1.1}{UC7.1.1}

\hyperlink{UC7.1}{UC7.1}

\hyperlink{UC7}{UC7}\tabularnewline
\hline
\begin{tikzpicture}
\draw [->, thick] (0.6,0.2) -- (0.6,0.1) -- (1,0.1);
\end{tikzpicture} & \hypertarget{R-2F13.1.1.1}{R-2F13.1.1.1} & Funzionale

Opzionale & Nel caso il nome della classe sia già presente il sistema deve segnalare un errore & \hyperlink{UC7.1.1}{UC7.1.1}

\hyperlink{UC7.1}{UC7.1}

\hyperlink{UC7.4}{UC7.4}

\hyperlink{UC7}{UC7}\tabularnewline
\hline
\begin{tikzpicture}
\draw [->, thick] (0.4,0.2) -- (0.4,0.1) -- (1,0.1);
\end{tikzpicture} & \hypertarget{R-2F13.1.2}{R-2F13.1.2} & Funzionale

Opzionale & Un docente deve poter specificare gli argomenti di una nuova classe & \hyperlink{UC7.1.2}{UC7.1.2}

\hyperlink{UC7.1}{UC7.1}

\hyperlink{UC7}{UC7}\tabularnewline
\hline
\begin{tikzpicture}
\draw [->, thick] (0.4,0.2) -- (0.4,0.1) -- (1,0.1);
\end{tikzpicture} & \hypertarget{R-2F13.1.3}{R-2F13.1.3} & Funzionale

Opzionale & Un docente deve poter specificare la password relativa ad una nuova classe & \hyperlink{UC7.1.3}{UC7.1.3}

\hyperlink{UC7.1}{UC7.1}

\hyperlink{UC7}{UC7}\tabularnewline
\hline
\begin{tikzpicture}
\draw [->, thick] (0.2,0.2) -- (0.2,0.1) -- (1,0.1);
\end{tikzpicture} & \hypertarget{R-2F13.2}{R-2F13.2} & Funzionale

Opzionale & Un docente deve avere la possibilità di eliminare una classe esistente & \hyperlink{UC7.3}{UC7.3}

\hyperlink{UC7}{UC7}\tabularnewline
\hline
\begin{tikzpicture}
\draw [->, thick] (0.2,0.2) -- (0.2,0.1) -- (1,0.1);
\end{tikzpicture} & \hypertarget{R-2F13.3}{R-2F13.3} & Funzionale

Opzionale & Un docente deve poter modificare una classe & \hyperlink{UC7.2}{UC7.2}

\hyperlink{UC7}{UC7}\tabularnewline
\hline
\begin{tikzpicture}
\draw [->, thick] (0.4,0.2) -- (0.4,0.1) -- (1,0.1);
\end{tikzpicture} & \hypertarget{R-2F13.3.1}{R-2F13.3.1} & Funzionale

Opzionale & Un docente deve poter modificare il nome di una propria classe esistente & \hyperlink{UC7.1.1}{UC7.1.1}

\hyperlink{UC7.2}{UC7.2}

\hyperlink{UC7}{UC7}\tabularnewline
\hline
\begin{tikzpicture}
\draw [->, thick] (0.4,0.2) -- (0.4,0.1) -- (1,0.1);
\end{tikzpicture} & \hypertarget{R-2F13.3.2}{R-2F13.3.2} & Funzionale

Opzionale & Un docente deve poter modificare gli argomenti di una sua classe esistente & \hyperlink{UC7.1.2}{UC7.1.2}

\hyperlink{UC7.2}{UC7.2}

\hyperlink{UC7}{UC7}\tabularnewline
\hline
\begin{tikzpicture}
\draw [->, thick] (0.4,0.2) -- (0.4,0.1) -- (1,0.1);
\end{tikzpicture} & \hypertarget{R-2F13.3.3}{R-2F13.3.3} & Funzionale

Opzionale & Un docente deve poter modificare la password di una delle proprie classi & \hyperlink{UC7.1.3}{UC7.1.3}

\hyperlink{UC7.2}{UC7.2}

\hyperlink{UC7}{UC7}\tabularnewline
\hline
 & \hypertarget{R-3F14}{R-3F14} & Funzionale

Obbligatorio & Il sistema deve consentire ad un utente di gestire i propri dati memorizzati nel sistema & \hyperlink{UC3}{UC3}\tabularnewline
\hline
\begin{tikzpicture}
\draw [->, thick] (0.2,0.2) -- (0.2,0.1) -- (1,0.1);
\end{tikzpicture} & \hypertarget{R-3F14.1}{R-3F14.1} & Funzionale

Obbligatorio & Il sistema deve consentire ad un utente di modificare la propria username & \hyperlink{UC3.2}{UC3.2}

\hyperlink{UC3}{UC3}\tabularnewline
\hline
\begin{tikzpicture}
\draw [->, thick] (0.4,0.2) -- (0.4,0.1) -- (1,0.1);
\end{tikzpicture} & \hypertarget{R-3F14.1.1}{R-3F14.1.1} & Funzionale

Obbligatorio & Il sistema deve gestire eventuali errori di inserimento username & \hyperlink{UC3.2}{UC3.2}

\hyperlink{UC3.4}{UC3.4}

\hyperlink{UC3.6}{UC3.6}

\hyperlink{UC3}{UC3}\tabularnewline
\hline
\begin{tikzpicture}
\draw [->, thick] (0.4,0.2) -- (0.4,0.1) -- (1,0.1);
\end{tikzpicture} & \hypertarget{R-3F14.1.2}{R-3F14.1.2} & Funzionale

Obbligatorio & Il sistema deve segnalare un errore in caso di inserimento di username già presente nel sistema & \hyperlink{UC3.6}{UC3.6}

\hyperlink{UC3}{UC3}\tabularnewline
\hline
\begin{tikzpicture}
\draw [->, thick] (0.2,0.2) -- (0.2,0.1) -- (1,0.1);
\end{tikzpicture} & \hypertarget{R-3F14.2}{R-3F14.2} & Funzionale

Obbligatorio & Il sistema deve consentire ad un utente di modificare la propria password di accesso al sistema & \hyperlink{UC3.3}{UC3.3}

\hyperlink{UC3}{UC3}\tabularnewline
\hline
\begin{tikzpicture}
\draw [->, thick] (0.4,0.2) -- (0.4,0.1) -- (1,0.1);
\end{tikzpicture} & \hypertarget{R-3F14.2.1}{R-3F14.2.1} & Funzionale

Obbligatorio & Il sistema deve richiedere l'inserimento della vecchia password all'utente per consentire la modifica della password & \hyperlink{UC3.3.1}{UC3.3.1}

\hyperlink{UC3.3}{UC3.3}

\hyperlink{UC3}{UC3}\tabularnewline
\hline
\begin{tikzpicture}
\draw [->, thick] (0.6,0.2) -- (0.6,0.1) -- (1,0.1);
\end{tikzpicture} & \hypertarget{R-3F14.2.1.1}{R-3F14.2.1.1} & Funzionale

Obbligatorio & Il sistema deve segnalare l'inserimento di una password scorretta & \hyperlink{UC3.3.2}{UC3.3.2}

\hyperlink{UC3.7}{UC3.7}

\hyperlink{UC3}{UC3}\tabularnewline
\hline
\begin{tikzpicture}
\draw [->, thick] (0.4,0.2) -- (0.4,0.1) -- (1,0.1);
\end{tikzpicture} & \hypertarget{R-3F14.2.2}{R-3F14.2.2} & Funzionale

Obbligatorio & Il sistema deve richiedere l'inserimento della nuova password & \hyperlink{UC3.3}{UC3.3}

\hyperlink{UC3}{UC3}\tabularnewline
\hline
\begin{tikzpicture}
\draw [->, thick] (0.6,0.2) -- (0.6,0.1) -- (1,0.1);
\end{tikzpicture} & \hypertarget{R-3F14.2.2.1}{R-3F14.2.2.1} & Funzionale

Obbligatorio & Il sistema deve verificare che la nuova password sia di almeno otto caratteri e segnalarlo in caso contrario & \hyperlink{UC3.3}{UC3.3}

\hyperlink{UC3.7}{UC3.7}

\hyperlink{UC3}{UC3}\tabularnewline
\hline
\begin{tikzpicture}
\draw [->, thick] (0.2,0.2) -- (0.2,0.1) -- (1,0.1);
\end{tikzpicture} & \hypertarget{R-3F14.3}{R-3F14.3} & Funzionale

Obbligatorio & Il sistema deve consentire all'utente di modificare i propri dati personali & \hyperlink{UC3.1}{UC3.1}

\hyperlink{UC3.2}{UC3.2}

\hyperlink{UC3}{UC3}\tabularnewline
\hline
\begin{tikzpicture}
\draw [->, thick] (0.4,0.2) -- (0.4,0.1) -- (1,0.1);
\end{tikzpicture} & \hypertarget{R-3F14.3.1}{R-3F14.3.1} & Funzionale

Obbligatorio & Il sistema deve consentire all'utente la modifica del proprio nome  & 

\hyperlink{UC3.1}{UC3.1}

\hyperlink{UC3}{UC3}\tabularnewline
\hline
\begin{tikzpicture}
\draw [->, thick] (0.4,0.2) -- (0.4,0.1) -- (1,0.1);
\end{tikzpicture} & \hypertarget{R-3F14.3.2}{R-3F14.3.2} & Funzionale

Obbligatorio & Il sistema deve consentire all'utente di modificare il proprio username & \hyperlink{UC3.2}{UC3.2}

\hyperlink{UC3}{UC3}\tabularnewline
\hline
\begin{tikzpicture}
\draw [->, thick] (0.4,0.2) -- (0.4,0.1) -- (1,0.1);
\end{tikzpicture} & \hypertarget{R-3F14.3.3}{R-3F14.3.3} & Funzionale

Obbligatorio & Il sistema deve controllare che il nome completo contenga almeno 2 caratteri e in caso contrario mostrare un errore & \hyperlink{UC3.1}{UC3.1}

\hyperlink{UC3.5}{UC3.5}

\hyperlink{UC3}{UC3}\tabularnewline
\hline
\begin{tikzpicture}
\draw [->, thick] (0.4,0.2) -- (0.4,0.1) -- (1,0.1);
\end{tikzpicture} & \hypertarget{R-3F14.3.4}{R-3F14.3.4} & Funzionale

Obbligatorio & Il sistema deve controllare che l'username contenga almeno 6 caratteri e in caso contrario mostrare un errore & \hyperlink{UC3.2}{UC3.2}

\hyperlink{UC3.4}{UC3.4}

\hyperlink{UC3}{UC3}\tabularnewline
\hline
 & \hypertarget{R-3F15}{R-3F15} & Funzionale

Obbligatorio & Il sistema deve gestire la ricerca di un questionario & \hyperlink{UC16}{UC16}\tabularnewline
\hline
\begin{tikzpicture}
\draw [->, thick] (0.2,0.2) -- (0.2,0.1) -- (1,0.1);
\end{tikzpicture} & \hypertarget{R-3F15.1}{R-3F15.1} & Funzionale

Obbligatorio & Il sistema deve gestire la ricerca di un questionario per titolo & \hyperlink{UC16}{UC16}

\hyperlink{UC17}{UC17}\tabularnewline
\hline
\begin{tikzpicture}
\draw [->, thick] (0.2,0.2) -- (0.2,0.1) -- (1,0.1);
\end{tikzpicture} & \hypertarget{R-2F15.2}{R-2F15.2} & Funzionale

Opzionale & Il sistema deve gestire la ricerca di un questionario per classe & \hyperlink{UC16}{UC16}

\hyperlink{UC18}{UC18}\tabularnewline
\hline
\begin{tikzpicture}
\draw [->, thick] (0.2,0.2) -- (0.2,0.1) -- (1,0.1);
\end{tikzpicture} & \hypertarget{R-3F15.3}{R-3F15.3} & Funzionale

Obbligatorio & Il sistema deve gestire la ricerca di un questionario per docente & \hyperlink{UC16}{UC16}

\hyperlink{UC20}{UC20}\tabularnewline
\hline
\begin{tikzpicture}
\draw [->, thick] (0.2,0.2) -- (0.2,0.1) -- (1,0.1);
\end{tikzpicture} & \hypertarget{R-3F15.4}{R-3F15.4} & Funzionale

Obbligatorio & Il sistema deve gestire la ricerca di un questionario per argomento & \hyperlink{UC16}{UC16}

\hyperlink{UC19}{UC19}\tabularnewline
\hline
\begin{tikzpicture}
\draw [->, thick] (0.2,0.2) -- (0.2,0.1) -- (1,0.1);
\end{tikzpicture} & \hypertarget{R-2F15.5}{R-2F15.5} & Funzionale

Opzionale & Il sistema deve gestire la ricerca di un questionario per difficoltà & \hyperlink{UC16}{UC16}

\hyperlink{UC21}{UC21}\tabularnewline
\hline
 & \hypertarget{R-2F16}{R-2F16} & Funzionale

Opzionale & Il sistema deve consentire l'iscrizione ad una classe da parte di uno studente & \hyperlink{UC14}{UC14}\tabularnewline
\hline
\begin{tikzpicture}
\draw [->, thick] (0.2,0.2) -- (0.2,0.1) -- (1,0.1);
\end{tikzpicture} & \hypertarget{R-2F16.1}{R-2F16.1} & Funzionale

Opzionale & Il sistema deve chiedere l'inserimento della password prima di confermare l'iscrizione ad una classe & \hyperlink{UC14.1}{UC14.1}

\hyperlink{UC14}{UC14}\tabularnewline
\hline
\begin{tikzpicture}
\draw [->, thick] (0.4,0.2) -- (0.4,0.1) -- (1,0.1);
\end{tikzpicture} & \hypertarget{R-2F16.1.1}{R-2F16.1.1} & Funzionale

Opzionale & Il sistema deve verificare la password prima di confermare l'iscrizione dello studente alla classe. & \hyperlink{UC14.2}{UC14.2}

\hyperlink{UC14}{UC14}\tabularnewline
\hline
\begin{tikzpicture}
\draw [->, thick] (0.2,0.2) -- (0.2,0.1) -- (1,0.1);
\end{tikzpicture} & \hypertarget{R-2F16.2}{R-2F16.2} & Funzionale

Opzionale & Il sistema deve consentire di ricercare la classe a cui iscriversi & \hyperlink{UC27}{UC27}

\hyperlink{UC28}{UC28}

\hyperlink{UC29}{UC29}\tabularnewline
\hline
 & \hypertarget{R-3F17}{R-3F17} & Funzionale

Obbligatorio & Il sistema deve dare la possibilità di eseguire un questionario ad uno studente & \hyperlink{UC13}{UC13}\tabularnewline
\hline
\begin{tikzpicture}
\draw [->, thick] (0.2,0.2) -- (0.2,0.1) -- (1,0.1);
\end{tikzpicture} & \hypertarget{R-3F17.1}{R-3F17.1} & Funzionale

Obbligatorio & Il sistema deve dare la possibilità ad uno studente di rispondere ad una domanda & \hyperlink{UC13.1}{UC13.1}

\hyperlink{UC13}{UC13}\tabularnewline
\hline
\begin{tikzpicture}
\draw [->, thick] (0.4,0.2) -- (0.4,0.1) -- (1,0.1);
\end{tikzpicture} & \hypertarget{R-3F17.1.1}{R-3F17.1.1} & Funzionale

Obbligatorio & Il sistema deve dare la possibilità ad uno studente di rispondere ad una domanda di tipo vero/falso & \hyperlink{UC13.1}{UC13.1}

\hyperlink{UC13.4}{UC13.4}\tabularnewline
\hline
\begin{tikzpicture}
\draw [->, thick] (0.4,0.2) -- (0.4,0.1) -- (1,0.1);
\end{tikzpicture} & \hypertarget{R-3F17.1.2}{R-3F17.1.2} & Funzionale

Obbligatorio & Il sistema deve dare la possibilità ad uno studente di rispondere ad una domanda a scelta multipla & \hyperlink{UC13.1}{UC13.1}

\hyperlink{UC13.5}{UC13.5}\tabularnewline
\hline
\begin{tikzpicture}
\draw [->, thick] (0.4,0.2) -- (0.4,0.1) -- (1,0.1);
\end{tikzpicture} & \hypertarget{R-1F17.1.3}{R-1F17.1.3} & Funzionale

Desiderabile & Il sistema deve dare la possibilità ad uno studente di rispondere ad una domanda a risposta multipla ovvero una domanda a scelta multipla  con più risposte giuste & \hyperlink{UC13.1}{UC13.1}

\hyperlink{UC13.6}{UC13.6}\tabularnewline
\hline
\begin{tikzpicture}
\draw [->, thick] (0.4,0.2) -- (0.4,0.1) -- (1,0.1);
\end{tikzpicture} & \hypertarget{R-1F17.1.4}{R-1F17.1.4} & Funzionale

Desiderabile & Il sistema deve dare la possibilità ad uno studente di rispondere ad una domanda di tipo testo con parole omesse & \hyperlink{UC13.1}{UC13.1}

\hyperlink{UC13.7}{UC13.7}\tabularnewline
\hline
\begin{tikzpicture}
\draw [->, thick] (0.4,0.2) -- (0.4,0.1) -- (1,0.1);
\end{tikzpicture} & \hypertarget{R-1F17.1.5}{R-1F17.1.5} & Funzionale

Desiderabile & Il sistema deve dare la possibilità ad uno studente di rispondere ad una domanda con associazione di parole & \hyperlink{UC13.1}{UC13.1}

\hyperlink{UC13.8}{UC13.8}\tabularnewline
\hline
\begin{tikzpicture}
\draw [->, thick] (0.4,0.2) -- (0.4,0.1) -- (1,0.1);
\end{tikzpicture} & \hypertarget{R-2F17.1.6}{R-2F17.1.6} & Funzionale

Opzionale & Il sistema deve dare la possibilità ad uno studente di rispondere ad una domanda a risposta aperta & \hyperlink{UC13.1}{UC13.1}

\hyperlink{UC13.9}{UC13.9}\tabularnewline
\hline
\begin{tikzpicture}
\draw [->, thick] (0.2,0.2) -- (0.2,0.1) -- (1,0.1);
\end{tikzpicture} & \hypertarget{R-3F17.2}{R-3F17.2} & Funzionale

Obbligatorio & Il sistema deve richiedere allo studente la conferma alla fine della compilazione di un questionario, perché questo sia considerato eseguito & \hyperlink{UC13.3}{UC13.3}

\hyperlink{UC13}{UC13}\tabularnewline
\hline
\begin{tikzpicture}
\draw [->, thick] (0.4,0.2) -- (0.4,0.1) -- (1,0.1);
\end{tikzpicture} & \hypertarget{R-3F17.2.1}{R-3F17.2.1} & Funzionale

Obbligatorio & Il sistema deve segnalare un errore nel caso una domanda del questionario non sia stata compilata & \hyperlink{UC13.2}{UC13.2}

\hyperlink{UC13.3}{UC13.3}

\hyperlink{UC13}{UC13}\tabularnewline
\hline
\begin{tikzpicture}
\draw [->, thick] (0.2,0.2) -- (0.2,0.1) -- (1,0.1);
\end{tikzpicture} & \hypertarget{R-3F17.3}{R-3F17.3} & Funzionale

Obbligatorio & Il sistema deve dare la possibilità ad uno studente di andare alla domanda successiva durante l'esecuzione del questionario & \hyperlink{UC13}{UC13}\tabularnewline
\hline
\begin{tikzpicture}
\draw [->, thick] (0.2,0.2) -- (0.2,0.1) -- (1,0.1);
\end{tikzpicture} & \hypertarget{R-3F17.4}{R-3F17.4} & Funzionale

Obbligatorio & Il sistema deve dare la possibilità ad uno studente di andare alla domanda precedente durante l'esecuzione del questionario & \hyperlink{UC13}{UC13}\tabularnewline
\hline
 & \hypertarget{R-3F18}{R-3F18} & Funzionale

Obbligatorio & Il sistema deve consentire all'utente autenticato di effettuare il logout & \hyperlink{UC4}{UC4}\tabularnewline
\hline
 & \hypertarget{R-3F19}{R-3F19} & Funzionale

Obbligatorio & Il sistema deve gestire le operazioni che può fare il proprietario del sistema & \hyperlink{UC30}{UC30}\tabularnewline
\hline
\begin{tikzpicture}
\draw [->, thick] (0.2,0.2) -- (0.2,0.1) -- (1,0.1);
\end{tikzpicture} & \hypertarget{R-3F19.1}{R-3F19.1} & Funzionale

Obbligatorio & Il sistema deve assicurare l'esistenza di un unico utente proprietario & Interno\tabularnewline
\hline
\begin{tikzpicture}
\draw [->, thick] (0.2,0.2) -- (0.2,0.1) -- (1,0.1);
\end{tikzpicture} & \hypertarget{R-3F19.2}{R-3F19.2} & Funzionale

Obbligatorio & Il sistema deve consentire solo all'utente proprietario di promuovere un utente al ruolo di amministratore & \hyperlink{UC30.1}{UC30.1}

\hyperlink{UC30}{UC30}\tabularnewline
\hline
\begin{tikzpicture}
\draw [->, thick] (0.2,0.2) -- (0.2,0.1) -- (1,0.1);
\end{tikzpicture} & \hypertarget{R-3F19.3}{R-3F19.3} & Funzionale

Obbligatorio & Il sistema deve consentire solo all'utente proprietario di eliminare un amministratore & \hyperlink{UC30.2}{UC30.2}

\hyperlink{UC30}{UC30}\tabularnewline
\hline
 & \hypertarget{R-3F20}{R-3F20} & Funzionale

Obbligatorio & Il sistema deve consentire ad un docente, amministratore o proprietario di ricercare una domanda presente nel sistema & \hyperlink{UC22}{UC22}\tabularnewline
\hline
\begin{tikzpicture}
\draw [->, thick] (0.2,0.2) -- (0.2,0.1) -- (1,0.1);
\end{tikzpicture} & \hypertarget{R-3F20.1}{R-3F20.1} & Funzionale

Obbligatorio & Il sistema deve consentire a un docente, amministratore o proprietario di ricercare una domanda per argomento & \hyperlink{UC22}{UC22}

\hyperlink{UC24}{UC24}\tabularnewline
\hline
\begin{tikzpicture}
\draw [->, thick] (0.2,0.2) -- (0.2,0.1) -- (1,0.1);
\end{tikzpicture} & \hypertarget{R-3F20.2}{R-3F20.2} & Funzionale

Obbligatorio & Il sistema deve consentire a un docente, amministratore o proprietario di ricercare una domanda per docente & \hyperlink{UC22}{UC22}

\hyperlink{UC26}{UC26}\tabularnewline
\hline
\begin{tikzpicture}
\draw [->, thick] (0.2,0.2) -- (0.2,0.1) -- (1,0.1);
\end{tikzpicture} & \hypertarget{R-2F20.3}{R-2F20.3} & Funzionale

Opzionale & Il sistema deve consentire a un docente, amministratore o proprietario di ricercare una domanda per difficoltà & \hyperlink{UC22}{UC22}

\hyperlink{UC25}{UC25}\tabularnewline
\hline
\begin{tikzpicture}
\draw [->, thick] (0.2,0.2) -- (0.2,0.1) -- (1,0.1);
\end{tikzpicture} & \hypertarget{R-3F20.4}{R-3F20.4} & Funzionale

Obbligatorio & Il sistema deve consentire a un docente, amministratore o proprietario di ricercare una domanda per keywords & \hyperlink{UC22}{UC22}

\hyperlink{UC23}{UC23}\tabularnewline
\hline
 & \hypertarget{R-2F21}{R-2F21} & Funzionale

Opzionale & Il sistema deve fornire la possibilità ad uno studente  di cercare una classe & \hyperlink{UC27}{UC27}\tabularnewline
\hline
\begin{tikzpicture}
\draw [->, thick] (0.2,0.2) -- (0.2,0.1) -- (1,0.1);
\end{tikzpicture} & \hypertarget{R-2F21.1}{R-2F21.1} & Funzionale

Opzionale & Il sistema deve fornire la possibilità ad uno studente  di cercare una classe per argomento & \hyperlink{UC27}{UC27}

\hyperlink{UC29}{UC29}\tabularnewline
\hline
\begin{tikzpicture}
\draw [->, thick] (0.2,0.2) -- (0.2,0.1) -- (1,0.1);
\end{tikzpicture} & \hypertarget{R-2F21.2}{R-2F21.2} & Funzionale

Opzionale & Il sistema deve fornire la possibilità ad uno studente  di cercare una classe per docente & \hyperlink{UC27}{UC27}

\hyperlink{UC28}{UC28}\tabularnewline
\hline
 & \hypertarget{R-2F22}{R-2F22} & Funzionale

Opzionale & Il sistema deve dare la possibilità di recupero password ad un utente non autenticato inviando una mail con la nuova password
 & \hyperlink{UC35}{UC35}\tabularnewline
\hline
\begin{tikzpicture}
\draw [->, thick] (0.2,0.2) -- (0.2,0.1) -- (1,0.1);
\end{tikzpicture} & \hypertarget{R-2F22.1}{R-2F22.1} & Funzionale

Opzionale & Il sistema deve dare la possibilità di inserire  mail per recupero password ad un utente non autenticato & \hyperlink{UC35}{UC35}\tabularnewline
\hline
\begin{tikzpicture}
\draw [->, thick] (0.2,0.2) -- (0.2,0.1) -- (1,0.1);
\end{tikzpicture} & \hypertarget{R-2F22.2}{R-2F22.2} & Funzionale

Opzionale & Il sistema deve segnalare ad un utente non autenticato che l'email inserito per recupero password non è presente nel sistema & \hyperlink{UC35.1}{UC35.1}

\hyperlink{UC35}{UC35}\tabularnewline
\hline
 & \hypertarget{R-3F23}{R-3F23} & Funzionale

Obbligatorio & Il sistema deve permettere ad un docente di gestire gli argomenti & \hyperlink{UC8}{UC8}\tabularnewline
\hline
\begin{tikzpicture}
\draw [->, thick] (0.2,0.2) -- (0.2,0.1) -- (1,0.1);
\end{tikzpicture} & \hypertarget{R-3F23.1}{R-3F23.1} & Funzionale

Obbligatorio & Il sistema deve consentire ad un docente di creare un nuovo argomento & \hyperlink{UC8.2}{UC8.2}

\hyperlink{UC8}{UC8}\tabularnewline
\hline
\begin{tikzpicture}
\draw [->, thick] (0.4,0.2) -- (0.4,0.1) -- (1,0.1);
\end{tikzpicture} & \hypertarget{R-3F23.1.1}{R-3F23.1.1} & Funzionale

Obbligatorio & Il sistema deve verificare se l'argomento sia già presente nel sistema e notificarlo in caso contrario & \hyperlink{UC8.2}{UC8.2}

\hyperlink{UC8.3}{UC8.3}

\hyperlink{UC8}{UC8}\tabularnewline
\hline
\begin{tikzpicture}
\draw [->, thick] (0.2,0.2) -- (0.2,0.1) -- (1,0.1);
\end{tikzpicture} & \hypertarget{R-3F23.2}{R-3F23.2} & Funzionale

Obbligatorio & Il sistema deve permettere ad un docente di eliminare un argomento & \hyperlink{UC8.5}{UC8.5}

\hyperlink{UC8}{UC8}\tabularnewline
\hline
\begin{tikzpicture}
\draw [->, thick] (0.4,0.2) -- (0.4,0.1) -- (1,0.1);
\end{tikzpicture} & \hypertarget{R-3F23.2.1}{R-3F23.2.1} & Funzionale

Obbligatorio & Il sistema deve permettere al docente di selezionare un argomento presente da eliminare & \hyperlink{UC8.1}{UC8.1}

\hyperlink{UC8.5}{UC8.5}

\hyperlink{UC8}{UC8}\tabularnewline
\hline
\begin{tikzpicture}
\draw [->, thick] (0.4,0.2) -- (0.4,0.1) -- (1,0.1);
\end{tikzpicture} & \hypertarget{R-3F23.2.2}{R-3F23.2.2} & Funzionale

Obbligatorio & Il sistema deve impedire al docente di eliminare un argomento corrisponde a domande o questionari presenti nel sistema & \hyperlink{UC8.5}{UC8.5}

\hyperlink{UC8.6}{UC8.6}

\hyperlink{UC8}{UC8}\tabularnewline
\hline
\begin{tikzpicture}
\draw [->, thick] (0.2,0.2) -- (0.2,0.1) -- (1,0.1);
\end{tikzpicture} & \hypertarget{R-3F23.3}{R-3F23.3} & Funzionale

Obbligatorio & Il sistema deve permettere ad un docente di modificare un argomento & \hyperlink{UC8.4}{UC8.4}

\hyperlink{UC8}{UC8}\tabularnewline
\hline
\begin{tikzpicture}
\draw [->, thick] (0.2,0.2) -- (0.2,0.1) -- (1,0.1);
\end{tikzpicture} & \hypertarget{R-3F23.4}{R-3F23.4} & Funzionale

Obbligatorio & Il sistema deve permettere ad un docente di esplorare gli argomenti presenti & \hyperlink{UC8.1}{UC8.1}

\hyperlink{UC8}{UC8}\tabularnewline
\hline
 & \hypertarget{R-2F24}{R-2F24} & Funzionale

Opzionale & Il sistema deve consentire ad un docente di visualizzare statistiche & \hyperlink{UC9}{UC9}\tabularnewline
\hline
\begin{tikzpicture}
\draw [->, thick] (0.2,0.2) -- (0.2,0.1) -- (1,0.1);
\end{tikzpicture} & \hypertarget{R-2F24.1}{R-2F24.1} & Funzionale

Opzionale & Il sistema deve permettere ad un docente di visualizzare le statistiche relative ad una domanda & \hyperlink{UC10}{UC10}

\hyperlink{UC9}{UC9}\tabularnewline
\hline
\begin{tikzpicture}
\draw [->, thick] (0.4,0.2) -- (0.4,0.1) -- (1,0.1);
\end{tikzpicture} & \hypertarget{R-2F24.1.1}{R-2F24.1.1} & Funzionale

Opzionale & Il sistema deve permettere ad un docente di ricercare la domanda di cui visualizzare le statistiche  & \hyperlink{UC10}{UC10}

\hyperlink{UC22}{UC22}

\hyperlink{UC9}{UC9}\tabularnewline
\hline
\begin{tikzpicture}
\draw [->, thick] (0.2,0.2) -- (0.2,0.1) -- (1,0.1);
\end{tikzpicture} & \hypertarget{R-2F24.2}{R-2F24.2} & Funzionale

Opzionale & Il sistema deve permettere ad un docente di visualizzare le statistiche relative ad un questionario & \hyperlink{UC11}{UC11}

\hyperlink{UC9}{UC9}\tabularnewline
\hline
\begin{tikzpicture}
\draw [->, thick] (0.4,0.2) -- (0.4,0.1) -- (1,0.1);
\end{tikzpicture} & \hypertarget{R-2F24.2.1}{R-2F24.2.1} & Funzionale

Opzionale & Il sistema deve consentire ad un docente di ricercare il questionario di cui visualizzare le statistiche & \hyperlink{UC11}{UC11}

\hyperlink{UC16}{UC16}

\hyperlink{UC9}{UC9}\tabularnewline
\hline
\begin{tikzpicture}
\draw [->, thick] (0.2,0.2) -- (0.2,0.1) -- (1,0.1);
\end{tikzpicture} & \hypertarget{R-2F24.3}{R-2F24.3} & Funzionale

Opzionale & Il sistema deve permettere ad un docente di visualizzare le statistiche relative ad una delle proprie classi & \hyperlink{UC12}{UC12}

\hyperlink{UC9}{UC9}\tabularnewline
\hline
\begin{tikzpicture}
\draw [->, thick] (0.4,0.2) -- (0.4,0.1) -- (1,0.1);
\end{tikzpicture} & \hypertarget{R-2F24.3.1}{R-2F24.3.1} & Funzionale

Opzionale & Un docente deve poter visualizzare i risultati dei suoi questionari in base alla classe & \hyperlink{UC12.2}{UC12.2}

\hyperlink{UC12}{UC12}

\hyperlink{UC9}{UC9}\tabularnewline
\hline
\begin{tikzpicture}
\draw [->, thick] (0.4,0.2) -- (0.4,0.1) -- (1,0.1);
\end{tikzpicture} & \hypertarget{R-2F24.3.2}{R-2F24.3.2} & Funzionale

Opzionale & Un docente deve poter selezionare tra le proprie classi quella di cui visualizzare le statistiche & \hyperlink{UC12}{UC12}

\hyperlink{UC9}{UC9}\tabularnewline
\hline
\begin{tikzpicture}
\draw [->, thick] (0.4,0.2) -- (0.4,0.1) -- (1,0.1);
\end{tikzpicture} & \hypertarget{R-2F24.3.3}{R-2F24.3.3} & Funzionale

Opzionale & Un docente deve poter visualizzare un sommario delle statistiche della classe & \hyperlink{UC12.3}{UC12.3}

\hyperlink{UC12}{UC12}

\hyperlink{UC9}{UC9}\tabularnewline
\hline
\begin{tikzpicture}
\draw [->, thick] (0.4,0.2) -- (0.4,0.1) -- (1,0.1);
\end{tikzpicture} & \hypertarget{R-2F24.3.4}{R-2F24.3.4} & Funzionale

Opzionale & Un docente deve poter visualizzare le statistiche relative ad un particolare studente della classe & \hyperlink{UC12.4}{UC12.4}

\hyperlink{UC12}{UC12}

\hyperlink{UC9}{UC9}\tabularnewline
\hline
\begin{tikzpicture}
\draw [->, thick] (0.6,0.2) -- (0.6,0.1) -- (1,0.1);
\end{tikzpicture} & \hypertarget{R-2F24.3.4.1}{R-2F24.3.4.1} & Funzionale

Opzionale & Un docente deve poter cercare e selezionare un particolare questionario tra quelli eseguiti dallo studente di cui visualizzare le statistiche & \hyperlink{UC12.4.1}{UC12.4.1}

\hyperlink{UC12.4}{UC12.4}

\hyperlink{UC12}{UC12}

\hyperlink{UC16}{UC16}

\hyperlink{UC9}{UC9}\tabularnewline
\hline
\begin{tikzpicture}
\draw [->, thick] (0.4,0.2) -- (0.4,0.1) -- (1,0.1);
\end{tikzpicture} & \hypertarget{R-2F24.3.5}{R-2F24.3.5} & Funzionale

Opzionale & Un docente deve poter visualizzare i risultati delle sue domande in base alla classe & \hyperlink{UC12.1}{UC12.1}

\hyperlink{UC12}{UC12}

\hyperlink{UC9}{UC9}\tabularnewline
\hline
\begin{tikzpicture}
\draw [->, thick] (0.6,0.2) -- (0.6,0.1) -- (1,0.1);
\end{tikzpicture} & \hypertarget{R-2F24.3.5.1}{R-2F24.3.5.1} & Funzionale

Opzionale & Il sistema deve consentire ad un docente di ricercare una domanda tra quelle della classe & \hyperlink{UC12.1}{UC12.1}

\hyperlink{UC12}{UC12}

\hyperlink{UC22}{UC22}

\hyperlink{UC9}{UC9}\tabularnewline
\hline
 & \hypertarget{R-2F25}{R-2F25} & Funzionale

Opzionale & Il sistema deve consentire ad uno studente di visualizzare le proprie statistiche & \hyperlink{UC15}{UC15}\tabularnewline
\hline
\begin{tikzpicture}
\draw [->, thick] (0.2,0.2) -- (0.2,0.1) -- (1,0.1);
\end{tikzpicture} & \hypertarget{R-2F25.1}{R-2F25.1} & Funzionale

Opzionale & Il sistema deve permettere ad uno studente di visualizzare un sommario delle proprie statistiche & \hyperlink{UC15.3}{UC15.3}

\hyperlink{UC15}{UC15}\tabularnewline
\hline
\begin{tikzpicture}
\draw [->, thick] (0.4,0.2) -- (0.4,0.1) -- (1,0.1);
\end{tikzpicture} & \hypertarget{R-2F25.1.1}{R-2F25.1.1} & Funzionale

Opzionale & Il sistema deve mostrare allo studente il totale delle domande eseguite e il numero di risposte corrette & \hyperlink{UC15.3}{UC15.3}

\hyperlink{UC15}{UC15}\tabularnewline
\hline
\begin{tikzpicture}
\draw [->, thick] (0.4,0.2) -- (0.4,0.1) -- (1,0.1);
\end{tikzpicture} & \hypertarget{R-2F25.1.2}{R-2F25.1.2} & Funzionale

Opzionale & Il sistema deve mostrare allo studente il totale dei questionari eseguiti e la media dei risultati ottenuti & \hyperlink{UC15.3}{UC15.3}

\hyperlink{UC15}{UC15}\tabularnewline
\hline
\begin{tikzpicture}
\draw [->, thick] (0.4,0.2) -- (0.4,0.1) -- (1,0.1);
\end{tikzpicture} & \hypertarget{R-2F25.1.3}{R-2F25.1.3} & Funzionale

Opzionale & Il sistema deve mostrare allo studente la media delle difficoltà delle domande a cui ha risposto correttamente & \hyperlink{UC15.3}{UC15.3}

\hyperlink{UC15}{UC15}\tabularnewline
\hline
\begin{tikzpicture}
\draw [->, thick] (0.4,0.2) -- (0.4,0.1) -- (1,0.1);
\end{tikzpicture} & \hypertarget{R-2F25.1.4}{R-2F25.1.4} & Funzionale

Opzionale & Il sistema deve mostrare allo studente la media delle difficoltà delle domande a cui non ha risposto correttamente & \hyperlink{UC15.3}{UC15.3}

\hyperlink{UC15}{UC15}\tabularnewline
\hline
\begin{tikzpicture}
\draw [->, thick] (0.4,0.2) -- (0.4,0.1) -- (1,0.1);
\end{tikzpicture} & \hypertarget{R-2F25.1.5}{R-2F25.1.5} & Funzionale

Opzionale & Per ogni argomento il sistema deve mostrare allo studente il totale delle domande eseguite, il numero di risposte corrette, la media delle difficoltà delle risposte corrette e la media delle difficoltà delle risposte non corrette. & \hyperlink{UC15.3}{UC15.3}

\hyperlink{UC15}{UC15}\tabularnewline
\hline
\begin{tikzpicture}
\draw [->, thick] (0.2,0.2) -- (0.2,0.1) -- (1,0.1);
\end{tikzpicture} & \hypertarget{R-2F25.2}{R-2F25.2} & Funzionale

Opzionale & Il sistema deve permettere ad uno studente di visualizzare le statistiche relative ad una domanda eseguita & \hyperlink{UC15.1}{UC15.1}

\hyperlink{UC15}{UC15}\tabularnewline
\hline
\begin{tikzpicture}
\draw [->, thick] (0.4,0.2) -- (0.4,0.1) -- (1,0.1);
\end{tikzpicture} & \hypertarget{R-2F25.2.1}{R-2F25.2.1} & Funzionale

Opzionale & Il sistema deve mostrare allo studente le risposte già date in passato alla domanda selezionata & \hyperlink{UC15.1}{UC15.1}

\hyperlink{UC15}{UC15}\tabularnewline
\hline
\begin{tikzpicture}
\draw [->, thick] (0.4,0.2) -- (0.4,0.1) -- (1,0.1);
\end{tikzpicture} & \hypertarget{R-2F25.2.2}{R-2F25.2.2} & Funzionale

Opzionale & Il sistema deve mostrare allo studente la difficoltà della domanda calcolata come funzione delle risposte date da tutti gli utenti & \hyperlink{UC15.1}{UC15.1}

\hyperlink{UC15}{UC15}\tabularnewline
\hline
\begin{tikzpicture}
\draw [->, thick] (0.2,0.2) -- (0.2,0.1) -- (1,0.1);
\end{tikzpicture} & \hypertarget{R-2F25.3}{R-2F25.3} & Funzionale

Opzionale & Il sistema deve permettere ad uno studente di visualizzare le statistiche e le risposte relative ad un questionario eseguito & \hyperlink{UC15.2}{UC15.2}

\hyperlink{UC15}{UC15}\tabularnewline
\hline
\begin{tikzpicture}
\draw [->, thick] (0.4,0.2) -- (0.4,0.1) -- (1,0.1);
\end{tikzpicture} & \hypertarget{R-2F25.3.1}{R-2F25.3.1} & Funzionale

Opzionale & Il sistema deve consentire ad uno studente di ricercare e selezionare uno dei questionari eseguiti & \hyperlink{UC15.2}{UC15.2}

\hyperlink{UC15}{UC15}

\hyperlink{UC16}{UC16}\tabularnewline
\hline
\begin{tikzpicture}
\draw [->, thick] (0.4,0.2) -- (0.4,0.1) -- (1,0.1);
\end{tikzpicture} & \hypertarget{R-2F25.3.2}{R-2F25.3.2} & Funzionale

Opzionale & Il sistema deve mostrare allo studente le statistiche di ogni singola domanda del questionario selezionato & \hyperlink{UC15.2}{UC15.2}

\hyperlink{UC15}{UC15}\tabularnewline
\hline
\begin{tikzpicture}
\draw [->, thick] (0.4,0.2) -- (0.4,0.1) -- (1,0.1);
\end{tikzpicture} & \hypertarget{R-2F25.3.3}{R-2F25.3.3} & Funzionale

Opzionale & Il sistema deve mostrare allo studente le proprie risposte al questionario selezionato e il punteggio totale & \hyperlink{UC15.2}{UC15.2}

\hyperlink{UC15}{UC15}\tabularnewline
\hline
\begin{tikzpicture}
\draw [->, thick] (0.4,0.2) -- (0.4,0.1) -- (1,0.1);
\end{tikzpicture} & \hypertarget{R-2F25.3.4}{R-2F25.3.4} & Funzionale

Opzionale & Il sistema deve mostrare allo studente la difficoltà del questionario calcolata come funzione delle risposte date da tutti gli utenti & \hyperlink{UC15.2}{UC15.2}

\hyperlink{UC15}{UC15}\tabularnewline
\hline
 & \hypertarget{R-3Q26}{R-3Q26} & Qualità

Obbligatorio & La progettazione rispetta le norme e le
metriche riportate nei documenti \NdP\ \textit{v2.0.0} e \PdQ\ \textit{v2.0.0} & Interno\tabularnewline
\hline
 & \hypertarget{R-3Q27}{R-3Q27} & Qualità

Obbligatorio & Devono essere rispettate tutte le norme e le metriche sulla stesura di codice riportate nei documenti \NdP\ \textit{v2.0.0} e \PdQ\ \textit{v2.0.0} & Interno\tabularnewline
\hline
 & \hypertarget{R-3Q28}{R-3Q28} & Qualità

Obbligatorio & Deve essere realizzato un manualeutente che descriva l’utilizzo del software & Interno\tabularnewline
\hline
 & \hypertarget{R-2Q29}{R-2Q29} & Qualità

Opzionale & Deve essere prodotta documentazione in jsDoc del codice sorgente del software scritto in Javascript & Interno\tabularnewline
\hline
 & \hypertarget{R-3F30}{R-3F30} & Funzionale

Obbligatorio & Il sistema deve permettere ad un docente di visualizzare una domanda & \hyperlink{UC5.12}{UC5.12}\tabularnewline
\hline
 & \hypertarget{R-3F31}{R-3F31} & Funzionale

Obbligatorio & Il sistema deve permettere ad un docente di visualizzare un questionario & \hyperlink{UC6.5}{UC6.5}\tabularnewline
\hline
\caption{Tabella requisiti / fonti} \tabularnewline
\end{longtable}


\newpage
\subsection{Tabella casi d'uso/requisiti}

% Output di tracy sotto a:
% tracy-tpiga.rhcloud.com/tracy/latex/SourceRequirement?type=latex&raw=1
\begin{longtable}{|r l p{5cm}|p{3cm}|}
\hline
\multicolumn{3}{|c|}{Fonte} & Requisito\tabularnewline
\hline
 &  & Capitolato & \hyperlink{R-3V1}{R-3V1}

\hyperlink{R-3V2}{R-3V2}

\hyperlink{R-3V3}{R-3V3}

\hyperlink{R-3V3.1}{R-3V3.1}

\hyperlink{R-3V3.2}{R-3V3.2}

\hyperlink{R-3V4}{R-3V4}

\hyperlink{R-3V3.3}{R-3V3.3}

\hyperlink{R-3V5}{R-3V5}

\hyperlink{R-3V6}{R-3V6}

\hyperlink{R-3V6.1}{R-3V6.1}

\hyperlink{R-3V6.2}{R-3V6.2}

\hyperlink{R-3F7.1}{R-3F7.1}

\hyperlink{R-3F7.5.1}{R-3F7.5.1}

\hyperlink{R-3F7.2}{R-3F7.2}

\hyperlink{R-3F7.3}{R-3F7.3}

\hyperlink{R-2F7.5.2}{R-2F7.5.2}

\hyperlink{R-3F7.4}{R-3F7.4}

\hyperlink{R-3F7.5}{R-3F7.5}

\hyperlink{R-3F7}{R-3F7}

\hyperlink{R-3F7.7}{R-3F7.7}

\hyperlink{R-3F8}{R-3F8}

\hyperlink{R-3F9}{R-3F9}

\hyperlink{R-3V11}{R-3V11}

\hyperlink{R-3V3.4}{R-3V3.4}

\hyperlink{R-2F7.8}{R-2F7.8}

\hyperlink{R-2F7.9}{R-2F7.9}

\hyperlink{R-2F7.10}{R-2F7.10}

\hyperlink{R-2F7.10.1}{R-2F7.10.1}

\hyperlink{R-2F7.10.2}{R-2F7.10.2}

\hyperlink{R-2F7.10.3}{R-2F7.10.3}

\hyperlink{R-3F7.5.5}{R-3F7.5.5}

\hyperlink{R-3F7.5.1.2}{R-3F7.5.1.2}\tabularnewline
\hline
 &  & Proponente & \hyperlink{R-3F7.5.3}{R-3F7.5.3}

\hyperlink{R-2F7.5.4}{R-2F7.5.4}

\hyperlink{R-2F7.6}{R-2F7.6}\tabularnewline
\hline
 &  & Committente & \tabularnewline
\hline
 &  & Interno & \hyperlink{R-3V11}{R-3V11}

\hyperlink{R-3F10.1}{R-3F10.1}

\hyperlink{R-3F19.1}{R-3F19.1}

\hyperlink{R-3Q26}{R-3Q26}

\hyperlink{R-3Q27}{R-3Q27}

\hyperlink{R-3Q28}{R-3Q28}

\hyperlink{R-2Q29}{R-2Q29}\tabularnewline
\hline
 & \hyperlink{UC1}{UC1} & \hyperlink{UC1}{Autenticazione} & \tabularnewline
\hline
\begin{tikzpicture}
\draw [->, thick] (0.2,0.2) -- (0.2,0.1) -- (1,0.1);
\end{tikzpicture} & \hyperlink{UC1.1}{UC1.1} & \hyperlink{UC1.1}{Inserimento username} & \tabularnewline
\hline
\begin{tikzpicture}
\draw [->, thick] (0.2,0.2) -- (0.2,0.1) -- (1,0.1);
\end{tikzpicture} & \hyperlink{UC1.2}{UC1.2} & \hyperlink{UC1.2}{Inserimento password} & \tabularnewline
\hline
\begin{tikzpicture}
\draw [->, thick] (0.2,0.2) -- (0.2,0.1) -- (1,0.1);
\end{tikzpicture} & \hyperlink{UC1.3}{UC1.3} & \hyperlink{UC1.3}{Errore username non presente} & \tabularnewline
\hline
\begin{tikzpicture}
\draw [->, thick] (0.2,0.2) -- (0.2,0.1) -- (1,0.1);
\end{tikzpicture} & \hyperlink{UC1.4}{UC1.4} & \hyperlink{UC1.4}{Errore password errata} & \tabularnewline
\hline
 & \hyperlink{UC2}{UC2} & \hyperlink{UC2}{Registrazione} & \hyperlink{R-3F10}{R-3F10}\tabularnewline
\hline
\begin{tikzpicture}
\draw [->, thick] (0.2,0.2) -- (0.2,0.1) -- (1,0.1);
\end{tikzpicture} & \hyperlink{UC2.1}{UC2.1} & \hyperlink{UC2.1}{Inserimento nome completo} & \hyperlink{R-3F10.2}{R-3F10.2}

\hyperlink{R-3F10.2.2}{R-3F10.2.2}\tabularnewline
\hline
\begin{tikzpicture}
\draw [->, thick] (0.2,0.2) -- (0.2,0.1) -- (1,0.1);
\end{tikzpicture} & \hyperlink{UC2.2}{UC2.2} & \hyperlink{UC2.2}{Inserimento username} & \hyperlink{R-3F10.2}{R-3F10.2}\tabularnewline
\hline
\begin{tikzpicture}
\draw [->, thick] (0.2,0.2) -- (0.2,0.1) -- (1,0.1);
\end{tikzpicture} & \hyperlink{UC2.3}{UC2.3} & \hyperlink{UC2.3}{Inserimento password} & \hyperlink{R-3F10.2}{R-3F10.2}\tabularnewline
\hline
\begin{tikzpicture}
\draw [->, thick] (0.2,0.2) -- (0.2,0.1) -- (1,0.1);
\end{tikzpicture} & \hyperlink{UC2.4}{UC2.4} & \hyperlink{UC2.4}{Errore nome completo troppo corto} & \hyperlink{R-3F10.2.1}{R-3F10.2.1}\tabularnewline
\hline
\begin{tikzpicture}
\draw [->, thick] (0.2,0.2) -- (0.2,0.1) -- (1,0.1);
\end{tikzpicture} & \hyperlink{UC2.5}{UC2.5} & \hyperlink{UC2.5}{Errore username troppo corto} & \hyperlink{R-3F10.2.1}{R-3F10.2.1}\tabularnewline
\hline
\begin{tikzpicture}
\draw [->, thick] (0.2,0.2) -- (0.2,0.1) -- (1,0.1);
\end{tikzpicture} & \hyperlink{UC2.6}{UC2.6} & \hyperlink{UC2.6}{Errore password troppo corta} & \hyperlink{R-3F10.4}{R-3F10.4}

\hyperlink{R-3F10.2.1}{R-3F10.2.1}\tabularnewline
\hline
\begin{tikzpicture}
\draw [->, thick] (0.2,0.2) -- (0.2,0.1) -- (1,0.1);
\end{tikzpicture} & \hyperlink{UC2.7}{UC2.7} & \hyperlink{UC2.7}{Errore username già utilizzato} & \hyperlink{R-3F10.3}{R-3F10.3}

\hyperlink{R-3F10.5}{R-3F10.5}

\hyperlink{R-3F10.2.1}{R-3F10.2.1}\tabularnewline
\hline
 & \hyperlink{UC3}{UC3} & \hyperlink{UC3}{Gestione profilo} & \hyperlink{R-3F14}{R-3F14}

\hyperlink{R-3F14.1}{R-3F14.1}

\hyperlink{R-3F14.2}{R-3F14.2}

\hyperlink{R-3F14.1.1}{R-3F14.1.1}

\hyperlink{R-3F14.2.1}{R-3F14.2.1}

\hyperlink{R-3F14.3}{R-3F14.3}

\hyperlink{R-3F14.3.1}{R-3F14.3.1}

\hyperlink{R-3F14.3.2}{R-3F14.3.2}

\hyperlink{R-3F14.2.1.1}{R-3F14.2.1.1}

\hyperlink{R-3F14.2.2}{R-3F14.2.2}

\hyperlink{R-3F14.2.2.1}{R-3F14.2.2.1}

\hyperlink{R-3F14.3.3}{R-3F14.3.3}

\hyperlink{R-3F14.1.2}{R-3F14.1.2}

\hyperlink{R-3F14.3.4}{R-3F14.3.4}\tabularnewline
\hline
\begin{tikzpicture}
\draw [->, thick] (0.2,0.2) -- (0.2,0.1) -- (1,0.1);
\end{tikzpicture} & \hyperlink{UC3.1}{UC3.1} & \hyperlink{UC3.1}{Modifica nome completo} & \hyperlink{R-3F14.3}{R-3F14.3}

\hyperlink{R-3F14.3.1}{R-3F14.3.1}

\hyperlink{R-3F14.3.3}{R-3F14.3.3}\tabularnewline
\hline
\begin{tikzpicture}
\draw [->, thick] (0.2,0.2) -- (0.2,0.1) -- (1,0.1);
\end{tikzpicture} & \hyperlink{UC3.2}{UC3.2} & \hyperlink{UC3.2}{Modifica username} & \hyperlink{R-3F14.1}{R-3F14.1}

\hyperlink{R-3F14.1.1}{R-3F14.1.1}

\hyperlink{R-3F14.3}{R-3F14.3}

\hyperlink{R-3F14.3.2}{R-3F14.3.2}

\hyperlink{R-3F14.3.4}{R-3F14.3.4}\tabularnewline
\hline
\begin{tikzpicture}
\draw [->, thick] (0.2,0.2) -- (0.2,0.1) -- (1,0.1);
\end{tikzpicture} & \hyperlink{UC3.3}{UC3.3} & \hyperlink{UC3.3}{Modifica password} & \hyperlink{R-3F14.2}{R-3F14.2}

\hyperlink{R-3F14.2.1}{R-3F14.2.1}

\hyperlink{R-3F14.2.2}{R-3F14.2.2}

\hyperlink{R-3F14.2.2.1}{R-3F14.2.2.1}\tabularnewline
\hline
\begin{tikzpicture}
\draw [->, thick] (0.4,0.2) -- (0.4,0.1) -- (1,0.1);
\end{tikzpicture} & \hyperlink{UC3.3.1}{UC3.3.1} & \hyperlink{UC3.3.1}{Inserimento vecchia password} & \hyperlink{R-3F14.2.1}{R-3F14.2.1}\tabularnewline
\hline
\begin{tikzpicture}
\draw [->, thick] (0.4,0.2) -- (0.4,0.1) -- (1,0.1);
\end{tikzpicture} & \hyperlink{UC3.3.2}{UC3.3.2} & \hyperlink{UC3.3.2}{Errore password non corrispondenti} & \hyperlink{R-3F14.2.1.1}{R-3F14.2.1.1}\tabularnewline
\hline
\begin{tikzpicture}
\draw [->, thick] (0.2,0.2) -- (0.2,0.1) -- (1,0.1);
\end{tikzpicture} & \hyperlink{UC3.4}{UC3.4} & \hyperlink{UC3.4}{Errore username troppo corto} & \hyperlink{R-3F14.1.1}{R-3F14.1.1}

\hyperlink{R-3F14.3.4}{R-3F14.3.4}\tabularnewline
\hline
\begin{tikzpicture}
\draw [->, thick] (0.2,0.2) -- (0.2,0.1) -- (1,0.1);
\end{tikzpicture} & \hyperlink{UC3.5}{UC3.5} & \hyperlink{UC3.5}{Errore nome completo troppo corto} & \hyperlink{R-3F14.3.3}{R-3F14.3.3}\tabularnewline
\hline
\begin{tikzpicture}
\draw [->, thick] (0.2,0.2) -- (0.2,0.1) -- (1,0.1);
\end{tikzpicture} & \hyperlink{UC3.6}{UC3.6} & \hyperlink{UC3.6}{Errore username già utilizzato} & \hyperlink{R-3F14.1.1}{R-3F14.1.1}

\hyperlink{R-3F14.1.2}{R-3F14.1.2}\tabularnewline
\hline
\begin{tikzpicture}
\draw [->, thick] (0.2,0.2) -- (0.2,0.1) -- (1,0.1);
\end{tikzpicture} & \hyperlink{UC3.7}{UC3.7} & \hyperlink{UC3.7}{Errore password troppo corta} & \hyperlink{R-3F14.2.1.1}{R-3F14.2.1.1}

\hyperlink{R-3F14.2.2.1}{R-3F14.2.2.1}\tabularnewline
\hline
 & \hyperlink{UC4}{UC4} & \hyperlink{UC4}{Logout} & \hyperlink{R-3F18}{R-3F18}\tabularnewline
\hline
 & \hyperlink{UC5}{UC5} & \hyperlink{UC5}{Gestione domande} & \hyperlink{R-3F7.2}{R-3F7.2}

\hyperlink{R-3F7.11}{R-3F7.11}

\hyperlink{R-3F7.11.1}{R-3F7.11.1}

\hyperlink{R-3F7.11.2}{R-3F7.11.2}

\hyperlink{R-3F7.11.3}{R-3F7.11.3}

\hyperlink{R-3F7.11.1.1}{R-3F7.11.1.1}

\hyperlink{R-3F7.11.1.2}{R-3F7.11.1.2}

\hyperlink{R-3F7.11.1.1.1}{R-3F7.11.1.1.1}

\hyperlink{R-3F7.11.1.2.1}{R-3F7.11.1.2.1}\tabularnewline
\hline
\begin{tikzpicture}
\draw [->, thick] (0.2,0.2) -- (0.2,0.1) -- (1,0.1);
\end{tikzpicture} & \hyperlink{UC5.1}{UC5.1} & \hyperlink{UC5.1}{Inserimento domanda} & \hyperlink{R-2F7.5.2}{R-2F7.5.2}

\hyperlink{R-3F7.5.3}{R-3F7.5.3}

\hyperlink{R-3F7.11.1}{R-3F7.11.1}

\hyperlink{R-3F7.11.1.1}{R-3F7.11.1.1}

\hyperlink{R-3F7.11.1.2.2}{R-3F7.11.1.2.2}

\hyperlink{R-3F7.11.1.2.3}{R-3F7.11.1.2.3}

\hyperlink{R-1F7.11.1.2.4}{R-1F7.11.1.2.4}

\hyperlink{R-1F7.11.1.2.5}{R-1F7.11.1.2.5}

\hyperlink{R-1F7.11.1.2.6}{R-1F7.11.1.2.6}

\hyperlink{R-2F7.11.1.2.7}{R-2F7.11.1.2.7}

\hyperlink{R-3F7.5.5}{R-3F7.5.5}\tabularnewline
\hline
\begin{tikzpicture}
\draw [->, thick] (0.4,0.2) -- (0.4,0.1) -- (1,0.1);
\end{tikzpicture} & \hyperlink{UC5.1.1}{UC5.1.1} & \hyperlink{UC5.1.1}{Selezione argomenti nuova domanda} & \hyperlink{R-3F7.11.1.1}{R-3F7.11.1.1}

\hyperlink{R-3F7.11.1.1.1}{R-3F7.11.1.1.1}\tabularnewline
\hline
\begin{tikzpicture}
\draw [->, thick] (0.4,0.2) -- (0.4,0.1) -- (1,0.1);
\end{tikzpicture} & \hyperlink{UC5.1.2}{UC5.1.2} & \hyperlink{UC5.1.2}{Scrittura domanda in QML della nuova domanda} & \hyperlink{R-3F7.5.1}{R-3F7.5.1}

\hyperlink{R-3F7.5}{R-3F7.5}

\hyperlink{R-3F7.11.1.2}{R-3F7.11.1.2}

\hyperlink{R-3F7.11.1.2.1}{R-3F7.11.1.2.1}

\hyperlink{R-2F7.5.1.1}{R-2F7.5.1.1}

\hyperlink{R-3F7.5.1.2}{R-3F7.5.1.2}

\hyperlink{R-2F7.5.1.3}{R-2F7.5.1.3}

\hyperlink{R-2F7.5.1.4}{R-2F7.5.1.4}

\hyperlink{R-2F7.5.1.5}{R-2F7.5.1.5}

\hyperlink{R-3F7.5.1.6}{R-3F7.5.1.6}

\hyperlink{R-2F7.5.1.7}{R-2F7.5.1.7}

\hyperlink{R-2F7.5.1.8}{R-2F7.5.1.8}\tabularnewline
\hline
\begin{tikzpicture}
\draw [->, thick] (0.4,0.2) -- (0.4,0.1) -- (1,0.1);
\end{tikzpicture} & \hyperlink{UC5.1.3}{UC5.1.3} & \hyperlink{UC5.1.3}{Scrittura nuova domanda da interfaccia grafica} & \hyperlink{R-2F7.5.4}{R-2F7.5.4}

\hyperlink{R-2F7.5.1.7}{R-2F7.5.1.7}

\hyperlink{R-2F7.5.1.8}{R-2F7.5.1.8}\tabularnewline
\hline
\begin{tikzpicture}
\draw [->, thick] (0.2,0.2) -- (0.2,0.1) -- (1,0.1);
\end{tikzpicture} & \hyperlink{UC5.2}{UC5.2} & \hyperlink{UC5.2}{Modifica domanda} & \hyperlink{R-3F7.11.2}{R-3F7.11.2}\tabularnewline
\hline
\begin{tikzpicture}
\draw [->, thick] (0.4,0.2) -- (0.4,0.1) -- (1,0.1);
\end{tikzpicture} & \hyperlink{UC5.2.1}{UC5.2.1} & \hyperlink{UC5.2.1}{Selezione argomenti modifica domanda} & \tabularnewline
\hline
\begin{tikzpicture}
\draw [->, thick] (0.4,0.2) -- (0.4,0.1) -- (1,0.1);
\end{tikzpicture} & \hyperlink{UC5.2.2}{UC5.2.2} & \hyperlink{UC5.2.2}{Scrittura domanda in QML della domanda da modificare} & \tabularnewline
\hline
\begin{tikzpicture}
\draw [->, thick] (0.4,0.2) -- (0.4,0.1) -- (1,0.1);
\end{tikzpicture} & \hyperlink{UC5.2.3}{UC5.2.3} & \hyperlink{UC5.2.3}{Modifica domanda da interfaccia grafica} & \tabularnewline
\hline
\begin{tikzpicture}
\draw [->, thick] (0.2,0.2) -- (0.2,0.1) -- (1,0.1);
\end{tikzpicture} & \hyperlink{UC5.3}{UC5.3} & \hyperlink{UC5.3}{Elimina domanda} & \hyperlink{R-3F7.11.3}{R-3F7.11.3}\tabularnewline
\hline
\begin{tikzpicture}
\draw [->, thick] (0.2,0.2) -- (0.2,0.1) -- (1,0.1);
\end{tikzpicture} & \hyperlink{UC5.4}{UC5.4} & \hyperlink{UC5.4}{Errore QML non valido} & \hyperlink{R-3F7.11.1.2.1}{R-3F7.11.1.2.1}\tabularnewline
\hline
\begin{tikzpicture}
\draw [->, thick] (0.2,0.2) -- (0.2,0.1) -- (1,0.1);
\end{tikzpicture} & \hyperlink{UC5.5}{UC5.5} & \hyperlink{UC5.5}{Errore argomento mancante} & \hyperlink{R-3F7.11.1.1.1}{R-3F7.11.1.1.1}\tabularnewline
\hline
\begin{tikzpicture}
\draw [->, thick] (0.2,0.2) -- (0.2,0.1) -- (1,0.1);
\end{tikzpicture} & \hyperlink{UC5.6}{UC5.6} & \hyperlink{UC5.6}{Inserimento domanda di tipo vero/falso} & \hyperlink{R-3F7.11.1.2.2}{R-3F7.11.1.2.2}

\hyperlink{R-3F7.5.5}{R-3F7.5.5}\tabularnewline
\hline
\begin{tikzpicture}
\draw [->, thick] (0.2,0.2) -- (0.2,0.1) -- (1,0.1);
\end{tikzpicture} & \hyperlink{UC5.7}{UC5.7} & \hyperlink{UC5.7}{Inserimento domanda a scelta multipla} & \hyperlink{R-2F7.5.2}{R-2F7.5.2}

\hyperlink{R-3F7.11.1.2.3}{R-3F7.11.1.2.3}\tabularnewline
\hline
\begin{tikzpicture}
\draw [->, thick] (0.2,0.2) -- (0.2,0.1) -- (1,0.1);
\end{tikzpicture} & \hyperlink{UC5.8}{UC5.8} & \hyperlink{UC5.8}{Inserimento domanda a risposta multipla} & \hyperlink{R-1F7.11.1.2.4}{R-1F7.11.1.2.4}

\hyperlink{R-3F7.5.5}{R-3F7.5.5}\tabularnewline
\hline
\begin{tikzpicture}
\draw [->, thick] (0.2,0.2) -- (0.2,0.1) -- (1,0.1);
\end{tikzpicture} & \hyperlink{UC5.9}{UC5.9} & \hyperlink{UC5.9}{Inserimento domanda di tipo testo con parole omesse} & \hyperlink{R-2F7.5.2}{R-2F7.5.2}

\hyperlink{R-1F7.11.1.2.5}{R-1F7.11.1.2.5}\tabularnewline
\hline
\begin{tikzpicture}
\draw [->, thick] (0.2,0.2) -- (0.2,0.1) -- (1,0.1);
\end{tikzpicture} & \hyperlink{UC5.10}{UC5.10} & \hyperlink{UC5.10}{Inserimento domanda con l'associazione di parole} & \hyperlink{R-2F7.5.2}{R-2F7.5.2}

\hyperlink{R-1F7.11.1.2.6}{R-1F7.11.1.2.6}\tabularnewline
\hline
\begin{tikzpicture}
\draw [->, thick] (0.2,0.2) -- (0.2,0.1) -- (1,0.1);
\end{tikzpicture} & \hyperlink{UC5.11}{UC5.11} & \hyperlink{UC5.11}{Inserimento domanda a risposta aperta} & \hyperlink{R-2F7.5.2}{R-2F7.5.2}

\hyperlink{R-2F7.11.1.2.7}{R-2F7.11.1.2.7}\tabularnewline
\hline
\begin{tikzpicture}
\draw [->, thick] (0.2,0.2) -- (0.2,0.1) -- (1,0.1);
\end{tikzpicture} & \hyperlink{UC5.12}{UC5.12} & \hyperlink{UC5.12}{Visualizza domanda} & \hyperlink{R-3F7.2}{R-3F7.2}

\hyperlink{R-3F30}{R-3F30}\tabularnewline
\hline
\begin{tikzpicture}
\draw [->, thick] (0.2,0.2) -- (0.2,0.1) -- (1,0.1);
\end{tikzpicture} & \hyperlink{UC5.13}{UC5.13} & \hyperlink{UC5.13}{Errore eliminazione domanda utilizzata} & \tabularnewline
\hline
 & \hyperlink{UC6}{UC6} & \hyperlink{UC6}{Gestione questionari} & \hyperlink{R-3F7.2}{R-3F7.2}

\hyperlink{R-3F7.3}{R-3F7.3}

\hyperlink{R-3F7}{R-3F7}

\hyperlink{R-3F7.7}{R-3F7.7}

\hyperlink{R-3F7.11}{R-3F7.11}

\hyperlink{R-3F7.12}{R-3F7.12}

\hyperlink{R-3F7.12.2}{R-3F7.12.2}

\hyperlink{R-3F7.12.3}{R-3F7.12.3}

\hyperlink{R-3F7.13}{R-3F7.13}

\hyperlink{R-3F7.12.4}{R-3F7.12.4}

\hyperlink{R-3F7.7.1}{R-3F7.7.1}\tabularnewline
\hline
\begin{tikzpicture}
\draw [->, thick] (0.2,0.2) -- (0.2,0.1) -- (1,0.1);
\end{tikzpicture} & \hyperlink{UC6.1}{UC6.1} & \hyperlink{UC6.1}{Inserisci questionario} & \hyperlink{R-3F7.7}{R-3F7.7}

\hyperlink{R-3F7.7.1}{R-3F7.7.1}\tabularnewline
\hline
\begin{tikzpicture}
\draw [->, thick] (0.4,0.2) -- (0.4,0.1) -- (1,0.1);
\end{tikzpicture} & \hyperlink{UC6.1.1}{UC6.1.1} & \hyperlink{UC6.1.1}{Aggiungi domanda in un nuovo questionario } & \hyperlink{R-3F7.11}{R-3F7.11}

\hyperlink{R-3F7.11.1.2}{R-3F7.11.1.2}

\hyperlink{R-3F7.11.1.1.1}{R-3F7.11.1.1.1}

\hyperlink{R-3F7.11.1.2.1}{R-3F7.11.1.2.1}\tabularnewline
\hline
\begin{tikzpicture}
\draw [->, thick] (0.4,0.2) -- (0.4,0.1) -- (1,0.1);
\end{tikzpicture} & \hyperlink{UC6.1.2}{UC6.1.2} & \hyperlink{UC6.1.2}{Elimina domanda da un nuovo questionario} & \hyperlink{R-3F7.11}{R-3F7.11}\tabularnewline
\hline
\begin{tikzpicture}
\draw [->, thick] (0.4,0.2) -- (0.4,0.1) -- (1,0.1);
\end{tikzpicture} & \hyperlink{UC6.1.3}{UC6.1.3} & \hyperlink{UC6.1.3}{Seleziona argomenti del nuovo questionario} & \tabularnewline
\hline
\begin{tikzpicture}
\draw [->, thick] (0.2,0.2) -- (0.2,0.1) -- (1,0.1);
\end{tikzpicture} & \hyperlink{UC6.2}{UC6.2} & \hyperlink{UC6.2}{Modifica questionario} & \hyperlink{R-3F7.12}{R-3F7.12}

\hyperlink{R-3F7.12.4}{R-3F7.12.4}\tabularnewline
\hline
\begin{tikzpicture}
\draw [->, thick] (0.4,0.2) -- (0.4,0.1) -- (1,0.1);
\end{tikzpicture} & \hyperlink{UC6.2.1}{UC6.2.1} & \hyperlink{UC6.2.1}{Aggiungi domanda in un questionario} & \hyperlink{R-3F7.12.2}{R-3F7.12.2}\tabularnewline
\hline
\begin{tikzpicture}
\draw [->, thick] (0.4,0.2) -- (0.4,0.1) -- (1,0.1);
\end{tikzpicture} & \hyperlink{UC6.2.2}{UC6.2.2} & \hyperlink{UC6.2.2}{Elimina domanda da un questionario da  modificare} & \hyperlink{R-3F7.12.3}{R-3F7.12.3}\tabularnewline
\hline
\begin{tikzpicture}
\draw [->, thick] (0.4,0.2) -- (0.4,0.1) -- (1,0.1);
\end{tikzpicture} & \hyperlink{UC6.2.3}{UC6.2.3} & \hyperlink{UC6.2.3}{Selezione argomenti modifica questionario} & \hyperlink{R-3F7.12.4}{R-3F7.12.4}\tabularnewline
\hline
\begin{tikzpicture}
\draw [->, thick] (0.2,0.2) -- (0.2,0.1) -- (1,0.1);
\end{tikzpicture} & \hyperlink{UC6.3}{UC6.3} & \hyperlink{UC6.3}{Elimina questionario} & \hyperlink{R-3F7.13}{R-3F7.13}\tabularnewline
\hline
\begin{tikzpicture}
\draw [->, thick] (0.2,0.2) -- (0.2,0.1) -- (1,0.1);
\end{tikzpicture} & \hyperlink{UC6.4}{UC6.4} & \hyperlink{UC6.4}{Errore questionario vuoto} & \hyperlink{R-3F7.7.1}{R-3F7.7.1}\tabularnewline
\hline
\begin{tikzpicture}
\draw [->, thick] (0.2,0.2) -- (0.2,0.1) -- (1,0.1);
\end{tikzpicture} & \hyperlink{UC6.5}{UC6.5} & \hyperlink{UC6.5}{Visualizza questionario} & \hyperlink{R-3F31}{R-3F31}\tabularnewline
\hline
 & \hyperlink{UC7}{UC7} & \hyperlink{UC7}{Gestione classi} & \hyperlink{R-2F13.1.1}{R-2F13.1.1}

\hyperlink{R-2F13}{R-2F13}

\hyperlink{R-2F13.1}{R-2F13.1}

\hyperlink{R-2F13.2}{R-2F13.2}

\hyperlink{R-2F13.3}{R-2F13.3}

\hyperlink{R-2F13.1.2}{R-2F13.1.2}

\hyperlink{R-2F13.1.3}{R-2F13.1.3}

\hyperlink{R-2F13.3.1}{R-2F13.3.1}

\hyperlink{R-2F13.3.2}{R-2F13.3.2}

\hyperlink{R-2F13.3.3}{R-2F13.3.3}

\hyperlink{R-2F13.1.1.1}{R-2F13.1.1.1}\tabularnewline
\hline
\begin{tikzpicture}
\draw [->, thick] (0.2,0.2) -- (0.2,0.1) -- (1,0.1);
\end{tikzpicture} & \hyperlink{UC7.1}{UC7.1} & \hyperlink{UC7.1}{Inserisci classe} & \hyperlink{R-2F7.6}{R-2F7.6}

\hyperlink{R-2F13.1.1}{R-2F13.1.1}

\hyperlink{R-2F13.1}{R-2F13.1}

\hyperlink{R-2F13.1.2}{R-2F13.1.2}

\hyperlink{R-2F13.1.3}{R-2F13.1.3}

\hyperlink{R-2F13.1.1.1}{R-2F13.1.1.1}\tabularnewline
\hline
\begin{tikzpicture}
\draw [->, thick] (0.4,0.2) -- (0.4,0.1) -- (1,0.1);
\end{tikzpicture} & \hyperlink{UC7.1.1}{UC7.1.1} & \hyperlink{UC7.1.1}{Inserisci nome classe} & \hyperlink{R-2F13.1.1}{R-2F13.1.1}

\hyperlink{R-2F13.3.1}{R-2F13.3.1}

\hyperlink{R-2F13.1.1.1}{R-2F13.1.1.1}\tabularnewline
\hline
\begin{tikzpicture}
\draw [->, thick] (0.4,0.2) -- (0.4,0.1) -- (1,0.1);
\end{tikzpicture} & \hyperlink{UC7.1.2}{UC7.1.2} & \hyperlink{UC7.1.2}{Inserisci argomenti classe} & \hyperlink{R-2F13.1.2}{R-2F13.1.2}

\hyperlink{R-2F13.3.2}{R-2F13.3.2}\tabularnewline
\hline
\begin{tikzpicture}
\draw [->, thick] (0.4,0.2) -- (0.4,0.1) -- (1,0.1);
\end{tikzpicture} & \hyperlink{UC7.1.3}{UC7.1.3} & \hyperlink{UC7.1.3}{Inserisci password classe} & \hyperlink{R-2F13.1.3}{R-2F13.1.3}

\hyperlink{R-2F13.3.3}{R-2F13.3.3}\tabularnewline
\hline
\begin{tikzpicture}
\draw [->, thick] (0.2,0.2) -- (0.2,0.1) -- (1,0.1);
\end{tikzpicture} & \hyperlink{UC7.2}{UC7.2} & \hyperlink{UC7.2}{Modifica classe} & \hyperlink{R-2F13.3}{R-2F13.3}

\hyperlink{R-2F13.3.1}{R-2F13.3.1}

\hyperlink{R-2F13.3.2}{R-2F13.3.2}

\hyperlink{R-2F13.3.3}{R-2F13.3.3}\tabularnewline
\hline
\begin{tikzpicture}
\draw [->, thick] (0.2,0.2) -- (0.2,0.1) -- (1,0.1);
\end{tikzpicture} & \hyperlink{UC7.3}{UC7.3} & \hyperlink{UC7.3}{Elimina classe} & \hyperlink{R-2F13.2}{R-2F13.2}\tabularnewline
\hline
\begin{tikzpicture}
\draw [->, thick] (0.2,0.2) -- (0.2,0.1) -- (1,0.1);
\end{tikzpicture} & \hyperlink{UC7.4}{UC7.4} & \hyperlink{UC7.4}{Errore nome classe già presente} & \hyperlink{R-2F13.1.1.1}{R-2F13.1.1.1}\tabularnewline
\hline
 & \hyperlink{UC8}{UC8} & \hyperlink{UC8}{Gestione argomenti} & \hyperlink{R-3F23}{R-3F23}

\hyperlink{R-3F23.1}{R-3F23.1}

\hyperlink{R-3F23.2}{R-3F23.2}

\hyperlink{R-3F23.3}{R-3F23.3}

\hyperlink{R-3F23.4}{R-3F23.4}

\hyperlink{R-3F23.1.1}{R-3F23.1.1}

\hyperlink{R-3F23.2.1}{R-3F23.2.1}

\hyperlink{R-3F23.2.2}{R-3F23.2.2}\tabularnewline
\hline
\begin{tikzpicture}
\draw [->, thick] (0.2,0.2) -- (0.2,0.1) -- (1,0.1);
\end{tikzpicture} & \hyperlink{UC8.1}{UC8.1} & \hyperlink{UC8.1}{Esplorazione argomenti} & \hyperlink{R-3F23.4}{R-3F23.4}

\hyperlink{R-3F23.2.1}{R-3F23.2.1}\tabularnewline
\hline
\begin{tikzpicture}
\draw [->, thick] (0.2,0.2) -- (0.2,0.1) -- (1,0.1);
\end{tikzpicture} & \hyperlink{UC8.2}{UC8.2} & \hyperlink{UC8.2}{Crea argomento} & \hyperlink{R-3F23.1}{R-3F23.1}

\hyperlink{R-3F23.1.1}{R-3F23.1.1}\tabularnewline
\hline
\begin{tikzpicture}
\draw [->, thick] (0.2,0.2) -- (0.2,0.1) -- (1,0.1);
\end{tikzpicture} & \hyperlink{UC8.3}{UC8.3} & \hyperlink{UC8.3}{Argomento già presente nel sistema} & \hyperlink{R-3F23.1.1}{R-3F23.1.1}\tabularnewline
\hline
\begin{tikzpicture}
\draw [->, thick] (0.2,0.2) -- (0.2,0.1) -- (1,0.1);
\end{tikzpicture} & \hyperlink{UC8.4}{UC8.4} & \hyperlink{UC8.4}{Modifica argomento} & \hyperlink{R-3F23.3}{R-3F23.3}\tabularnewline
\hline
\begin{tikzpicture}
\draw [->, thick] (0.2,0.2) -- (0.2,0.1) -- (1,0.1);
\end{tikzpicture} & \hyperlink{UC8.5}{UC8.5} & \hyperlink{UC8.5}{Eliminazione argomento} & \hyperlink{R-3F23.2}{R-3F23.2}

\hyperlink{R-3F23.2.1}{R-3F23.2.1}

\hyperlink{R-3F23.2.2}{R-3F23.2.2}\tabularnewline
\hline
\begin{tikzpicture}
\draw [->, thick] (0.2,0.2) -- (0.2,0.1) -- (1,0.1);
\end{tikzpicture} & \hyperlink{UC8.6}{UC8.6} & \hyperlink{UC8.6}{Errore l'argomento ha domande o questionari} & \hyperlink{R-3F23.2.2}{R-3F23.2.2}\tabularnewline
\hline
 & \hyperlink{UC9}{UC9} & \hyperlink{UC9}{Visualizza statistiche} & \hyperlink{R-2F24.3.1}{R-2F24.3.1}

\hyperlink{R-2F24}{R-2F24}

\hyperlink{R-2F24.1}{R-2F24.1}

\hyperlink{R-2F24.2}{R-2F24.2}

\hyperlink{R-2F24.3}{R-2F24.3}

\hyperlink{R-2F24.1.1}{R-2F24.1.1}

\hyperlink{R-2F24.2.1}{R-2F24.2.1}

\hyperlink{R-2F24.3.2}{R-2F24.3.2}

\hyperlink{R-2F24.3.3}{R-2F24.3.3}

\hyperlink{R-2F24.3.4}{R-2F24.3.4}

\hyperlink{R-2F24.3.4.1}{R-2F24.3.4.1}

\hyperlink{R-2F24.3.5}{R-2F24.3.5}

\hyperlink{R-2F24.3.5.1}{R-2F24.3.5.1}\tabularnewline
\hline
 & \hyperlink{UC10}{UC10} & \hyperlink{UC10}{Visualizza statistiche domanda} & \hyperlink{R-2F24.1}{R-2F24.1}

\hyperlink{R-2F24.1.1}{R-2F24.1.1}\tabularnewline
\hline
 & \hyperlink{UC11}{UC11} & \hyperlink{UC11}{Visualizza statistiche questionario} & \hyperlink{R-2F24.2}{R-2F24.2}

\hyperlink{R-2F24.2.1}{R-2F24.2.1}\tabularnewline
\hline
 & \hyperlink{UC12}{UC12} & \hyperlink{UC12}{Visualizza statistiche classe} & \hyperlink{R-2F24.3.1}{R-2F24.3.1}

\hyperlink{R-2F24.3}{R-2F24.3}

\hyperlink{R-2F24.3.2}{R-2F24.3.2}

\hyperlink{R-2F24.3.3}{R-2F24.3.3}

\hyperlink{R-2F24.3.4}{R-2F24.3.4}

\hyperlink{R-2F24.3.4.1}{R-2F24.3.4.1}

\hyperlink{R-2F24.3.5}{R-2F24.3.5}

\hyperlink{R-2F24.3.5.1}{R-2F24.3.5.1}\tabularnewline
\hline
\begin{tikzpicture}
\draw [->, thick] (0.2,0.2) -- (0.2,0.1) -- (1,0.1);
\end{tikzpicture} & \hyperlink{UC12.1}{UC12.1} & \hyperlink{UC12.1}{Visualizza risultati domande della classe} & \hyperlink{R-2F24.3.5}{R-2F24.3.5}

\hyperlink{R-2F24.3.5.1}{R-2F24.3.5.1}\tabularnewline
\hline
\begin{tikzpicture}
\draw [->, thick] (0.2,0.2) -- (0.2,0.1) -- (1,0.1);
\end{tikzpicture} & \hyperlink{UC12.2}{UC12.2} & \hyperlink{UC12.2}{Visualizza risultati questionari della classe} & \hyperlink{R-2F24.3.1}{R-2F24.3.1}\tabularnewline
\hline
\begin{tikzpicture}
\draw [->, thick] (0.2,0.2) -- (0.2,0.1) -- (1,0.1);
\end{tikzpicture} & \hyperlink{UC12.3}{UC12.3} & \hyperlink{UC12.3}{Visualizza sommario statistiche classe} & \hyperlink{R-2F24.3.3}{R-2F24.3.3}\tabularnewline
\hline
\begin{tikzpicture}
\draw [->, thick] (0.2,0.2) -- (0.2,0.1) -- (1,0.1);
\end{tikzpicture} & \hyperlink{UC12.4}{UC12.4} & \hyperlink{UC12.4}{Visualizza statistiche studente della classe} & \hyperlink{R-2F24.3.4}{R-2F24.3.4}

\hyperlink{R-2F24.3.4.1}{R-2F24.3.4.1}\tabularnewline
\hline
\begin{tikzpicture}
\draw [->, thick] (0.4,0.2) -- (0.4,0.1) -- (1,0.1);
\end{tikzpicture} & \hyperlink{UC12.4.1}{UC12.4.1} & \hyperlink{UC12.4.1}{Visualizza risultati questionari dello studente} & \hyperlink{R-2F24.3.4.1}{R-2F24.3.4.1}\tabularnewline
\hline
 & \hyperlink{UC13}{UC13} & \hyperlink{UC13}{Esegui questionario} & \hyperlink{R-3F7.4}{R-3F7.4}

\hyperlink{R-3F17}{R-3F17}

\hyperlink{R-3F17.1}{R-3F17.1}

\hyperlink{R-3F17.2}{R-3F17.2}

\hyperlink{R-3F17.3}{R-3F17.3}

\hyperlink{R-3F17.4}{R-3F17.4}

\hyperlink{R-3F17.2.1}{R-3F17.2.1}\tabularnewline
\hline
\begin{tikzpicture}
\draw [->, thick] (0.2,0.2) -- (0.2,0.1) -- (1,0.1);
\end{tikzpicture} & \hyperlink{UC13.1}{UC13.1} & \hyperlink{UC13.1}{Rispondi domanda} & \hyperlink{R-3F7.2}{R-3F7.2}

\hyperlink{R-3F8}{R-3F8}

\hyperlink{R-3F17.1}{R-3F17.1}

\hyperlink{R-3F17.1.1}{R-3F17.1.1}

\hyperlink{R-3F17.1.2}{R-3F17.1.2}

\hyperlink{R-1F17.1.3}{R-1F17.1.3}

\hyperlink{R-1F17.1.4}{R-1F17.1.4}

\hyperlink{R-1F17.1.5}{R-1F17.1.5}

\hyperlink{R-2F17.1.6}{R-2F17.1.6}\tabularnewline
\hline
\begin{tikzpicture}
\draw [->, thick] (0.2,0.2) -- (0.2,0.1) -- (1,0.1);
\end{tikzpicture} & \hyperlink{UC13.2}{UC13.2} & \hyperlink{UC13.2}{Errore domanda non risposta} & \hyperlink{R-3F17.2.1}{R-3F17.2.1}\tabularnewline
\hline
\begin{tikzpicture}
\draw [->, thick] (0.2,0.2) -- (0.2,0.1) -- (1,0.1);
\end{tikzpicture} & \hyperlink{UC13.3}{UC13.3} & \hyperlink{UC13.3}{Conferma questionario} & \hyperlink{R-3F17.2}{R-3F17.2}

\hyperlink{R-3F17.2.1}{R-3F17.2.1}\tabularnewline
\hline
\begin{tikzpicture}
\draw [->, thick] (0.2,0.2) -- (0.2,0.1) -- (1,0.1);
\end{tikzpicture} & \hyperlink{UC13.4}{UC13.4} & \hyperlink{UC13.4}{Rispondi domanda vero/falso} & \hyperlink{R-3F17.1.1}{R-3F17.1.1}\tabularnewline
\hline
\begin{tikzpicture}
\draw [->, thick] (0.2,0.2) -- (0.2,0.1) -- (1,0.1);
\end{tikzpicture} & \hyperlink{UC13.5}{UC13.5} & \hyperlink{UC13.5}{Rispondi domanda a scelta multipla} & \hyperlink{R-3F17.1.2}{R-3F17.1.2}\tabularnewline
\hline
\begin{tikzpicture}
\draw [->, thick] (0.2,0.2) -- (0.2,0.1) -- (1,0.1);
\end{tikzpicture} & \hyperlink{UC13.6}{UC13.6} & \hyperlink{UC13.6}{Rispondi domanda a risposta multipla} & \hyperlink{R-1F17.1.3}{R-1F17.1.3}\tabularnewline
\hline
\begin{tikzpicture}
\draw [->, thick] (0.2,0.2) -- (0.2,0.1) -- (1,0.1);
\end{tikzpicture} & \hyperlink{UC13.7}{UC13.7} & \hyperlink{UC13.7}{Rispondi domanda di tipo testo con parole omesse} & \hyperlink{R-1F17.1.4}{R-1F17.1.4}\tabularnewline
\hline
\begin{tikzpicture}
\draw [->, thick] (0.2,0.2) -- (0.2,0.1) -- (1,0.1);
\end{tikzpicture} & \hyperlink{UC13.8}{UC13.8} & \hyperlink{UC13.8}{Rispondi domanda con associazione di parole} & \hyperlink{R-1F17.1.5}{R-1F17.1.5}\tabularnewline
\hline
\begin{tikzpicture}
\draw [->, thick] (0.2,0.2) -- (0.2,0.1) -- (1,0.1);
\end{tikzpicture} & \hyperlink{UC13.9}{UC13.9} & \hyperlink{UC13.9}{Rispondi domanda a risposta aperta} & \hyperlink{R-2F17.1.6}{R-2F17.1.6}\tabularnewline
\hline
\begin{tikzpicture}
\draw [->, thick] (0.2,0.2) -- (0.2,0.1) -- (1,0.1);
\end{tikzpicture} & \hyperlink{UC13.10}{UC13.10} & \hyperlink{UC13.10}{Feedback questionario} & \hyperlink{R-2F7.9}{R-2F7.9}\tabularnewline
\hline
\begin{tikzpicture}
\draw [->, thick] (0.2,0.2) -- (0.2,0.1) -- (1,0.1);
\end{tikzpicture} & \hyperlink{UC13.11}{UC13.11} & \hyperlink{UC13.11}{Visualizza valutazione questionario} & \hyperlink{R-3F7.4}{R-3F7.4}\tabularnewline
\hline
 & \hyperlink{UC14}{UC14} & \hyperlink{UC14}{Iscrizione ad una classe} & \hyperlink{R-2F7.6}{R-2F7.6}

\hyperlink{R-2F16}{R-2F16}

\hyperlink{R-2F16.1}{R-2F16.1}

\hyperlink{R-2F16.1.1}{R-2F16.1.1}\tabularnewline
\hline
\begin{tikzpicture}
\draw [->, thick] (0.2,0.2) -- (0.2,0.1) -- (1,0.1);
\end{tikzpicture} & \hyperlink{UC14.1}{UC14.1} & \hyperlink{UC14.1}{Inserisci password classe} & \hyperlink{R-2F16.1}{R-2F16.1}\tabularnewline
\hline
\begin{tikzpicture}
\draw [->, thick] (0.2,0.2) -- (0.2,0.1) -- (1,0.1);
\end{tikzpicture} & \hyperlink{UC14.2}{UC14.2} & \hyperlink{UC14.2}{Errore password classe} & \hyperlink{R-2F16.1.1}{R-2F16.1.1}\tabularnewline
\hline
 & \hyperlink{UC15}{UC15} & \hyperlink{UC15}{Visualizza storico studente} & \hyperlink{R-3F7.4}{R-3F7.4}

\hyperlink{R-2F7.10}{R-2F7.10}

\hyperlink{R-2F7.10.1}{R-2F7.10.1}

\hyperlink{R-2F7.10.2}{R-2F7.10.2}

\hyperlink{R-2F7.10.3}{R-2F7.10.3}

\hyperlink{R-2F25}{R-2F25}

\hyperlink{R-2F25.1}{R-2F25.1}

\hyperlink{R-2F25.2}{R-2F25.2}

\hyperlink{R-2F25.3}{R-2F25.3}

\hyperlink{R-2F25.3.1}{R-2F25.3.1}

\hyperlink{R-2F25.1.1}{R-2F25.1.1}

\hyperlink{R-2F25.1.2}{R-2F25.1.2}

\hyperlink{R-2F25.1.3}{R-2F25.1.3}

\hyperlink{R-2F25.1.4}{R-2F25.1.4}

\hyperlink{R-2F25.1.5}{R-2F25.1.5}

\hyperlink{R-2F25.2.1}{R-2F25.2.1}

\hyperlink{R-2F25.2.2}{R-2F25.2.2}

\hyperlink{R-2F25.3.2}{R-2F25.3.2}

\hyperlink{R-2F25.3.3}{R-2F25.3.3}

\hyperlink{R-2F25.3.4}{R-2F25.3.4}\tabularnewline
\hline
\begin{tikzpicture}
\draw [->, thick] (0.2,0.2) -- (0.2,0.1) -- (1,0.1);
\end{tikzpicture} & \hyperlink{UC15.1}{UC15.1} & \hyperlink{UC15.1}{Visualizza statistiche domande studente} & \hyperlink{R-2F25.2}{R-2F25.2}

\hyperlink{R-2F25.2.1}{R-2F25.2.1}

\hyperlink{R-2F25.2.2}{R-2F25.2.2}\tabularnewline
\hline
\begin{tikzpicture}
\draw [->, thick] (0.2,0.2) -- (0.2,0.1) -- (1,0.1);
\end{tikzpicture} & \hyperlink{UC15.2}{UC15.2} & \hyperlink{UC15.2}{Visualizza statistiche questionari studente} & \hyperlink{R-2F7.10.2}{R-2F7.10.2}

\hyperlink{R-2F25.3}{R-2F25.3}

\hyperlink{R-2F25.3.1}{R-2F25.3.1}

\hyperlink{R-2F25.3.2}{R-2F25.3.2}

\hyperlink{R-2F25.3.3}{R-2F25.3.3}

\hyperlink{R-2F25.3.4}{R-2F25.3.4}\tabularnewline
\hline
\begin{tikzpicture}
\draw [->, thick] (0.2,0.2) -- (0.2,0.1) -- (1,0.1);
\end{tikzpicture} & \hyperlink{UC15.3}{UC15.3} & \hyperlink{UC15.3}{Visualizza sommario statistiche studente} & \hyperlink{R-2F7.10.3}{R-2F7.10.3}

\hyperlink{R-2F25.1}{R-2F25.1}

\hyperlink{R-2F25.1.1}{R-2F25.1.1}

\hyperlink{R-2F25.1.2}{R-2F25.1.2}

\hyperlink{R-2F25.1.3}{R-2F25.1.3}

\hyperlink{R-2F25.1.4}{R-2F25.1.4}

\hyperlink{R-2F25.1.5}{R-2F25.1.5}\tabularnewline
\hline
 & \hyperlink{UC16}{UC16} & \hyperlink{UC16}{Ricerca questionario} & \hyperlink{R-3F7.12.1}{R-3F7.12.1}

\hyperlink{R-3F15}{R-3F15}

\hyperlink{R-3F15.1}{R-3F15.1}

\hyperlink{R-2F15.2}{R-2F15.2}

\hyperlink{R-3F15.3}{R-3F15.3}

\hyperlink{R-3F15.4}{R-3F15.4}

\hyperlink{R-2F15.5}{R-2F15.5}

\hyperlink{R-2F24.2.1}{R-2F24.2.1}

\hyperlink{R-2F24.3.4.1}{R-2F24.3.4.1}

\hyperlink{R-2F25.3.1}{R-2F25.3.1}\tabularnewline
\hline
 & \hyperlink{UC17}{UC17} & \hyperlink{UC17}{Ricerca questionario per titolo} & \hyperlink{R-3F15.1}{R-3F15.1}\tabularnewline
\hline
 & \hyperlink{UC18}{UC18} & \hyperlink{UC18}{Ricerca questionario per classe} & \hyperlink{R-2F15.2}{R-2F15.2}\tabularnewline
\hline
 & \hyperlink{UC19}{UC19} & \hyperlink{UC19}{Ricerca questionario per argomento} & \hyperlink{R-3F15.4}{R-3F15.4}\tabularnewline
\hline
 & \hyperlink{UC20}{UC20} & \hyperlink{UC20}{Ricerca questionario per docente} & \hyperlink{R-3F15.3}{R-3F15.3}\tabularnewline
\hline
 & \hyperlink{UC21}{UC21} & \hyperlink{UC21}{Ricerca questionario per difficoltà} & \hyperlink{R-2F15.5}{R-2F15.5}\tabularnewline
\hline
 & \hyperlink{UC22}{UC22} & \hyperlink{UC22}{Ricerca domanda} & \hyperlink{R-3F20}{R-3F20}

\hyperlink{R-3F20.1}{R-3F20.1}

\hyperlink{R-3F20.2}{R-3F20.2}

\hyperlink{R-2F20.3}{R-2F20.3}

\hyperlink{R-3F20.4}{R-3F20.4}

\hyperlink{R-2F24.1.1}{R-2F24.1.1}

\hyperlink{R-2F24.3.5.1}{R-2F24.3.5.1}\tabularnewline
\hline
 & \hyperlink{UC23}{UC23} & \hyperlink{UC23}{Ricerca domanda per keywords} & \hyperlink{R-3F20.4}{R-3F20.4}\tabularnewline
\hline
 & \hyperlink{UC24}{UC24} & \hyperlink{UC24}{Ricerca domanda per argomento} & \hyperlink{R-3F20.1}{R-3F20.1}\tabularnewline
\hline
 & \hyperlink{UC25}{UC25} & \hyperlink{UC25}{Ricerca domanda per difficoltà} & \hyperlink{R-2F20.3}{R-2F20.3}\tabularnewline
\hline
 & \hyperlink{UC26}{UC26} & \hyperlink{UC26}{Ricerca domanda per docente} & \hyperlink{R-3F20.2}{R-3F20.2}\tabularnewline
\hline
 & \hyperlink{UC27}{UC27} & \hyperlink{UC27}{Ricerca classe} & \hyperlink{R-2F16.2}{R-2F16.2}

\hyperlink{R-2F21}{R-2F21}

\hyperlink{R-2F21.1}{R-2F21.1}

\hyperlink{R-2F21.2}{R-2F21.2}\tabularnewline
\hline
 & \hyperlink{UC28}{UC28} & \hyperlink{UC28}{Ricerca classe per docente} & \hyperlink{R-2F16.2}{R-2F16.2}

\hyperlink{R-2F21.2}{R-2F21.2}\tabularnewline
\hline
 & \hyperlink{UC29}{UC29} & \hyperlink{UC29}{Ricerca classe per argomento} & \hyperlink{R-2F16.2}{R-2F16.2}

\hyperlink{R-2F21.1}{R-2F21.1}\tabularnewline
\hline
 & \hyperlink{UC30}{UC30} & \hyperlink{UC30}{Azioni Amministratore} & \hyperlink{R-3F12}{R-3F12}

\hyperlink{R-3F12.1}{R-3F12.1}

\hyperlink{R-3F12.2}{R-3F12.2}

\hyperlink{R-3F19}{R-3F19}

\hyperlink{R-3F19.2}{R-3F19.2}

\hyperlink{R-3F19.3}{R-3F19.3}\tabularnewline
\hline
\begin{tikzpicture}
\draw [->, thick] (0.2,0.2) -- (0.2,0.1) -- (1,0.1);
\end{tikzpicture} & \hyperlink{UC30.1}{UC30.1} & \hyperlink{UC30.1}{Cambia ruolo} & \hyperlink{R-3F12.1}{R-3F12.1}

\hyperlink{R-3F19.2}{R-3F19.2}\tabularnewline
\hline
\begin{tikzpicture}
\draw [->, thick] (0.2,0.2) -- (0.2,0.1) -- (1,0.1);
\end{tikzpicture} & \hyperlink{UC30.2}{UC30.2} & \hyperlink{UC30.2}{Rimozione utente} & \hyperlink{R-3F12.2}{R-3F12.2}

\hyperlink{R-3F19.3}{R-3F19.3}\tabularnewline
\hline
 & \hyperlink{UC31}{UC31} & \hyperlink{UC31}{Ricerca utente} & \tabularnewline
\hline
 & \hyperlink{UC32}{UC32} & \hyperlink{UC32}{Ricerca utente per nome completo} & \tabularnewline
\hline
 & \hyperlink{UC33}{UC33} & \hyperlink{UC33}{Ricerca utente per username} & \tabularnewline
\hline
 & \hyperlink{UC34}{UC34} & \hyperlink{UC34}{Ricerca utente per ruolo} & \tabularnewline
\hline
 & \hyperlink{UC35}{UC35} & \hyperlink{UC35}{Recupero password} & \hyperlink{R-2F22}{R-2F22}

\hyperlink{R-2F22.1}{R-2F22.1}

\hyperlink{R-2F22.2}{R-2F22.2}\tabularnewline
\hline
\begin{tikzpicture}
\draw [->, thick] (0.2,0.2) -- (0.2,0.1) -- (1,0.1);
\end{tikzpicture} & \hyperlink{UC35.1}{UC35.1} & \hyperlink{UC35.1}{Errore mail per recupero password non presente} & \hyperlink{R-2F22.2}{R-2F22.2}\tabularnewline
\hline
\caption{Tabella fonti/requisiti} \tabularnewline
\end{longtable}


\end{document}
